\chapter{Centers, actions, and representation tensors}
\label{v4-arith:centers:chapter}

A center records operations that commute with a specified action.
Its definition therefore depends on the objects allowed to act and
on the meaning of commutation. An algebra homomorphism need not
transport even an ordinary center. An equivalence of module
categories can transport a derived center, provided it retains the
coefficient action. These two facts determine the first comparison
we need.

For collision coefficients there is another choice. Convolution
combines inputs on disjoint configurations, whereas diagonal tensor
combines coefficients at the same configuration. Their units differ.
The center of a convolution algebra and the universal diagonal actor
thus solve different problems. Factorization and a specified local
vacuum impose further conditions on an actor.

We first construct scalar Hochschild centers and their Morita
comparisons. Normal differential modules then give explicit centers
whose products, brackets, and coefficient dependence can be computed.
We next compare the word and collision constructions and examine
diagonal actions. The finite fusion example retains both its braided
center and the distinguished algebra that recovers its boundary
category. These data specify what the arithmetic comparison must
preserve.

\section{Maps and centers}
\label{v4-arith:centers:sec:maps}

Let $k$ be a field and let $A$ be a unital associative $k$-algebra.
Its \emph{ordinary center} is
\[
 Z_{\mathrm{alg}}(A)=\{a\in A:ab=ba\text{ for every }b\in A\}.
\]
It is a commutative algebra. A homomorphism $f:A\to B$ sends it
into the centralizer of $f(A)$ in $B$. That centralizer can be
larger than the center of $B$.

\begin{example}[A central element with noncentral image]
\label{v4-arith:centers:ex:diagonal}
Take
\[
 A=k\times k,\qquad B=M_2(k),\qquad
 f(a,b)=
 \begin{pmatrix}a&0\\0&b\end{pmatrix}.
\]
The element $(1,0)$ belongs to $Z_{\mathrm{alg}}(A)$.
Its image is $E_{11}$, and
\[
 E_{11}E_{12}-E_{12}E_{11}=E_{12}\ne0.
\]
Thus $f$ does not restrict to a map
$Z_{\mathrm{alg}}(A)\to Z_{\mathrm{alg}}(B)$.
The centralizer of the diagonal subalgebra is the diagonal
subalgebra itself, whereas $Z_{\mathrm{alg}}(B)=kI_2$.
\end{example}

The derived construction retains this distinction. Put
$A^e=A\otimes_k A^{\mathrm{op}}$, so that left $A^e$-modules
are $A$-bimodules. The \emph{scalar Hochschild center} is
\begin{equation}\label{v4-arith:centers:eq:scalar}
 Z_k(A)=R\operatorname{Hom}_{A^e}(A,A).
\end{equation}
For a differential graded algebra the tensor and Hom in this
formula are derived. We use cohomological grading, with
$M[r]^i=M^{i+r}$.

For an ordinary algebra, a Hochschild $p$-cochain is a
$k$-linear map $A^{\otimes p}\to A$. Its differential is
\begin{equation}\label{v4-arith:centers:eq:hoch-differential}
 \begin{split}
 (bf)(a_1,\ldots,a_{p+1})={}&
 a_1f(a_2,\ldots,a_{p+1})\\
 &+\sum_{i=1}^p(-1)^i
 f(a_1,\ldots,a_ia_{i+1},\ldots,a_{p+1})\\
 &+(-1)^{p+1}f(a_1,\ldots,a_p)a_{p+1}.
 \end{split}
\end{equation}
In degree zero this is $b(a)(x)=xa-ax$. Hence
$H^0Z_k(A)=Z_{\mathrm{alg}}(A)$.

\begin{proposition}[The first two Hochschild groups]
\label{v4-arith:centers:prop:low}
For an ordinary unital algebra,
\[
 H^0Z_k(A)=Z_{\mathrm{alg}}(A),\qquad
 H^1Z_k(A)=\operatorname{Der}_k(A)/\operatorname{Inn}(A).
\]
Here a derivation satisfies $D(ab)=aD(b)+D(a)b$, and an
inner derivation has the form $b\mapsto ab-ba$.
\end{proposition}

\begin{proof}
The degree-zero calculation was just made. For a one-cochain $D$,
\[
 (bD)(a,b)=aD(b)-D(ab)+D(a)b.
\]
Its vanishing is the derivation identity. The image from degree
zero consists exactly of the inner derivations, with an immaterial
change of sign in their parameter.
\end{proof}

A map $f:A\to B$ gives a mixed complex
$C^\bullet(A,B)$, in which $A$ acts on $B$ through $f$.
It receives two natural cochain maps:
\begin{equation}\label{v4-arith:centers:eq:mixed}
 C^\bullet(A,A)\longrightarrow C^\bullet(A,B)
 \longleftarrow C^\bullet(B,B).
\end{equation}
The first applies $f$ to the output. The second restricts all
inputs along $f$. Both preserve cup products. They need not
give a map between the two endpoint centers.

The mixed complex also explains a limitation of insertion
operations. A cochain $A^{\otimes p}\to B$ outputs an element
of $B$. Another such cochain expects an element of $A$ in
each input. Inserting the first into the second is therefore
undefined without additional data. The cup product uses
multiplication in $B$ and has no such typing problem.

\section{The full cochain structure}
\label{v4-arith:centers:sec:cochains}

The derived center remembers more than its cohomology ring.
The bar resolution makes its additional operations explicit.
For an ordinary unital algebra its normalized cochains are
those that vanish whenever an input is $1$. Normalization
does not remove degree-zero cochains.

The cup product of cochains of arities $p$ and $q$ is
\begin{equation}\label{v4-arith:centers:eq:cup}
 (f\smile g)(a_1,\ldots,a_{p+q})
 =f(a_1,\ldots,a_p)g(a_{p+1},\ldots,a_{p+q}).
\end{equation}
For an ungraded algebra, define
\[
 f\circ g=\sum_{i=0}^{p-1}(-1)^{i(q-1)}
 f(1^{\otimes i}\otimes g\otimes1^{\otimes(p-1-i)}).
\]
The \emph{Gerstenhaber bracket} is
\begin{equation}\label{v4-arith:centers:eq:bracket}
 [f,g]_{\mathrm G}
 =f\circ g-(-1)^{(p-1)(q-1)}g\circ f.
\end{equation}
It has cohomological degree $-1$. On one-cochains it is
the commutator of endomorphisms.

Higher \emph{braces} $f\{g_1,\ldots,g_r\}$ insert the
$g_i$ into disjoint slots of $f$, in the indicated order.
The sum runs over all allowed positions. For a differential
graded algebra, suspend the inputs and use the Koszul sign
of each insertion. This convention includes the internal
differential and avoids assigning signs by arity alone.

An $E_n$-algebra is an algebra over chains on the little
$n$-disks operad, or its equivalent derived operad.

\begin{theorem}[The cochain center]
\label{v4-arith:centers:thm:deligne}
The Hochschild center has an $E_2$-algebra structure.
Its product and degree-$-1$ bracket on cohomology are
\eqref{v4-arith:centers:eq:cup} and
\eqref{v4-arith:centers:eq:bracket}. In particular, the cup
product is graded commutative on cohomology, and the bracket
is a graded derivation of that product.
\end{theorem}

This is the cochain form of the Deligne theorem. McClure--Smith identify the unit, cup, and brace operad
with little two-cubes chains. Their result is Theorem 1.1 of
\emph{A solution of Deligne's Hochschild cohomology conjecture}. Their construction is integral. In the
derived setting, apply the construction to a suitable
cofibrant algebra model.

The $E_2$ structure is not the assertion that a selected
cochain cup product is strictly commutative.

\subsection{The obstruction to a higher disk structure}

Let $A=k[x,y,z]$ in characteristic zero. The derivation
$D=\partial_x$ is a Hochschild one-cocycle, and $x$ is a
degree-zero cocycle. Direct insertion gives
\begin{equation}\label{v4-arith:centers:eq:nonzero-bracket}
 [D,x]_{\mathrm G}=1.
\end{equation}
There are no inner derivations in this commutative algebra,
and $1$ is nonzero in its center.

\begin{proposition}[A restriction obstruction]
\label{v4-arith:centers:prop:no-e3}
The specified Hochschild $E_2$ structure on $Z_k(k[x,y,z])$
cannot be obtained by restriction of an $E_3$ structure
along the standard inclusion of little disks.
\end{proposition}

\begin{proof}
The space of two little two-disks has the homotopy type
of $S^1$. Its degree-one homology class gives the
Gerstenhaber bracket. The corresponding space for
three-disks has the homotopy type of $S^2$, whose
degree-one homology is zero. Restriction of an $E_3$
structure would therefore make that bracket zero.
Equation~\eqref{v4-arith:centers:eq:nonzero-bracket}
contradicts this consequence.
\end{proof}

A volume form can identify polyvectors with differential
forms and furnish additional operators. It does not change
the degree of this nonzero bracket. The proposition concerns
the specified $E_2$ structure and the specified restriction
map. It does not exclude every differently structured
object with the same underlying graded vector space.

\subsection{Completion in cochain length}

For an augmented differential graded algebra, write
$\bar A$ for the augmentation ideal and
$BA=\bigoplus_{q\geq0}(\bar A[1])^{\otimes q}$ for its bar
coalgebra. Its differential combines the internal
differential with multiplication, using suspended signs.
The complete Hochschild cochain complex is
\begin{equation}\label{v4-arith:centers:eq:complete}
 \operatorname{Hom}(BA,A)
 =\prod_{q\geq0}\operatorname{Hom}((\bar A[1])^{\otimes q},A).
\end{equation}
The product is forced by Hom from a direct sum.
Replacing it by a direct sum can remove cochains.

Let $\tau:BA\to A$ be the universal twisting cochain.
Convolution uses deconcatenation in $BA$ and multiplication
in $A$. The Hochschild differential is
\begin{equation}\label{v4-arith:centers:eq:twist}
 d_{\mathrm{Hoch}}=\partial+[\tau,-].
\end{equation}
The equation $\partial\tau+\tau\ast\tau=0$ implies
$d_{\mathrm{Hoch}}^2=0$. The two exterior terms in
\eqref{v4-arith:centers:eq:hoch-differential} are part of
this twist.

At each fixed input length, convolution has finitely many
cuts. The complete product therefore defines these
operations without a convergence assumption on the
sequence of cochain lengths. A completion in state weight,
pole order, or coefficient parameters is a separate
choice.

\section{Coefficient-linear Morita comparison}
\label{v4-arith:centers:sec:morita}

We now specify the structure that transports the center.
Let $(\mathcal C,\star,K)$ be a presentable stable
$k$-linear symmetric monoidal category. Assume that $\star$
preserves colimits separately. An algebra and its module
actions are unital, and relative tensors are derived.

For algebras $A,B$ in $\mathcal C$, let
$\!{}_B\operatorname{Bimod}_A(\mathcal C)$ denote the
category of objects with a left $B$-action and a commuting
right $A$-action. Write $\operatorname{Mod}_A(\mathcal C)$
for left $A$-modules. A $B$-$A$ bimodule $P$ defines
\[
 P\star_A-:
 \operatorname{Mod}_A(\mathcal C)\longrightarrow
 \operatorname{Mod}_B(\mathcal C).
\]
It preserves colimits and is $\mathcal C$-linear: tensoring
by a coefficient object commutes coherently with this
functor.

\begin{theorem}[Coefficient-linear Eilenberg--Watts comparison]
\label{v4-arith:centers:thm:watts}
The assignment $P\mapsto P\star_A-$ gives an equivalence
\begin{equation}\label{v4-arith:centers:eq:watts}
 {}_B\operatorname{Bimod}_A(\mathcal C)
 \simeq
 \operatorname{Fun}^{L}_{\mathcal C}
 \bigl(\operatorname{Mod}_A(\mathcal C),
       \operatorname{Mod}_B(\mathcal C)\bigr).
\end{equation}
For a $\mathcal C$-linear colimit-preserving functor $F$,
the inverse sends $F$ to $F(A)$ with its induced right
$A$-action. Composition corresponds to relative tensor.
\end{theorem}

\begin{proof}
The free left $A$-module on $X\in\mathcal C$ is
$A\star X$. Every module $M$ is the geometric realization
of its free simplicial bar resolution. Its terms have the
form $A\star A^{\star q}\star M$, with faces supplied by
multiplication and the action.

Put $P=F(A)$. Right multiplication on the regular module
$A$ is an $A$-module action. Applying $F$ gives the right
$A$-action on $P$. The left $B$-action is part of its
value in $\operatorname{Mod}_B(\mathcal C)$.
Coefficient linearity identifies
$F(A\star X)$ with $P\star X$.

Apply $F$ to the free simplicial resolution. Preservation
of colimits identifies its realization with $F(M)$.
The resulting simplicial object is the two-sided bar
construction computing $P\star_A M$. This gives a
natural equivalence $P\star_A M\simeq F(M)$.

A natural transformation between two such functors is
determined on the free modules by its value at $A$ and
coefficient linearity. Compatibility with right
multiplication says precisely that this value is a
bimodule map. The free resolutions extend it to all
modules, including its higher homotopies.

Finally, the two iterated bar realizations for a
composite can be realized in either order. Their common
bisimplicial object gives
$(Q\star_B P)\star_A M\simeq Q\star_B(P\star_A M)$.
Its associativity is inherited from the same multibar
construction.
\end{proof}

The identity functor corresponds to the regular
$A$-bimodule. Its endomorphisms have composition and the
coherent multiplication induced by composition of
functors. Their interaction is the $E_2$ structure on
the center. Thus the coefficient action is part of the
comparison, even when the underlying module categories
are equivalent as abstract categories.

\begin{definition}\label{v4-arith:centers:def:internal}
The \emph{internal convolution center} of $A$, denoted
$Z^{\mathrm{int}}_\star(A)$, represents the functor
\begin{equation}\label{v4-arith:centers:eq:internal}
 W\longmapsto
 \operatorname{Map}_{A\text{-}A}(W\star A,A).
\end{equation}
The bimodule structure on $W\star A$ uses the symmetry
to move $A$ past $W$. Representation means a natural
equivalence with
$\operatorname{Map}_{\mathcal C}(W,Z^{\mathrm{int}}_\star(A))$
for every $W\in\mathcal C$.
\end{definition}

The functor $W\mapsto W\star A$ preserves colimits.
Its right adjoint therefore supplies the representing
object. The center theorem gives this object its
$E_2$ structure. In Lurie's \emph{Higher Algebra},
Corollary 5.3.1.15 gives existence, and Theorem 5.3.1.30
identifies its underlying object through module maps.
For associative algebras these modules are bimodules.

The coefficient action gives the scalar center
\begin{equation}\label{v4-arith:centers:eq:internal-scalar}
 Z_k(A)=R\operatorname{Hom}_{\mathcal C}
                  (K,Z^{\mathrm{int}}_\star(A)).
\end{equation}
The internal center is an object of $\mathcal C$.
The scalar center is a complex over $k$. Evaluating at
$K$ does not identify these two recipients.

\begin{theorem}[Morita invariance with coefficients]
\label{v4-arith:centers:thm:morita}
Suppose that $P$ is a $B$-$A$ bimodule and $Q$ an
$A$-$B$ bimodule, with coherent equivalences
\[
 P\star_AQ\simeq B,\qquad Q\star_BP\simeq A.
\]
Then conjugation gives equivalences of internal and
scalar $E_2$ centers:
\[
 Z^{\mathrm{int}}_\star(A)\simeq
 Z^{\mathrm{int}}_\star(B),\qquad
 Z_k(A)\simeq Z_k(B).
\]
The equivalences depend on the supplied Morita data and
commute with their composition.
\end{theorem}

\begin{proof}
The functor
$M\mapsto P\star_AM\star_AQ$ identifies the categories
of $A$- and $B$-bimodules. It takes $A$ to $B$.
For every coefficient object $W$, symmetry and
coefficient linearity identify its value on $W\star A$
with $W\star B$. It therefore identifies the functors
represented in
\eqref{v4-arith:centers:eq:internal}.
The internal equivalence follows by Yoneda.

Conjugation also identifies composition of
coefficient-linear endofunctors. Under
\eqref{v4-arith:centers:eq:watts}, both centers are
endomorphisms of the corresponding identity functor.
Thus the equivalence preserves the coherent operations
giving the $E_2$ structure. Evaluation at $K$ gives the
scalar assertion.

For a composite Morita equivalence, both constructions
are the same iterated relative tensor. The multibar
associativity in the preceding theorem identifies them,
with its full coherence.
\end{proof}

For differential graded algebras this also has a
cochain description. Keller,
\emph{Derived invariance of higher structures on the
Hochschild complex}, Section 4.6, constructs the
comparison in the homotopy category of brace algebras.
Its hypotheses include the required fully faithful
derived tensor functors. An arbitrary homomorphism
such as Example~\ref{v4-arith:centers:ex:diagonal}
does not satisfy this conclusion.

\subsection{Matrices retain the scalar center}

Let $B=M_n(A)$. The column module $P=A^n$ is a $B$-$A$ bimodule.
Its right-$A$ dual is the row $A$-$B$ bimodule
\[
 Q=\operatorname{Hom}_{\operatorname{Mod}\text{-}A}(P,A).
\] Matrix multiplication gives
$P\otimes_AQ\simeq B$ and $Q\otimes_BP\simeq A$.
Hence $Z_k(M_n(A))\simeq Z_k(A)$ as $E_2$ algebras.

An explicit cochain comparison uses matrix paths.
For a $p$-cochain $f$ on $A$, define
\begin{equation}\label{v4-arith:centers:eq:matrix-path}
 \bigl(\widetilde f(X_1,\ldots,X_p)\bigr)_{i_0i_p}
 =\sum_{i_1,\ldots,i_{p-1}}
 f\bigl((X_1)_{i_0i_1},\ldots,(X_p)_{i_{p-1}i_p}\bigr).
\end{equation}
For $p=0$ put $\widetilde a=aI_n$.
Matrix multiplication is concatenation of these paths.

\begin{proposition}\label{v4-arith:centers:prop:matrix-path}
The matrix-path map preserves the Hochschild
differential, cup product, and all braces. It induces
the Morita equivalence on scalar centers.
\end{proposition}

\begin{proof}
Each inner term of the Hochschild differential
contracts two neighboring edges. Summing their common
index is matrix multiplication. The first and last
terms attach an edge at the corresponding endpoint.
These descriptions prove compatibility with the
differential.

The cup product splits a path into two consecutive
pieces and multiplies their values. A brace inserts
shorter paths into selected slots of a longer one.
The index sums are the same before and after either
operation. Matrix units have degree zero, so the
suspended signs are unchanged.

The corner $E_{11}BE_{11}=A$ and the two bimodules
above implement the inverse derived equivalence.
The matrix-path formula is its map on the bar
cochains. One can also compare ordinary cochains with
cochains relative to $M_n(k)$ using the separability
element
\[
 \sum_{i=1}^n E_{i1}\otimes E_{1i}.
\]
Its product is $I_n$, and it commutes with the two
actions of $M_n(k)$. Inserting it contracts the
complement of the relative bar resolution. This
identifies the displayed map with the Morita
quasi-isomorphism.
\end{proof}

No factor $1/n$ occurs. The comparison is valid even
when the coefficient characteristic divides $n$.
It concerns the scalar Hochschild center of the
matrix algebra. If a geometric algebra is required
to have scalar empty coefficient, replacing that
coefficient by $M_n(k)$ changes the geometric class.

\section{Relative scalar extension}
\label{v4-arith:centers:sec:relative}

Let $\mathcal C$ now carry a central action of
$D(k)$, compatible with its tensor. Let $A_0$ be an
algebra in $\mathcal C$. For a differential graded
$k$-algebra $B$, set
\[
 A_B=A_0\otimes_k B.
\]
For a $B'$-$B$ bimodule $P$, set
$J(P)=A_0\otimes_k P$ with its induced
$A_{B'}$-$A_B$ actions.

Assume the coefficient action has the relative-tensor
comparison
\begin{equation}\label{v4-arith:centers:eq:relative-tensor}
 (A_0\otimes Q)\star_{A_{B'}}
 (A_0\otimes P)
 \simeq A_0\otimes(Q\otimes_{B'}^{\mathbf L}P).
\end{equation}
Its coherence is part of the central coefficient
action. It can be checked on the common two-sided
bar construction.

Suppose also that empty evaluation
$e_\varnothing:\mathcal C\to D(k)$ is strong monoidal,
preserves the realizations in question, and satisfies
$e_\varnothing(A_0)=k$. It then sends $J(P)$ back to
$P$, including its two coefficient actions.

\begin{theorem}[A relative center retract]
\label{v4-arith:centers:thm:relative-retract}
Under these hypotheses there are maps of scalar
$E_2$ algebras
\begin{equation}\label{v4-arith:centers:eq:retract}
 Z_k(B)\xrightarrow{J_B}Z_{\mathcal C}(A_B)
 \xrightarrow{e_\varnothing}Z_k(B),
 \qquad e_\varnothing J_B=\mathrm{id}.
\end{equation}
Here
$Z_{\mathcal C}(A_B)=
R\operatorname{Hom}_{A_B\text{-}A_B}(A_B,A_B)$
is formed in $\mathcal C$. The statement makes
$Z_k(B)$ a retract. It does not assert that $J_B$
is essentially surjective on geometric bimodules.
\end{theorem}

\begin{proof}
The relative-tensor comparison makes $J$ a monoidal
functor on the endomorphism category of the object
$B$ in the bimodule bicategory. It sends the regular
$B$-bimodule to the regular $A_B$-bimodule.
It therefore induces a map on endomorphisms of these
units, preserving their $E_2$ operations.

Empty evaluation gives the reverse monoidal functor
on these objects. The identity $e_\varnothing J(P)=P$
is natural and compatible with composition.
Its value on unit endomorphisms proves the retract
identity.

On bar cochains the forward map multiplies the
$A_0$ inputs and applies the coefficient cochain to
the $B$ inputs. Centrality moves the coefficients
past the $A_0$ factors. The resulting Koszul signs
are exactly those in the tensor product and
suspended insertion conventions. The same formula
commutes with every cut, cup, and brace.
\end{proof}

In particular, a nonzero coefficient-center class
remains nonzero after extension. Its empty
evaluation recovers the original class. A class
supported on nonempty configurations may lie
outside the image of $J_B$.

This construction is useful for arithmetic
coefficient algebras: it places their actual
Hochschild operations in a geometric recipient.
It does not construct an identification with every
internal center, every diagonal actor, or a chosen
physical bulk algebra. Each such identification
has its own mapping property.

\section{The center of a homomorphism diagram}
\label{v4-arith:centers:sec:triangle}

The two arrows in
\eqref{v4-arith:centers:eq:mixed} have a common target.
Their homotopy fiber can therefore be formed as a
complex. To obtain its higher operations, we realize
the diagram by one algebra.

Let $f:A\to B$ be a unital homomorphism of ordinary
$k$-algebras. Regard $B$ as a $B$-$A$ bimodule by
right multiplication through $f$. Define the
\emph{triangular algebra}
\begin{equation}\label{v4-arith:centers:eq:triangle}
 T_f=
 \begin{pmatrix}A&0\\ B&B\end{pmatrix},
 \qquad
 \begin{pmatrix}a&0\\ m&b\end{pmatrix}
 \begin{pmatrix}a'&0\\ m'&b'\end{pmatrix}
 =
 \begin{pmatrix}
 aa'&0\\ mf(a')+bm'&bb'
 \end{pmatrix}.
\end{equation}
Its unit is $\operatorname{diag}(1,1)$.
The two diagonal idempotents are denoted $e_A,e_B$.

Put $u\mapsto f_*u$ for the first map of
\eqref{v4-arith:centers:eq:mixed}, and
$v\mapsto f^*v$ for the second. Define a cochain
complex $\mathcal F_f$ by
\[
 \mathcal F_f^n=C^n(A,A)\oplus C^n(B,B)
                    \oplus C^{n-1}(A,B),
\]
with differential
\begin{equation}\label{v4-arith:centers:eq:triangle-d}
 d(u,v,h)=(bu,bv,f_*u-f^*v-bh).
\end{equation}
The sign on $bh$ is the shift sign. Since both
comparison maps commute with $b$, one has $d^2=0$.

\begin{theorem}[A triangular realization]
\label{v4-arith:centers:thm:triangle}
There is a quasi-isomorphism of complexes
\[
 Z_k(T_f)\simeq\mathcal F_f.
\]
Consequently the homotopy fiber has an $E_2$ structure
through the center of $T_f$. This assertion does not
assign the ordinary brace formulas to mixed cochains.
\end{theorem}

\begin{proof}
The subalgebra $S=ke_A\oplus ke_B$ is separable over
$k$. Its separability element is
$e_A\otimes e_A+e_B\otimes e_B$.
Inserting this element gives the usual contraction
from the ordinary bar resolution to the relative
bar resolution over $S$.
Thus relative Hochschild cochains compute $Z_k(T_f)$.

An $S$-balanced tensor is a composable sequence of
matrix entries. Such a sequence either stays at one
diagonal vertex or crosses once through the
$B$-$A$ entry. It cannot cross twice, because the
opposite off-diagonal entry is zero. The two
constant-vertex parts are $C^\bullet(A,A)$ and
$C^\bullet(B,B)$.

The part with one crossing has terms
\[
 \operatorname{Hom}_k
       (B^{\otimes q}\otimes B\otimes A^{\otimes p},B),
 \qquad p,q\geq0.
\]
It occupies total Hochschild degree $p+q+1$.
Its differential is the two-sided bar differential
for the bimodule $B$, together with the two endpoint
actions. Without those endpoint terms, it computes
$R\operatorname{Hom}_{B\text{-}A}(B,B)$ shifted by
one.

Resolve the right $A$-action on $B$ by the bar
resolution, retaining its free left $B$-action.
Evaluation on $1\in B$ then identifies this derived
Hom with $C^\bullet(A,B)$.
The extra left-$B$ bar direction contracts by its
unit degeneracy. The contraction respects the
right-$A$ differential, so the same comparison
applies to the total complex.

The remaining endpoint maps apply an $A$-cochain
through $f$ and restrict a $B$-cochain along $f$.
Their orientations are opposite. With the shifted
mixed term, the total differential is exactly
\eqref{v4-arith:centers:eq:triangle-d}.
This proves the quasi-isomorphism. Transport the
$E_2$ structure from $Z_k(T_f)$ along this
equivalence. No mixed insertion with mismatched
source and target is required.
\end{proof}

The degree-zero part can be read directly from
the algebra.

\begin{proposition}\label{v4-arith:centers:prop:triangle-zero}
There is an isomorphism
\[
 Z_{\mathrm{alg}}(T_f)
 =
 \{(a,b)\in Z_{\mathrm{alg}}(A)\times
                 Z_{\mathrm{alg}}(B): f(a)=b\}.
\]
For the diagonal homomorphism of
Example~\ref{v4-arith:centers:ex:diagonal}, this
center is $k$.
\end{proposition}

\begin{proof}
Commutation with $e_A$ and $e_B$ forces the
off-diagonal entry of a central element to vanish.
Commutation with the two diagonal subalgebras
makes its remaining entries central.
Commutation with an off-diagonal element $m\in B$
is
$mf(a)=bm$ for every $m$. Taking $m=1$ gives
$f(a)=b$, and centrality of $b$ makes this condition
sufficient.

For the stated example, $b=\lambda I_2$.
The equality $\operatorname{diag}(a_1,a_2)=b$
forces $a_1=a_2=\lambda$.
\end{proof}

The triangular center thus retains the compatibility
of the two actions. Merely taking the product of
the endpoint centers would lose that condition.
Conversely, it would be incorrect to call every
ordinary fiber of cochain complexes a strict brace
algebra. The triangular model supplies the
operations that make the comparison meaningful.

\section{The regular and punctured normal lines}
\label{v4-arith:centers:sec:normal}

The simplest collision coordinate already gives a center
with a nontrivial bracket. Let $k$ have characteristic zero,
and put
\[
 D=k\langle u,\partial\rangle/(\partial u-u\partial-1),
 \qquad R=k[u],\qquad L=k[u,u^{-1}].
\]
Both $R$ and $L$ are left $D$-modules with their derivative
connections. Write $\mathcal C_2=\operatorname{thick}_D(R\oplus L)$
for the full differential graded category obtained by shifts,
finite cones, and direct summands. Its scalar center is the
derived endomorphism complex of its diagonal bimodule.
We will compute the center together with its cup product
and Gerstenhaber bracket.

For $v\in D$, let $r_v$ denote right multiplication by $v$.
This is a map of free left $D$-modules. The two resolutions
\begin{equation}\label{v4-arith:centers:eq:normal-resolutions}
 P_0=[D\xrightarrow{r_\partial}D]\longrightarrow R,
 \qquad
 P_1=[D\xrightarrow{r_{u\partial+1}}D]\longrightarrow L
\end{equation}
have terms in degrees $-1,0$ and augmentations $1,u^{-1}$.
The term $1$ in the second differential is essential:
$(u\partial+1)u^{-1}=0$, whereas $u\partial(u^{-1})=-u^{-1}$.

\begin{lemma}\label{v4-arith:centers:lem:normal-corners}
The resolutions in
\eqref{v4-arith:centers:eq:normal-resolutions} are exact.
They give
\[
\begin{array}{c|cc}
 R\operatorname{Hom}_D(-,-)&R&L\\\hline
 R&k&k\oplus k[-1]\\
 L&0&k\oplus k[-1].
\end{array}
\]
\end{lemma}

\begin{proof}
The Weyl algebra is a domain, so the two differentials
are injective. The first cokernel is the polynomial module.
In the second cokernel the Euler relation reduces mixed
monomials to powers of $u$ and powers of $\partial$.
Their images are
\[
 u^{m-1}\quad(m\geq0),\qquad
 (-1)^n n!u^{-n-1}\quad(n\geq1).
\]
These form a Laurent basis, proving the second assertion
about cokernels.

Applying Hom to the resolutions gives the operators
$\partial$ and $u\partial+1$ on $R$ and $L$.
Polynomial integration makes $\partial:R\to R$
surjective with kernel $k$. On $L$ its cokernel is the
residue line represented by $u^{-1}$. Finally,
$(u\partial+1)u^m=(m+1)u^m$. This operator is invertible
on $R$, and on $L$ it has the single exceptional mode
$m=-1$ in both kernel and cokernel.
\end{proof}

The four complexes alone do not specify their composition.
To retain composition, take the graded path algebra $T$ on two
objects. Give it an arrow $a:0\to1$ of degree zero and a loop
$\theta:1\to1$ of degree one. Set
\[
 \theta^2=0,\qquad b=\theta a,\qquad
 T^0=ke_0\oplus ke_1\oplus ka,\qquad
 T^1=k\theta\oplus kb.
\]
The differential is zero.

Products are written in
composition order. Apart from identities, the sole
nonzero product of basis elements is $\theta a=b$.

\begin{proposition}\label{v4-arith:centers:prop:normal-model}
There is a strict quasi-isomorphism
\[
 T\longrightarrow\operatorname{End}_D(P_0\oplus P_1).
\]
Consequently $\mathcal C_2$ is equivalent to
$\operatorname{Perf}(T)$, using right $T$-modules.
\end{proposition}

\begin{proof}
Send $e_i$ to the two projections. The arrow $a$ is
represented by a map $\alpha:P_0\to P_1$ with
$\alpha^{-1}=1$ and $\alpha^0=r_u$. Its chain equation
is $\partial u=u\partial+1$. Represent $\theta$ by the
degree-one identity $\Theta:P_1^{-1}\to P_1^0$ and
$b$ by $\Theta\alpha$.

These maps are closed and satisfy the multiplication
relations strictly. After augmentation, $\alpha$ is the
inclusion $R\to L$, while $\Theta$ and $\Theta\alpha$
represent the two residue classes. Thus the map induces
isomorphisms in all four Hom corners. Hom from
$P_0\oplus P_1$ and its tensor adjoint identify the
generators and their Hom complexes. The identification
extends to their thick closures.
\end{proof}

\subsection{The completed center and its bracket}
\label{v4-arith:centers:subsec:logarithmic}

Use the normalized two-sided bar relative to the separable
algebra $S=ke_0\oplus ke_1$. Its middle term is the direct
sum of finite composable words in $(T/S)[1]$. Taking Hom
therefore gives a product over arities. We denote the
resulting complete cochain algebra by $C(T)$.

Here are explicit sign conventions for the calculation.
An $n$-cochain of internal degree $r$ has total degree
$n+r$. Its differential is $D=(-1)^r b_H$, where
\[
\begin{split}
(b_Hf)(v_1,\ldots,v_{n+1})={}&
 (-1)^{r|v_1|}v_1f(v_2,\ldots,v_{n+1})\\
&+\sum_{i=1}^n(-1)^i
 f(v_1,\ldots,v_iv_{i+1},\ldots,v_{n+1})\\
&+(-1)^{n+1}f(v_1,\ldots,v_n)v_{n+1}.
\end{split}
\]
For an $m$-cochain $g$ of internal degree $s$, the cup
product is
\[
 (f\smile g)(v_1,\ldots,v_{n+m})
 =
 (-1)^{s(\sum_{i=1}^n|v_i|-n)}
 f(v_1,\ldots,v_n)g(v_{n+1},\ldots,v_{n+m}).
\]
These formulas express the suspended-bar operations on
unsuspended inputs.

\begin{theorem}\label{v4-arith:centers:thm:normal-center}
The scalar center of $\mathcal C_2$ has a continuous
associative differential graded model
\[
 k[[x]]\otimes\Lambda(\eta),\qquad |x|=0,\quad|\eta|=1,
 \qquad d=0.
\]
Its Gerstenhaber bracket is determined by
\[
 [\eta,x]=-x,\qquad [x,x]=[\eta,\eta]=0.
\]
\end{theorem}

\begin{proof}
Give $\theta,b$ weight one and $e_0,e_1,a$ weight zero.
For a cochain let $d$ be input weight minus output weight.
Every nonempty composable word has one of the forms
\[
 \theta^{\otimes n},\qquad
 \theta^{\otimes(n-1)}\otimes a,\qquad
 \theta^{\otimes(n-1)}\otimes b.
\]
The output has the same endpoints.

For each $d\geq1$,
the cochain block in total degrees $0,1,2$ is
\begin{equation}\label{v4-arith:centers:eq:normal-block}
 k^2\xrightarrow{M}k^3\xrightarrow{N}k,\qquad
 M=\begin{pmatrix}0&0\\-1&1\\-1&1\end{pmatrix},
 \quad N=\begin{pmatrix}1&1&-1\end{pmatrix}.
\end{equation}
The degree-zero basis sends $\theta^{\otimes d}$ to
$e_1$ and $\theta^{\otimes(d-1)}\otimes b$ to $a$.
The degree-one basis has outputs $\theta,a,b$ on the
corresponding words of input weights $d+1,d,d+1$.
The degree-two basis sends
$\theta^{\otimes(d+1)}\otimes a$ to $b$.
The multiplication $\theta a=b$ gives the interior
faces, and the endpoint faces give the other entries
of \eqref{v4-arith:centers:eq:normal-block}.

For $d=0$ the degree-zero basis is $e_0,e_1$ and the
first matrix is $-M$. Weight $-1$ consists of an
identity map $k\to k$ in degrees one and two. Its source
is the nullary $\theta$ and its target sends $a$ to $b$.
There are no further blocks.

For $d\geq0$ choose representatives and projections
\[
\begin{array}{ll}
 i_0(c)=(c,c),&p_0(p,q)=p,\\
 i_1(c)=(c,0,c),&p_1(r,s,t)=t-s.
\end{array}
\]
A homotopy on the complement is
\[
 h_1(r,s,t)=
 \begin{cases}(0,s),&d>0,\\(0,-s),&d=0,\end{cases}
 \qquad h_2(z)=(z,0,0).
\]
The matrices give $Dh+hD=1-ip$ and $h^2=hi=ph=0$.
Use the inverse identity in weight $-1$. The maps
preserve weight, so their product contracts the
complete cochain complex.

To identify the product, define derivations $X,Y$ by
\[
\begin{array}{c|ccccc}
 &e_0&e_1&\theta&a&b\\\hline
 X&0&0&e_1&0&a\\
 Y&0&0&\theta&0&b.
\end{array}
\]
Their internal degrees are $-1,0$ and their total
degrees are $0,1$. They are closed,
$X\smile Y=Y\smile X$, and $Y\smile Y=0$.
The cochains $X^{\smile d}$ and
$Y\smile X^{\smile d}$ give the representatives just
chosen. Thus $x\mapsto X$, $\eta\mapsto Y$ is a
continuous multiplicative quasi-isomorphism.
On each fixed word, only finitely many terms of a
formal series contribute.

The unary insertion bracket is the graded
commutator. Since $YX-XY=-X$, it gives
$[\eta,x]=-x$. The biderivation identity determines
the remaining brackets.
\end{proof}

In this description $\eta$ acts as $-x\partial_x$.
For formal series $f,g$,
\[
 [\eta f,g]=-xfg',\qquad
 [\eta f,\eta g]=\eta x(gf'-fg').
\]
The associative model and this cohomology bracket
do not assert formality of the full $E_2$ structure.
The braces still live on the complete bar cochains.
Replacing their product by a direct sum would give
$k[x]$ in place of $k[[x]]$ and would omit legitimate
maps from the bar.

Restriction to the full subcategory generated by
$L$ has target $\Lambda(\theta)[[x]]$. It is
\[
 x\longmapsto x,\qquad \eta\longmapsto\theta x.
\]
Restriction to $R$ sends $x,\eta$ to zero and
preserves constants. The regular object therefore
forces the odd coefficient at the punctured object
to vanish at the formal origin.

The same conclusion follows from the relative
bar triangle
\[
 C(T)\longrightarrow k\oplus C(\Lambda(\theta))
 \longrightarrow R\operatorname{Hom}_{\Lambda(\theta)}
          (\Lambda(\theta)a,\Lambda(\theta)a).
\]
A composable word crosses from $0$ to $1$ at most
once. Its crossing terms form the shifted module
bar. The module $\Lambda(\theta)a$ is free, so this
bar contracts to $\Lambda(\theta)$. On the diagonal
models, the last map is $(c,f)\mapsto c-f(0)$.
Its kernel is $k\oplus x\Lambda(\theta)[[x]]$,
which recovers $\eta=\theta x$.

\subsection{Removing nullary cochains}
\label{v4-arith:centers:subsec:nullary}

Deformation problems can impose conditions that exclude
the primitives used by the full center. Fix the two
objects, hence the inclusion $S\subset T$, and retain
only positive-arity cochains. These describe uncurved
$S$-linear deformations with zero nullary comparison
term. Write their complex as $C_{\geq1,S}(T)$.

\begin{proposition}\label{v4-arith:centers:prop:normal-positive}
The only nonzero groups of $C_{\geq1,S}(T)$ are
\[
 H^0=xk[[x]],\qquad
 H^1=k[[x]]\eta\oplus k\sigma,\qquad H^2=kq.
\]
Here $\sigma(a)=a$, $\sigma(b)=b$, and $q(a)=b$,
with other generator values zero.
For the normalized positive-arity complex relative
to $k$, they are instead
\[
\begin{split}
 H^0&=xk[[x]],\\
 H^1&=k[[x]]\eta\oplus k\sigma\oplus kDa,\\
 H^2&=kD\theta\oplus kDb.
\end{split}
\]
The symbol $D$ in the last expressions is the
Hochschild differential of a nullary element in
the full complex.
\end{proposition}

\begin{proof}
The positive-weight blocks do not change.
Removing the two nullary identities leaves their
former boundary $\sigma$ in degree one and removes
the constant degree-zero class. Removing the
nullary $\theta$ leaves $q$ in degree two.

For the $k$-relative complex, use the exact sequence
\[
 0\to C_{\geq1,k}(T)\to C_k(T)\to T\to0.
\]
The map on degree-zero cohomology is
$f(x)\mapsto f(0)1_T$, and the classes
$\eta k[[x]]$ have zero nullary component.
The long exact sequence contributes
$T^0/k1_T$ in degree one and $T^1$ in degree two.
Represent the former by $e_0,a$. The connecting
maps give the displayed classes. They are distinct:
$(Db)(e_0)=b$, while $(D\theta)(e_0)=0$ and
$(D\theta)(a)=b$.
\end{proof}

Let $A=k[[s]]$ and give $T\otimes_k A$ the
differential
\begin{equation}\label{v4-arith:centers:eq:normal-family}
 d_sa=sb,\qquad d_s\theta=d_sb=d_se_i=0.
\end{equation}
This is a differential graded algebra, since
$d_s=s[\theta,-]$ and $\theta^2=0$. Its underlying
graded module is free, but
\[
 H^0(T_s)=Ae_0\oplus Ae_1,\qquad
 H^1(T_s)=A\theta\oplus(A/(s))b.
\]

Even modulo $s^2$, there is no uncurved
$A_\infty$ map $T_s\to T\otimes A$ with
$F_0=0$ and $F_1=1\bmod s$. Its unary equation
would give $F_1d_s=0$, whereas its value on $a$
is $sb\ne0\bmod s^2$. Higher components do not
enter this equation. The identity $D\theta=q$
in the full center does not remove the obstruction
when the nullary term is forbidden.

\subsection{Changing the connection residue}
\label{v4-arith:centers:subsec:kummer}

The deformation has a geometric realization. Let
$D_A=A\langle u,\partial\rangle/(\partial u-u\partial-1)$,
take $R_A=A[u]$ with its derivative connection, and put
\[
 L_s=A[u,u^{-1}],\qquad
 \partial f=f'-\frac{s}{u}f.
\]
These are polynomial and Laurent polynomial modules
over $A$. There is no completion in $u$.

\begin{theorem}\label{v4-arith:centers:thm:kummer}
The category generated by $R_A\oplus L_s$ is
equivalent to $\operatorname{Perf}(T_s)$.
For $K=k((s))$,
\[
 \operatorname{Perf}(T_s\otimes_A K)
 \simeq
 \operatorname{Perf}(K)\oplus
 \operatorname{Perf}(\Lambda_K(\theta)).
\]
The family $T_s$ is not $A$-linearly Morita equivalent
to the constant family $T\otimes_k A$.
\end{theorem}

\begin{proof}
Replace $P_1$ by
\[
 P_{1,s}=[D_A\xrightarrow{r_{u\partial+1+s}}D_A],
 \qquad P_{1,s}\longrightarrow L_s,\quad1\longmapsto u^{-1}.
\]
Its cokernel calculation uses
\[
 \partial^n u^{-1}
 =(-1)^n\prod_{j=1}^n(j+s)u^{-n-1}.
\]
All $j+s$ are units of $A$, so these vectors
and the nonnegative powers give a Laurent basis.

Keep the maps $\alpha,\Theta,\Theta\alpha$
used above. Now
\[
 d_{\operatorname{Hom}}\alpha
 =r_{u\partial+1+s}-r_{\partial u}
 =s\Theta\alpha.
\]
Thus they define a strict map from $T_s$ to
the endomorphism algebra of the two resolutions.

In the corner from $R_A$ to $L_s$, all Laurent
modes except $m=0$ contract under
$u^m\mapsto(m-s)u^{m-1}$. The remaining complex
is $[A\xrightarrow{s}A]$, with the Hom sign
chosen as in \eqref{v4-arith:centers:eq:normal-family}.
The reverse corner is acyclic because $m+1+s$
is a unit for $m\geq0$. The two diagonal
corners are $A$ and $A\oplus A[-1]$.
The map is therefore a quasi-isomorphism,
and the generator argument proves the
categorical equivalence.

Over $K$, the ideal spanned by $a,b$ is
contractible. Quotienting gives the algebra
$K\times\Lambda_K(\theta)$. Its degree-zero
Hochschild center is $K\times K[[x]]$.
The constant family, computed directly over
$K$, has center $K[[x]]$.

The former ring has
a nontrivial idempotent and the latter has
none. Morita equivalence would preserve these
rings. An $A$-linear equivalence would remain
one over $K$, so it is also excluded.
\end{proof}

The ordinary inclusion $R_A\to L_s$ fails to
respect the connections by $-s/u$. The equation
$d_sa=sb$ records that failure. The proof
computed the two generic centers directly.
It did not interchange completed Hochschild
cochains with extension of scalars.

\section{Independent directions and a triple collision}
\label{v4-arith:centers:sec:three-points}

Independent normal coordinates allow the preceding
calculation to tensor. Let $D_d$ be the Weyl algebra
of $R_d=k[u_1,\ldots,u_d]$ and set
\[
 G_d=\bigoplus_{I\subseteq\{1,\ldots,d\}}
 R_d\left[\left(\prod_{i\in I}u_i\right)^{-1}\right].
\]
Every summand has its derivative connection, and
the empty product is one.

\begin{theorem}\label{v4-arith:centers:thm:crossings}
The scalar center of $\operatorname{thick}_{D_d}(G_d)$,
with product-total Hochschild cochains, has the
associative differential graded model
\[
 k[[x_1,\ldots,x_d]]
 \otimes\Lambda(\eta_1,\ldots,\eta_d),\qquad d=0,
\]
where $|x_i|=0$, $|\eta_i|=1$. The cohomology
bracket is
\[
 [\eta_i,x_j]=-\delta_{ij}x_j,\qquad
 [\eta_i,\eta_j]=[x_i,x_j]=0.
\]
\end{theorem}

\begin{proof}
The generator is the external tensor product of
the $d$ modules $k[u_i]\oplus k[u_i,u_i^{-1}]$.
Tensor the bounded free resolutions and their
strict endomorphism models. They identify the
generated category with $\operatorname{Perf}(T^{\otimes d})$.

Resolve the diagonal bimodule by the tensor product
of the relative bars. At each multiweight the
cochain spaces are finite dimensional and bounded.
Their contractions tensor: for two factors the
homotopy is
$h_1\otimes1+i_1p_1\otimes h_2$, with the graded
tensor sign. Taking the product over multiweights
gives precisely the displayed formal series.

The multiplicative representatives tensor as well.
The shuffle comparison with the bar of
$T^{\otimes d}$ identifies insertions up to
coherent homotopy. The bracket is the sum of the
factor brackets, with zero brackets between
different factors.
\end{proof}

This theorem applies to two independent pair
collisions, for example $12|34$ with constant
tangential coefficients. At a triple collision
the equations $u$, $v$, and $u-v$ are dependent.
There is an additional shape coordinate.

\subsection{The separated shape}
\label{v4-arith:centers:subsec:shape}

Write three distinct points as
\[
 z_3=z,\qquad z_1=z+u,\qquad z_2=z+ut.
\]
Their configuration space is
$\mathbb A^1_z\times\mathbb G_{m,u}\times V$, where
\[
 V=\operatorname{Spec}k[t,t^{-1},(t-1)^{-1}].
\]
The coordinate $u$ measures scale, and $t$
retains the two finite collision channels.

The algebraic de Rham algebra of $V$ has a
strict multiplicative model
\[
 B=k\oplus k\epsilon_0\oplus k\epsilon_1,\qquad
 |\epsilon_i|=1,\qquad\epsilon_i\epsilon_j=0,
\]
via $\epsilon_0\mapsto dt/t$ and
$\epsilon_1\mapsto dt/(t-1)$. Partial fractions
integrate all polynomial terms and all poles of
order at least two. The two simple residues
remain independent. This proves the
quasi-isomorphism. The scale factor contributes
$\Lambda(du/u)$, giving de Rham dimensions
$1,3,2$.

\begin{proposition}\label{v4-arith:centers:prop:shape}
Let $F=k\langle\!\langle x,y\rangle\!\rangle$
be the completed free associative algebra.
The complete normalized Hochschild cochains of
$B$ identify with
\[
 F\xrightarrow{d}F\oplus F,\qquad
 df=([x,f],[y,f]),
\]
in degrees zero and one. The degree-one ideal
has square zero and componentwise left and
right $F$-actions. Consequently
\[
 H^0C(B)=k,\qquad
 H^1C(B)=
 (F\oplus F)/\{([x,f],[y,f]):f\in F\}.
\]
There is no other cohomology. The component
of $H^1$ of output word length $m$ has
dimension
\[
 2\quad(m=0),\qquad4\quad(m=1),\qquad
 3\cdot2^{m-1}\quad(m\geq2).
\]
\end{proposition}

\begin{proof}
The two suspended augmentation generators have
degree zero. Scalar outputs give the series
$F$ in total degree zero, and the two odd
outputs give $F\oplus F$ in degree one.
All interior bar faces vanish. The endpoint
faces are the two commutators, while the cup
product concatenates words and multiplies
outputs.

A homogeneous word polynomial commuting with
$x$ is a polynomial in $x$. Comparing $xf$ with $fx$ first excludes words ending in $y$.
Moving each final $x$ to the front then excludes every occurrence
of $y$. Commutation with $y$ leaves only
constants. Applying this in each length
gives $Z(F)=k$.

At output length $m\geq2$, the differential
is injective on scalar words of length
$m-1$, hence has rank $2^{m-1}$.
Its rank is zero for $m=1$, and it has no
source for $m=0$. Subtract these ranks from
$2\cdot2^m$. All maps respect the finite
length components, so the image is the
product of its images in these components.
The computation therefore respects the
completion.
\end{proof}

This is the center of the separated shape
corner. Identifying the full triple coefficient
category requires its remaining objects and
collision maps.

\subsection{A class invisible to every pair}
\label{v4-arith:centers:subsec:resonance}

The difference also appears in ordinary global
de Rham cohomology. For charges $q_i\in k$,
consider the connection
\[
 \nabla_n=d+c\sum_{i<j}q_iq_j\,d\log(z_i-z_j)
\]
on the distinct $n$-point configuration.
After separating a scale coordinate $u$, its
radial residue is
$\lambda_n=c\sum_{i<j}q_iq_j$.
The differential on the Laurent mode $u^m$
is multiplication by $m+\lambda_n$, followed
by $du/u$. If $\lambda_n\notin\mathbb Z$,
division by these scalars contracts the
radial complex.

\begin{proposition}\label{v4-arith:centers:prop:thirds}
For $c=1/3$ and $q_1=q_2=q_3=1$, all two-point
global algebraic de Rham complexes are acyclic.
The three-point complex has exactly
\[
 H^1=k\left[u^{-1}\frac{dt}{t}\right],\qquad
 H^2=k\left[u^{-1}\frac{du}{u}\wedge\frac{dt}{t}\right].
\]
\end{proposition}

\begin{proof}
The pair residue is $1/3$ and the triple residue
is one. In the coordinates above, the connection
is
\[
 d+\frac{du}{u}
 +\frac13\left(\frac{dt}{t}+\frac{dt}{t-1}\right).
\]
Let
$D_{\mathrm{sh}}f=f'+(1/(3t)+1/(3(t-1)))f$.

A pole of order $j\geq2$ at $p=0,1$ with
leading coefficient $a$ is removed by
subtracting
\[
 D_{\mathrm{sh}}\left(
 \frac{a(t-p)^{1-j}}{1/3-(j-1)}\right).
\]
This introduces at most a simple pole at the
other point. Remove a polynomial leading
term $at^j$ with
$D_{\mathrm{sh}}(at^{j+1}/(j+1+2/3))$.

The remainder is $A_0/t+A_1/(t-1)$.
Subtracting $D_{\mathrm{sh}}(3A_1)$ leaves
$(A_0-A_1)/t$. For example,
\begin{equation}\label{v4-arith:centers:eq:shape-identity}
 \frac1{t^2(t-1)^2}
 =
 D_{\mathrm{sh}}\left(\frac{3-6t^2}{2t(t-1)}\right)
 +\frac2t.
\end{equation}

There is no nonzero rational horizontal function:
its order at zero would have to be $-1/3$.
Nor is $1/t$ in the image. A primitive with
a pole would produce a pole of higher order
with nonzero coefficient. It must therefore
be polynomial. A positive degree would produce
a polynomial part, and a constant gives equal
coefficients at the two simple poles.
Thus the shape cohomology is exactly
$k\,dt/t$ in degree one.

The radial contraction removes all modes
except $u^{-1}$ and $u^{-1}du/u$.
Polynomial integration contracts the
translation coordinate. Tensoring the
three contractions gives the stated result.
\end{proof}

The exceptional integral residues can all be
included in the calculation. Let
$a_{12},a_{13},a_{23}\in k$ and
$\lambda=a_{12}+a_{13}+a_{23}$. All functions
and forms in the following statement have
finite Laurent support.

\begin{proposition}\label{v4-arith:centers:prop:residue-strata}
For the connection
$d+\sum_{i<j}a_{ij}\,d\log(z_i-z_j)$, the
global de Rham dimensions are
\[
\begin{array}{c|ccc}
 &H^0&H^1&H^2\\\hline
 \lambda\notin\mathbb Z&0&0&0\\
 \lambda\in\mathbb Z,\ (a_{12},a_{13},a_{23})
                  \notin\mathbb Z^3&0&1&1\\
 (a_{12},a_{13},a_{23})\in\mathbb Z^3&1&3&2.
\end{array}
\]
All other groups vanish.
\end{proposition}

\begin{proof}
The scale differential is $m+\lambda$ on
$u^m$. When $\lambda$ is integral it retains
one class in degrees zero and one. The shape
operator is
\[
 Df=\left(f'+\left(\frac a t+\frac b{t-1}\right)f\right)dt,
 \qquad a=a_{23},\quad b=a_{12}.
\]

Choose an integer $M\geq1$ such that
$a-j$, $b-j$, and $j+a+b$ are nonzero for
every integer $j>M$. Let $V_M$ consist of
functions with poles of order at most $M$
at each puncture and polynomial degree at
most $M$. Let $W_M$ consist of one-forms
with pole orders at most $M+1$ and
polynomial degree at most $M-1$.
Then
\[
 D(V_M)\subset W_M,\qquad
 \dim V_M=3M+1,\qquad \dim W_M=3M+2.
\]

The leading coefficients just listed remove
every pole or polynomial term outside this
finite complex. Conversely, a function outside
$V_M$ has a leading term whose derivative
lies outside $W_M$ with nonzero coefficient.
Thus the finite inclusion induces an
isomorphism on cohomology.

A nonzero horizontal rational function has
orders $-a,-b$ at the two punctures, so it
exists exactly when both residues are
integral. In that case it is a scalar
multiple of $t^{-a}(t-1)^{-b}$.
The dimension difference gives shape
dimensions $(1,2)$ in the integral case
and $(0,1)$ otherwise. Tensor with the
two scale classes. For integral $\lambda$,
integrality of $a_{12},a_{23}$ is equivalent
to integrality of all three residues.
\end{proof}

Pairwise global cohomology therefore does not
recover the complete coefficient diagram.
The example does not negate factorization
at the coefficient level. Its conclusion
concerns the loss caused by taking global
cohomology separately on pairs.

\subsection{The relative residue class}
\label{v4-arith:centers:subsec:relative-residue}

Localization also gives a relative complex
\[
 \mathcal R=\operatorname{fib}\left(
 R\operatorname{Hom}_D(R,R)\longrightarrow
 R\operatorname{Hom}_D(R,L)\right).
\]
It records a map to the regular coefficient
together with a nullhomotopy after localization.
The constant terms cancel, while the residue
line survives one degree higher:
\[
 \mathcal R\simeq k[-2].
\]
In $\operatorname{Cone}(k\to\Omega^\bullet(L))[-1]$
the generator is $(0,du/u)$. It cannot be a
boundary because no Laurent polynomial
has derivative $u^{-1}$.

Thus a relative map can be nonzero even
when forgetting its specified nullhomotopy
kills the underlying map. An application
to pointed actions must supply its actual
input maps into $\mathcal R$ and their
normalizations. The one-dimensional
cohomology does not choose them.

\section{Centers over a parameter ring}
\label{v4-arith:centers:sec:parameters}

A finite bar model makes coefficient change especially
explicit. It also separates the cup product from the
insertion bracket. Let
\[
 R=k[h_1,\ldots,h_r],\qquad I=(h_1,\ldots,h_r),
 \qquad S=R/I,
\]
where $r\geq1$ and $k$ has characteristic zero.
Let $W$ have basis $u_1,\ldots,u_r$, all in
cohomological degree one. Write $L_i$ for exterior
multiplication by $u_i$, and form the complex
\begin{equation}\label{v4-arith:centers:eq:parameter-arrow}
 M=R\otimes_k\Lambda W,\qquad d_M=\sum_i h_iL_i.
\end{equation}

The anticommutation of the $L_i$ gives $d_M^2=0$.
The complex $M$ is an arrow complex. Its differential
is not a derivation of the auxiliary exterior
multiplication. Its cohomology is $S$ in degree $r$.

Take the triangular differential graded algebra
\[
 T_R=R\varepsilon_0\oplus R\varepsilon_1\oplus M,
 \qquad \varepsilon_1M\varepsilon_0=M,\qquad M^2=0.
\]
We compute its relative scalar center
$Z_R(T_R)=R\operatorname{Hom}_{T_R\otimes_R T_R^{\mathrm{op}}}
(T_R,T_R)$. All cochains here are $R$-linear.

\subsection{A finite model with both operations}
\label{v4-arith:centers:subsec:finite-parameter}

Put $E=\operatorname{End}_R(M)$, with differential
$\delta f=d_Mf-(-1)^{|f|}fd_M$. The normalized
bar relative to
$S_0=R\varepsilon_0\oplus R\varepsilon_1$
has just two terms,
\[
 T_R\otimes_{S_0}T_R,\qquad
 T_R\otimes_{S_0}M[1]\otimes_{S_0}T_R.
\]
There are no longer terms because
$M\otimes_{S_0}M=0$. Reversing the sign of its
one-arrow generator gives the finite free
cochain model
\begin{equation}\label{v4-arith:centers:eq:parameter-complex}
\begin{split}
 Z^n&=(R^2)^n\oplus E^{n-1},\\
 d(a,b;f)&=(0,0;(b-a)\operatorname{id}-\delta f).
\end{split}
\end{equation}
The endpoint summands occur only in degree zero.
The model has total rank $2+4^r$ over $R$.

Use the source-first opposite-cup convention.
The cup product and unit are
\begin{equation}\label{v4-arith:centers:eq:parameter-cup}
 (a,b;f)(a',b';f')=(aa',bb';fb'+af'),
 \qquad 1=(1,1;0).
\end{equation}
The differential identity on endpoints is
$bb'-aa'=(b-a)b'+a(b'-a')$.
Together with $R$-linearity of $\delta$, it
proves the differential Leibniz rule.
Associativity follows from the commuting
endpoint scalar actions.

Let $\sigma f=(0,0;f)$, of degree $|f|+1$.
The insertion convention is
\begin{equation}\label{v4-arith:centers:eq:parameter-bracket}
 [\sigma f,\sigma g]
 =\sigma\bigl(fg-(-1)^{|f||g|}gf\bigr).
\end{equation}
Endpoint scalars have zero bracket with
one-arrow cochains.

These formulas follow
by inserting an arrow operation into
another arrow operation. They give a
bracket of degree $-1$. Its differential
and Jacobi identities reduce to the
graded commutator identities in $E$.
In particular, composition in $E$ remains
visible even though two pure one-arrow
cochains have zero cup product.

\subsection{Integral cohomology and contraction brackets}
\label{v4-arith:centers:subsec:integral-parameter}

Let $\iota_i$ be contraction by the dual of
$u_i$. The operators satisfy
\[
 L_iL_j+L_jL_i=0,\qquad
 \iota_i\iota_j+\iota_j\iota_i=0,\qquad
 \iota_iL_j+L_j\iota_i=\delta_{ij}\operatorname{id}.
\]
The ordered operators $L_A\iota_B$, for
subsets $A,B\subseteq\{1,\ldots,r\}$, form
an $R$-basis of $E$. For one generator,
$1,L,\iota,L\iota$ span all two-by-two
matrices. Graded tensor products give
the general assertion.

Moreover,
\[
 \delta L_i=0,\qquad
 \delta\iota_i=h_i\operatorname{id}.
\]
For fixed $L_A$, the contraction factors
therefore form a Koszul resolution. Hence
\begin{equation}\label{v4-arith:centers:eq:end-parameter}
 H^q(E)=S\otimes_k\Lambda^qW\quad(0\leq q\leq r),
\end{equation}
with no other cohomology.

\begin{theorem}\label{v4-arith:centers:thm:integral-parameter}
The relative Hochschild groups are
\[
\begin{split}
 \operatorname{HH}^0_R(T_R)&=R\times_S R,\\
 \operatorname{HH}^1_R(T_R)&=0,\\
 \operatorname{HH}^{q+1}_R(T_R)&=S\otimes_k\Lambda^qW
       \quad(1\leq q\leq r).
\end{split}
\]
They vanish in all other degrees. The
positive-degree ideal has square-zero cup
product, and a degree-zero pair acts through
its common residue in $S$.

Writing $t_i=[(0,h_i;\iota_i)]$, the
degree-zero ring has presentation
\[
 R[t_1,\ldots,t_r]/
 (h_it_j-h_jt_i,\ t_it_j-h_it_j:1\leq i,j\leq r).
\]
\end{theorem}

\begin{proof}
The fiber complex
\eqref{v4-arith:centers:eq:parameter-complex}
gives the map
$R^2\to H^0(E)=S$, $(a,b)\mapsto b-a$.
It is surjective with kernel $R\times_S R$.
Its remaining cohomology is the positive
cohomology of $E$ shifted by one.
The cup formula proves the product
statements.

Each representative $t_i$ is closed and
maps to $(0,h_i)$. Every element of the
fiber product is a diagonal scalar plus
$(0,c)$ with $c\in I$. Thus these elements
generate. The quadratic relations reduce
a polynomial to $a+\sum_i b_it_i$.
It maps to zero exactly when $a=0$ and
$\sum_i b_ih_i=0$. The Koszul syzygies
of the regular sequence generate these
relations, giving the presentation.
\end{proof}

The exterior powers in the theorem describe
modules of classes. Their exterior product
is not the cup product in the center.
The proof works for any regular sequence
in a commutative $\mathbb Q$-algebra, with
the corresponding free exterior module.

For $A=\{a_1<\cdots<a_q\}$, let
$\alpha_A=[\sigma L_{a_1}\cdots L_{a_q}]$.
It has degree $q+1$. The contraction
relations give, for $q\geq2$,
\begin{equation}\label{v4-arith:centers:eq:contraction-bracket}
 [t_i,\alpha_A]=
 \begin{cases}
 (-1)^{j-1}\alpha_{A\setminus\{i\}},&i=a_j,\\
 0,&i\notin A.
 \end{cases}
\end{equation}
For $q=1$ the result is zero in cohomology:
the commutator is a scalar identity, which
is the boundary of an endpoint difference.
The $L_i$ supercommute, as do the
$\iota_i$, so brackets between two positive
classes or two degree-zero classes vanish.
The bracket with $R$ also vanishes.

For two parameters, put
$\omega=\alpha_{\{1,2\}}$. Then
\[
 [t_1,\omega]=\alpha_{\{2\}},\qquad
 [t_2,\omega]=-\alpha_{\{1\}}.
\]
The positive classes have zero mutual
cup products, and the $t_i$ annihilate
them under cup product. The bracket still
detects the contraction of an arrow
operation.

\subsection{Derived specialization}
\label{v4-arith:centers:subsec:specialization}

The finite model gives an operation-preserving
base-change theorem. For every commutative
derived $R$-algebra $B$,
\begin{equation}\label{v4-arith:centers:eq:parameter-base}
 Z_R(T_R)\otimes_R^{\mathbb L}B
 \simeq
 Z_B(T_R\otimes_R^{\mathbb L}B).
\end{equation}
Indeed, both terms of the bimodule bar are
induced from bounded finite free
$R$-complexes. Tensor--Hom base change
holds on each term and commutes with
the bar differential, evaluation, and
insertions. If one $h_i$ is inverted,
$h_i^{-1}\iota_i$ contracts $M$. The
center then reduces to the two scalar
endpoint sectors.

At the common zero of the parameters,
put $V=\Lambda W$ and $T_0=T_R\otimes_R k$.
The arrow differential becomes zero.

\begin{theorem}\label{v4-arith:centers:thm:singular-fiber}
The entire singular-fiber center has
cohomology
\[
 \operatorname{HH}^*_k(T_0)
 =
 k\ltimes
 \bigl(\operatorname{End}_k(V)/k\operatorname{id}\bigr)[-1].
\]
The ideal has square-zero cup product and
the shifted graded-commutator bracket.
The unit has zero bracket. Its dimensions
are
\[
 \dim_k\operatorname{HH}^n_k(T_0)=
 \begin{cases}
 1+\binom{2r}{r-1},&n=0,\\
 \binom{2r}{r}-1,&n=1,\\
 \binom{2r}{r+n-1},&\text{otherwise},
 \end{cases}
\]
where out-of-range binomial coefficients
are zero. The total dimension is $4^r$.
\end{theorem}

\begin{proof}
In \eqref{v4-arith:centers:eq:parameter-complex},
the remaining differential sends the
endpoint difference to $\operatorname{id}$,
now in center degree one. The diagonal
scalar is the unit. All other classes
are one-arrow cochains modulo this
identity. The formulas for cup products
and brackets descend to the stated
operations.

For internal degree $p$,
\[
 \dim_k\operatorname{End}^p(\Lambda W)
 =\sum_j\binom rj\binom r{j+p}
 =\binom{2r}{r+p}.
\]
Shift by one, remove the scalar identity
in degree one, and add the unit in
degree zero. This gives the formula.
\end{proof}

For $r=2$, the dimensions in degrees
$-1,0,1,2,3$ are $1,5,5,4,1$.
Ordinary specialization of the integral
degree-zero ring has dimension only
three, with basis $1,t_1,t_2$.
Thus \eqref{v4-arith:centers:eq:parameter-base}
requires derived specialization before
cohomology.

There is also no $R$-linear deformation
retraction of the displayed finite
complex onto its integral cohomology
with zero differential. The nonzero
group $\operatorname{HH}^2_R(T_R)$ is
annihilated by $I$, while $Z^2$ is
free. Every $R$-linear map from that
group into $Z^2$ is zero, since $h_1$
acts injectively on $Z^2$. Such a map
cannot have a left inverse.
This excludes the specified contraction,
without deciding all possible formality
zigzags.

\subsection{A nonzero bracket obstruction}
\label{v4-arith:centers:subsec:parameter-obstruction}

Take $r=2$ and the basis
$v_0=1,v_1=u_1,v_2=u_2,v_3=u_1u_2$,
of degrees $0,1,1,2$. Define a degree-one
endomorphism by
\[
 f(v_0)=v_1,\qquad f(v_1)=v_3,\qquad
 f(v_2)=f(v_3)=0.
\]
The arrow differential $\epsilon f$
defines a first-order deformation over
$k[\epsilon]/(\epsilon^2)$.
Its square vanishes modulo $\epsilon^2$,
and no two arrows compose.

\begin{proposition}\label{v4-arith:centers:prop:parameter-obstruction}
This first-order deformation has no
extension over $k[\epsilon]/(\epsilon^3)$
on the same graded module, even with
uncurved $A_\infty$ corrections having
the same first-order Hochschild class.
The obstruction is
\[
 \frac12[\sigma f,\sigma f]=\sigma f^2
 \ne0\quad\text{in }\operatorname{HH}^3_k(T_0).
\]
Its cup square is zero.
\end{proposition}

\begin{proof}
The endomorphism square is nonzero:
$f^2(v_0)=v_3$. Adding $\epsilon^2g$
to the arrow differential leaves its
square $\epsilon^2f^2$ modulo
$\epsilon^3$.

Allowing higher-arity operations gives
the second-order equation
\[
 d_{\mathrm{Hoch}}g+
 \tfrac12[\sigma f,\sigma f]=0.
\]
The class of $\sigma f^2$ is nonzero
in Theorem~\ref{v4-arith:centers:thm:singular-fiber}:
its internal degree is two, so it
cannot be a scalar identity.
The obstruction class depends only
on the first-order cohomology class.
No such correction exists. The cup
square vanishes by
\eqref{v4-arith:centers:eq:parameter-cup}.
\end{proof}

The original Koszul directions do integrate:
$(xL_1+yL_2)^2=0$ over $k[x,y]$.
The insertion bracket distinguishes these
directions from $f$. Under the scalar
extension and empty-coefficient retraction
constructed earlier, the nonzero obstruction
remains nonzero in the scalar center of
the extended algebra. Applying the
retraction recovers it, including its
bracket. Interpreting that image as a
diagonal action still requires a map
between the two action problems.

\section{Duality changes the variance of a center}
\label{v4-arith:centers:sec:polarity}

Return to the ordered collision category
$(\mathcal C,\star,K)$. Its convolution combines
coefficients on disjoint subsets, and $K$ is
supported at the empty configuration. Let $A$
be an augmented associative algebra with
$A_\varnothing=k$, and let $C=BA$.
The full coefficient duality of
Chapter~\ref{v4-arith:md:chapter} is an
anti-equivalence
\[
 (-)^\sharp:\mathcal C^{\mathrm{op}}\simeq\mathcal C^\sharp,
 \qquad
 (V\star W)^\sharp\simeq
 V^\sharp\star_\sharp W^\sharp.
\]
It retains the entire pro-perfect diagram,
including its transition maps. Write
$\mathsf S$ for the inverse anti-equivalence
and $E=C^\sharp$ for the transposed bar algebra.

A central operation is a bimodule map
$V\star A\to A$. The coefficient-linear
bar comparison carries it to
$V\star C\to C$ in bicomodules.
Transposition gives
\[
 E\longrightarrow V^\sharp\star_\sharp E.
\]
Thus the coefficient has moved from the
source to the target. This is the variance
of a central coaction.

\begin{theorem}\label{v4-arith:centers:thm:coaction}
Let $Z^{\mathrm{int}}(A)$ be the internal
center and put $Q_A=Z^{\mathrm{int}}(A)^\sharp$.
For every $W\in\mathcal C^\sharp$,
\[
 R\operatorname{Hom}_{\mathcal C^\sharp}(Q_A,W)
 \simeq
 R\operatorname{Hom}_{E\text{-}E}
       (E,W\star_\sharp E).
\]
The object $Q_A$ has the counital coalgebra
structure dual to composition in the center.
Its universal map is a coherent coaction
in $E$-bimodules,
\[
 \gamma_A:E\longrightarrow Q_A\star_\sharp E.
\]
\end{theorem}

\begin{proof}
The three natural equivalences are
\[
\begin{split}
 R\operatorname{Hom}(Q_A,W)
 &\simeq
 R\operatorname{Hom}_{\mathcal C}
       (\mathsf SW,Z^{\mathrm{int}}(A))\\
 &\simeq
 R\operatorname{Hom}^{C\text{-}C}
       ((\mathsf SW)\star C,C)\\
 &\simeq
 R\operatorname{Hom}_{E\text{-}E}
       (E,W\star_\sharp E).
\end{split}
\]
The middle step is the coefficient-linear
bar comparison, and the last transposes
the actual bicomodule maps.
The identity of $Q_A$ gives $\gamma_A$.
Naturality says that the map associated
with $q:Q_A\to W$ is
$(q\star_\sharp1_E)\gamma_A$.

Composition makes the evaluation of the
internal center an associative unital
action. Transpose its identities and
coherent homotopies. They become the
coassociativity and counit identities
of $\gamma_A$. Every comparison respects
the original collision diagram.
\end{proof}

An arbitrary map $q:Q_A\to W$ gives a
bimodule map with coefficient $W$.
Suppose $W$ is counital. The displayed map is a counital coaction
precisely when $q$ is a coalgebra map. Explicitly, this requires
\[
 \epsilon_Wq=\epsilon_{Q_A},\qquad
 \Delta_Wq=(q\star_\sharp q)\Delta_{Q_A},
\]
together with their higher coherences.
Apply the natural equivalence of the theorem
to $K$, $W\star_\sharp W$, and all higher
tensor powers to obtain this equivalence
of coherent structures. The zero bimodule
map, for example, fails the counit identity
when $E\ne0$.

An internal center of $E$ would represent
the different functor
$W\mapsto R\operatorname{Hom}_{E\text{-}E}
(W\star_\sharp E,E)$.
The existence of $Q_A$ does not settle
that representability question.

\subsection{Scalar centers and reflection}
\label{v4-arith:centers:subsec:reflection}

Scalarization tests only the tensor unit.
For bounded holonomic coefficients,
Verdier duality gives the algebra
$\mathsf K_*A=\mathbb D BA$ for the
transported product $\star_!$.
There is an equivalence
\[
 Z^{\mathrm{sc}}(A)
 \simeq
 Z^{\mathrm{sc}}(\mathsf K_*A)^\sigma
\]
of $E_2$-algebras, where $\sigma$ reflects
one multiplication direction.

To see the reflection, first use the
coefficient-linear equivalence between
bimodules over $A$ and bicomodules over
$BA$. The regular bimodule becomes the
regular bicomodule. The derived cotensor
of two bicomodules is computed by
normalized terms
$P\star\overline C^{\star q}\star Q$.
At a configuration $I$, positivity of
$\overline C$ leaves only $q\leq|I|$.

Verdier duality therefore transposes a
finite limit into the finite realization
that computes the balanced tensor
product over $\mathbb DC$. It preserves
the units and iterated interface
compositions with their coherences.
It reverses the vertical composition
of maps. The two compositions that
give the $E_2$ structure consequently
have one direction reflected.

On cohomology this preserves the cup
product and reverses the Gerstenhaber
bracket. Multiplication by $(-1)^n$ in
degree $n$ identifies the resulting
graded Gerstenhaber algebras abstractly.
That observation alone supplies no
chain-level $E_2$ comparison without
the reflection. Internal polarity,
scalar reflection, and opposite
algebra formation remain separately
typed constructions.

\subsection{A point-supported test}
\label{v4-arith:centers:subsec:point-test}

Fix a point $x$ of the curve. Put the
same differential-module point direct
image on each positive ordered chart,
at the tuple $(x,\ldots,x)$, with its
Kashiwara collision identifications.
Such systems, together with an arbitrary
empty coefficient, form a category
$D(k)\times D(k)$. Write an object as
$(U_0,U_1)$. Its tensor is
\[
 (U_0,U_1)\star(V_0,V_1)
 =
 \bigl(U_0\otimes V_0,\,
 (U_0\otimes V_1)\oplus(U_1\otimes V_0)\bigr).
\]

Two positive factors have zero product:
the disjoint open misses their common
support. The remaining collision
identifications have no extra mapping
data, since positive ordered sets and
monotone surjections have the singleton
as terminal object.

For a vector space $M$, take
\[
 A_M=(k,M[-1]),\qquad
 E_M=(BA_M)^\sharp=(k,M^\sharp).
\]
The bar has only the empty and singleton
terms. Here $M^\sharp$ is the pro-object
of duals of finite-dimensional subspaces
of $M$.

\begin{lemma}\label{v4-arith:centers:lem:point-test}
For $W(V)=(V^\sharp,0)$, there is a
natural equivalence
\[
 R\operatorname{Hom}_{E_M\text{-}E_M}
       (W(V)\star_\sharp E_M,E_M)
 \simeq
 V\oplus\operatorname{Hom}_k(M,V\otimes M)[-1].
\]
For vector spaces $V,M$, these summands
lie in degrees zero and one with zero
differential.
\end{lemma}

\begin{proof}
Let $N=(0,M^\sharp)$. Since
$N\star_\sharp N=0$, the bimodule bar
has two free terms. Its positive part
is the split exact sequence
\[
 M^\sharp\xrightarrow{n\mapsto(n,-n)}
 M^\sharp\oplus M^\sharp
 \xrightarrow{(a,b)\mapsto a+b}M^\sharp.
\]
The first free term has resolution
degree $-1$.

Tensoring by $W(V)$ and
applying free bimodule adjunction
gives the two coefficient mapping
objects $V$ and
$\operatorname{Hom}_k(M,V\otimes M)$.
The left and right scalar actions
agree, so the Hom differential is
zero. The resolution degree places
the second summand in degree one.
\end{proof}

For $M=k$, all relevant internal
objects exist within these point
categories, but they have different
coefficients:
\[
\begin{array}{c|c}
 Z^{\mathrm{int}}(A_k)&(k\oplus k[-1],k[-1])\\
 Q_{A_k}&(k\oplus k[1],k[1])\\
 Z^{\mathrm{int}}(E_k)&(k\oplus k[-1],k).
\end{array}
\]
The positive test has no additional
bar term, while the empty test acquires
the shifted endomorphism term.
Duality reverses its shift. Maps from
the unit to the center and maps from
the coHom to the unit still give the
same scalar complex $k\oplus k[-1]$.

\begin{proposition}\label{v4-arith:centers:prop:point-obstruction}
For $M=\bigoplus_{r\geq0}ke_r$, the
algebra $E_M$ has no internal center
representing all pro-perfect
coefficient tests. Its universal
central coaction still exists.
\end{proposition}

\begin{proof}
A representing object would turn the
restricted tests into
$R\operatorname{Hom}_k(T,V)$, where
$T$ is its transposed empty coefficient.
This functor preserves products in
$V$. So would its degree-one
cohomology.

The product comparison in the lemma
would then be surjective:
\[
 \operatorname{Hom}_k
 \left(M,\left(\prod_{r\geq0}k\right)\otimes M\right)
 \longrightarrow
 \prod_{r\geq0}\operatorname{End}_k(M).
\]
It misses the family
$f_r(e_0)=e_r$, $f_r(e_j)=0$ for $j>0$.
The value at $e_0$ on the left is a
finite tensor sum. All its coordinates
lie in one finite-dimensional subspace
of $M$, whereas the $e_r$ do not.
This contradicts product preservation.
The coaction follows independently
from Theorem~\ref{v4-arith:centers:thm:coaction}.
\end{proof}

\section{A holonomic representing kernel}
\label{v4-arith:centers:sec:holonomic}

Finite-dimensional duality suggests a positive
criterion:
$\operatorname{Hom}(M,V\otimes N)
\simeq\operatorname{Hom}(M\otimes N^*,V)$.
The output $N$ is dualized, while the
test $V$ remains arbitrary. For holonomic
differential modules, direct image and
Verdier duality give the corresponding
relative kernel.

Work over $\mathbb C$ on smooth schemes
of finite type. Let
$\mathsf H(Y)=D^b_{\mathrm{hol}}(\mathcal D_Y)$.
Use
$P\otimes^!Q=\Delta_Y^!(P\boxtimes Q)$,
with its density and shift conventions.
The two required identities are
\[
\begin{split}
 R\operatorname{Hom}_{\mathcal D_Y}(M,L)
 &\simeq
 R\Gamma_{\mathrm{dR}}(Y,\mathbb D_YM\otimes^!L),\\
 p_*(P\otimes^!p^!L)&\simeq p_*P\otimes^!L.
\end{split}
\]
The first allows bounded coherent $M$
and arbitrary $L$. The second is the
projection formula for quasi-compact
morphisms.

These are the Verdier
pairing and projection formula of
Drinfeld--Gaitsgory,
\emph{On some finiteness questions for
algebraic stacks}, Sections 5.2.6 and
5.3.1. We use preservation of bounded
holonomicity under inverse image,
direct image, and duality. On affine
space the first two are Mili\v{c}i\'c,
\emph{Lectures on Algebraic Theory of
$D$-Modules}, Chapter I, Theorems 14.6
and 14.7.

\begin{lemma}\label{v4-arith:centers:lem:holonomic-kernel}
For $p:U\to Y$, $q:U\to Z$, and
$M\in\mathsf H(U)$, $N\in\mathsf H(Z)$,
put
\[
 P_{p,q}(M,N)
 =\mathbb D_Yp_*
       (\mathbb D_UM\otimes^!q^!N).
\]
This object is bounded holonomic and,
for every $V\in\operatorname{DMod}(Y)$,
\[
 R\operatorname{Hom}_{\mathcal D_U}
       (M,p^!V\otimes^!q^!N)
 \simeq
 R\operatorname{Hom}_{\mathcal D_Y}
       (P_{p,q}(M,N),V).
\]
\end{lemma}

\begin{proof}
External product, inverse image, and
duality preserve bounded holonomicity.
So does algebraic direct image.
The tensor product is the inverse
image of an external product along
the diagonal. This proves the first
assertion.

For the mapping equivalence, apply
the Verdier pairing, then the
projection formula:
\[
\begin{split}
 &R\Gamma_{\mathrm{dR}}
   (U,\mathbb D_UM\otimes^!p^!V\otimes^!q^!N)\\
 &\qquad\simeq
 R\Gamma_{\mathrm{dR}}
   (Y,V\otimes^!p_*(\mathbb D_UM\otimes^!q^!N)).
\end{split}
\]
Apply the pairing again to the bounded
holonomic kernel. At no step was $V$
required to be holonomic. Naturality
gives the coherent composition
comparisons.
\end{proof}

In particular,
$f_!M=\mathbb D_Zf_*\mathbb D_YM$
is left adjoint to $f^!$ on these
source objects, against arbitrary
targets. Properness is not required
for this assertion.

\subsection{Retaining the collision diagram}
\label{v4-arith:centers:subsec:holonomic-diagram}

Call a coefficient system
\emph{configurationwise holonomic} if
its empty coefficient is perfect and
each nonempty coefficient is bounded
holonomic. The bounds may depend on
the configuration.
Let $\mathcal C_{\mathrm h}$ be the
ordered collision category formed
from $\operatorname{Ind}\mathsf H(\mathscr R_I)$,
with a $D(k)$ empty coefficient.

Bounded holonomic modules are compact
in $\operatorname{DMod}(\mathscr R_I)$:
they are bounded coherent, and the
compactness statement is
Drinfeld--Gaitsgory, Section 5.1.17.
Thus the inclusion
$b_I:\operatorname{Ind}\mathsf H(\mathscr R_I)
\to\operatorname{DMod}(\mathscr R_I)$
is fully faithful, preserves colimits,
and has a right adjoint $r_I$.
For every collision morphism $f$,
\[
 r_Yf^!\simeq f^!r_Z.
\]
Indeed, testing both sides on bounded
holonomic $H$ gives
\[
 R\operatorname{Hom}(H,f^!L)
 \simeq R\operatorname{Hom}(f_!H,L)
 \simeq R\operatorname{Hom}(f_!H,r_ZL).
\]
These tests generate the ind-category.
The equivalences are adjunction mates
of the coherent composition maps.

On ind-holonomic modules, $f^!$
preserves colimits and is also a
right adjoint to the ind-extension
of $f_!$, so it preserves limits.
The coherent limit category
$\mathcal C_{\mathrm h}$ is therefore
presentable with componentwise limits
and colimits. Its inclusion $b$ has
right adjoint $r$, and evaluation
$e_I$ has a left adjoint $e_{I,!}$.

\begin{proposition}\label{v4-arith:centers:prop:holonomic-cohom}
For configurationwise holonomic
$M,N$, there is an object
$\mathsf T(M,N)\in\mathcal C_{\mathrm h}$
such that, for every $V\in\mathcal C$,
\[
 R\operatorname{Hom}_{\mathcal C}
       (b\mathsf T(M,N),V)
 \simeq
 R\operatorname{Hom}_{\mathcal C}
       (M,V\star N).
\]
\end{proposition}

\begin{proof}
For a split $I=A\sqcup B$, let $j$
be its disjoint open immersion and
$p,q$ the projections from that
open to $\mathscr R_A,\mathscr R_B$.
The split summand of the mapping
object on the right is
\[
 R\operatorname{Hom}_{\mathcal D_{\mathscr R_A}}
       (P_{p,q}(j^!M_I,N_B),V_A).
\]
This follows from open adjunction
and the preceding lemma.
The kernel is bounded holonomic,
so these maps depend only on
$r_AV_A$.

Define the finite sum
\[
 T_I=\bigoplus_{I=A\sqcup B}
       e_{A,!}P_{p,q}(j^!M_I,N_B).
\]
Its mapping functor gives the
mapping object at $I$.
The collision maps induce, by
enriched Yoneda, the contravariant
coherent diagram of the $T_I$.
Take its colimit, including the
separate empty term, to obtain
$\mathsf T(M,N)$.

Mapping out of the colimit gives
the required coherent limit over
configurations. The adjunction
$b\dashv r$ extends the equivalence
from ind-holonomic tests to all
$V$. The construction retains the
collision homotopies.
\end{proof}

\subsection{The center of the transposed algebra}
\label{v4-arith:centers:subsec:holonomic-center}

Let $A$ be configurationwise holonomic,
augmented, with $A_\varnothing=k$.
For $s\geq0$ form
$T_s=\mathsf T(A^{\star s},A)$,
where $A^{\star0}=K$.
Finite convolution powers remain
configurationwise holonomic.
The two-sided bar gives a simplicial
object $T_\bullet$, with faces and
degeneracies induced by multiplication
and unit. Put $T_A=b|T_\bullet|$.

\begin{theorem}\label{v4-arith:centers:thm:holonomic-center}
For $E=(BA)^\sharp$, the object
$T_A^\sharp$ represents the internal
center on every pro-perfect probe:
\[
 R\operatorname{Hom}_{\mathcal C^\sharp}
       (W,T_A^\sharp)
 \simeq
 R\operatorname{Hom}_{E\text{-}E}
       (W\star_\sharp E,E).
\]
The equivalence includes evaluation,
the internal $E_2$ operations, and
the ordered collision coherences.
\end{theorem}

\begin{proof}
The augmented two-sided bar contracts
by its unit. Free bimodule adjunction
and the proposition give
\[
\begin{split}
 R\operatorname{Hom}_{A\text{-}A}(A,V\star A)
 &\simeq
 \operatorname*{Tot}_s
 R\operatorname{Hom}(A^{\star s},V\star A)\\
 &\simeq
 \operatorname*{Tot}_s
 R\operatorname{Hom}(bT_s,V)\\
 &\simeq R\operatorname{Hom}(T_A,V).
\end{split}
\]
This keeps the entire realization.
It assumes neither a uniform arity
bound nor bounded holonomicity of
$T_A$ itself.

Set $V=\mathsf SW$. Coefficient-linear
bar comparison and transposition
identify the first line with
$R\operatorname{Hom}_{E\text{-}E}
(W\star_\sharp E,E)$.
The last line becomes
$R\operatorname{Hom}_{\mathcal C^\sharp}
(W,T_A^\sharp)$.
These identifications are natural
for all probes.

The identity of $T_A$ corresponds
to a universal map
$\rho:A\to T_A\star A$.
Its iteration gives the coproduct
of $T_A$, and the identity of $A$
gives the counit. There is a second
composition from
$\rho\star_A^{\mathbb L}\rho$, using
\[
 (T_A\star A)\star_A^{\mathbb L}(T_A\star A)
 \simeq(T_A\star T_A)\star A.
\]

Convolution preserves the bar
realization of the balanced product.
The two compositions share a unit
and satisfy coherent interchange.
They give the $E_2$-coalgebra
operations on $T_A$, which transpose
to the center operations on
$T_A^\sharp$. The same adjunctions
retain collision coherences.
\end{proof}

The universal central coaction and
this internal center now coexist.
They solve the two mapping problems
with opposite variance. Neither
construction identifies the actor
of an independently specified
bulk theory.

\section{Free generators and their support}
\label{v4-arith:centers:sec:free}

A two-term bimodule resolution computes the center
of a free algebra without suppressing the support
of its generators. First work in the word category
$(\mathcal W_X,\otimes_{\mathrm{cut}},K)$, where
the tensor uses consecutive cuts. For an object
$M$, put
\[
 A=T_{\mathrm{cut}}(M)
   =\bigoplus_{n\geq0}M^{\otimes_{\mathrm{cut}}n},
 \qquad M^{\otimes_{\mathrm{cut}}0}=K.
\]
This is an algebraic direct sum. Multiplication
concatenates the factors, with no permutation.

\begin{proposition}\label{v4-arith:centers:prop:free-resolution}
There is a cofiber sequence of bimodules
\[
 A\otimes_{\mathrm{cut}}M\otimes_{\mathrm{cut}}A
 \xrightarrow{\partial}
 A\otimes_{\mathrm{cut}}A
 \xrightarrow{\mu}A,
\]
where
$\partial(a\otimes m\otimes b)=am\otimes b-a\otimes mb$.
Consequently the scalar center is
\[
 Z_{\mathrm{word}}(A)\simeq
 \operatorname{fib}\left(
 R\operatorname{Hom}_{\mathcal W_X}(K,A)
 \longrightarrow
 R\operatorname{Hom}_{\mathcal W_X}(M,A)\right).
\]
On a degree-zero coefficient $c:K\to A$,
the arrow is $m\mapsto mc-cm$.
\end{proposition}

\begin{proof}
Fix tensor length $n$. The middle term has
one copy of the word coefficient for each
cut $c_0,\ldots,c_n$. The first term has
one copy for each marked factor
$e_0,\ldots,e_{n-1}$. The maps are
\[
 \partial e_i=c_{i+1}-c_i,\qquad
 \mu c_i=1.
\]

This is the augmented interval complex.
The section $1\mapsto c_0$ has homotopy
$h(c_i)=e_0+\cdots+e_{i-1}$.
Tensoring with the word coefficient
preserves the contraction, including
its internal differential. Summing
over $n$ proves the cofiber sequence.
Apply the derived free bimodule
adjunction to obtain the fiber formula.
\end{proof}

For a differential module $N$, let $iN$
be its positive diagonal insertion,
with coefficient $\Delta_{n*}N$ on a
word of length $n$. Kashiwara
equivalence identifies all these
coefficients with $N$. The collision
index has the singleton as terminal
object, so
\[
 R\operatorname{Hom}_{\mathcal W_X}(iN,G)
 \simeq R\operatorname{Hom}_{\mathcal D_X}(N,G_1).
\]

For $A=T_{\mathrm{cut}}(iN)$,
$A_\varnothing=k$ and $A_1=N$.
The scalar commutator is zero.
Thus
\begin{equation}\label{v4-arith:centers:eq:diagonal-free}
 Z_{\mathrm{word}}(T_{\mathrm{cut}}(iN))
 \simeq
 k\oplus R\operatorname{Hom}_{\mathcal D_X}(N,N)[-1].
\end{equation}
Any internal self-extension of $N$
remains in this formula, shifted by
one center degree.

\begin{example}[Two diagonal generators]
\label{v4-arith:centers:ex:diagonal-generators}
On $X=\mathbb A^1_k$, take
$N=\omega_X\otimes_k k^2$ with trivial
connection. As a right differential
module,
$\omega_X=D/\partial_zD$.
Its free resolution has differential
$P\mapsto\partial_zP$.
Applying Hom into $\omega_X$ gives
$k[z]\xrightarrow{-\partial_z}k[z]$.
Polynomial integration therefore gives
\[
 H^0Z_{\mathrm{word}}(A)=k,\qquad
 H^1Z_{\mathrm{word}}(A)=M_2(k),
\]
and no further groups.

A matrix acts by the derivation that
applies it in each tensor position.
The bracket is the matrix commutator.
No nonzero matrix is an inner boundary,
because the only degree-zero coefficients
are scalars. Positive cup products
vanish, since $H^2=0$.

On $\mathbb G_m$, the derivative on
$k[z,z^{-1}]$ has the additional
cokernel generated by $z^{-1}$.
The same two generators then give
$H^2=M_2(k)$ as well, from the
de Rham residue $dz/z$.
\end{example}

Placing the two generators at the
empty word gives a different result.
Let $M=K\otimes_k(kx\oplus ky)$.
Then $A=K\otimes_k F$ for the ordinary
free algebra $F=k\langle x,y\rangle$,
and its center complex is
\[
 F\longrightarrow F\oplus F,\qquad
 c\longmapsto(xc-cx,yc-cy).
\]
It has $H^0=k$ and
$H^1=\operatorname{Der}(F)/\operatorname{Inn}(F)$,
with no higher groups.
The degree-one space is infinite
dimensional: the derivations
$x\mapsto y^m$, $y\mapsto0$ are
independent modulo inner derivations.
Abelianization kills commutators
and sends their first components
to the independent polynomials
$y^m$.

The difference comes from support.
A diagonal generator can only take
values compatible with the diagonal
mapping adjunction. A generator
at the empty word can take values
in the entire ordinary free algebra.

\subsection{Perfect differential modules}
\label{v4-arith:centers:subsec:perfect-current}

The same interval resolution applies
to the full collision convolution.
Let $M$ be a perfect right
$\mathcal D_X$-complex on
$X=\mathbb A^1_{\mathbb C}$.
Let $G_M$ be its positive diagonal
family, with one small-diagonal
direct image of $M$ on each ordered
chart. Define
\[
 \mathfrak A_M=T_\star(G_M),\qquad
 \mathfrak C_M=K\oplus G_M[1],\qquad
 \mathfrak E_M=\mathfrak C_M^\sharp.
\]

The reduced coproduct of
$\mathfrak C_M$ is zero.
Its cobar is $\mathfrak A_M$, and
its map to $B\mathfrak A_M$ sends
$g[1]$ to the singleton bar word.
Bar--cobar comparison identifies
these coalgebras. At a fixed
configuration the word filtration
is finite.

Write $[G,N]_\star$ for internal
Hom and
$B_M=R\operatorname{Hom}_{\mathcal D_X}(M,M)$
for the scalar endomorphism complex.
The interval resolution, tested on
every coefficient, gives
\[
 Z^{\mathrm{int}}(\mathfrak A_M)
 \simeq
 \operatorname{fib}\left(
 \mathfrak A_M\xrightarrow{\mathrm{ad}}
       [G_M,\mathfrak A_M]_\star\right).
\]
Here $\mathrm{ad}$ is adjoint to
the graded commutator. Transposing
this fiber gives the coHom
\[
 Q_{\mathfrak A_M}\simeq
 \operatorname{cofib}\left(
 [G_M,\mathfrak A_M]_\star^\sharp
 \longrightarrow\mathfrak A_M^\sharp\right),
\]
with its universal central coaction.

\begin{lemma}\label{v4-arith:centers:lem:current-test}
For every complex $V$ over $\mathbb C$,
\[
 R\operatorname{Hom}_{\mathfrak A_M\text{-}\mathfrak A_M}
       (\mathfrak A_M,(VK)\star\mathfrak A_M)
 \simeq V\oplus(V\otimes B_M)[-1].
\]
\end{lemma}

\begin{proof}
Apply the free bimodule resolution
to the indicated target. The
resolution-degree-zero term is $V$.
The two endpoint actions agree, so
the map to the other term is zero.

Only the length-one part of
$\mathfrak A_M$ contributes to
maps from $G_M$. Every longer
word is extended by $j_*$ from
an open missing the small diagonal.
Open adjunction makes its mapping
complex from $G_M$ zero.
The empty word also receives no
map from this positive object.
These vanishings respect every
collision map.

Kashiwara equivalence and the
contractible index of ordered
surjections identify the remaining
maps with
$R\operatorname{Hom}_{\mathcal D_X}(M,V\otimes M)$.
On the affine line, a perfect $M$
is a retract of a finite complex
of finite free modules. Thus $V$
can be moved through this finite
Hom complex, giving $V\otimes B_M$.
The bar degree shifts it by one.
\end{proof}

\begin{theorem}\label{v4-arith:centers:thm:current-obstruction}
If $\mathfrak E_M$ has an internal
center on all pro-perfect probes,
then $B_M$ is perfect over $\mathbb C$.
In particular, $\mathfrak E_{\mathcal D_X}$
has no such center, although each
configuration coefficient is the
dual of a perfect differential-module
complex.
\end{theorem}

\begin{proof}
Coefficient-linear bar comparison
and transposition turn the lemma
into the tests of a proposed center
on $W(V)=(VK)^\sharp$.
Representability would make
$V\mapsto V\oplus(V\otimes B_M)[-1]$
preserve products. Its second
summand would also preserve them.

The complex $B_M$ is bounded.
If $H^q(B_M)$ contained independent
vectors $b_0,b_1,\ldots$, the family
$(b_i)$ would not lie in the image
of
\[
 \left(\prod_{i\geq0}\mathbb C\right)
       \otimes H^q(B_M)
 \longrightarrow
 \prod_{i\geq0}H^q(B_M).
\]
Every tensor on the left uses one
finite-dimensional span for all
coordinates. Products of vector
spaces are exact, so this is also
the product obstruction on
cohomology. Every $H^q(B_M)$ must
therefore be finite dimensional.

For $M=\mathcal D_X$, the complex
$B_M$ is the Weyl algebra
$\mathbb C\langle z,\partial\rangle/
(\partial z-z\partial-1)$ in degree
zero. The sequence $\partial^i$
violates the common finite
differential-order bound.

The coefficient perfectness is a
different statement. On a smooth
chart the direct image of the
free module along a small diagonal
has a finite Koszul resolution
in its normal coordinates.
The normal coordinate operators
form a regular sequence by Weyl
normal ordering. Finite sums,
shifts, cones, retracts, and the
finite-cell dual preserve
perfectness. This proves the
coefficient assertion.
\end{proof}

For scalar probes alone the condition
is also sufficient. If $B$ is a
bounded complex over $k$, the
functor
$V\mapsto V\oplus(V\otimes B)[-1]$
is corepresented by
$k\oplus B^\vee[1]$ exactly when
$B$ is perfect. Finite duality
proves sufficiency, and the same
product argument proves necessity.
Extending this scalar representation
to all collision probes still
requires compatible geometric
kernels.

The scalar center of
$\mathfrak A_{\mathcal D_X}$ remains
$k\oplus B_M[-1]$. A right-linear
endomorphism $a$ extends to a
derivation $D_a$ of the free
convolution algebra. Its bracket
is
\[
 [D_a,D_b]=D_{ab-ba},\qquad
 [D_\partial,D_z]=D_1\ne0.
\]
Here $D_1$ multiplies a word by
its length. It is not an inner
boundary, since degree-zero
coefficients are scalars.

Coefficient transposition sends
$z$ to $z$ and $\partial$ to
$-\partial$. Hence the transposed
commutator is $-1$, displaying
the reflected bracket sign
on a nonzero class.

These free-algebra examples have
no prescribed field realization
or normalized local vacuum.
To apply the same test to a
field algebra, its actual
relations must enter the
endomorphism and representability
calculation.

\section{Heisenberg states and the word center}
\label{v4-arith:centers:sec:heisenberg}

The Heisenberg residue algebra has the entire
vacuum module at one position. It is therefore
different from a free word algebra on two
diagonal current generators. The vacuum,
products of currents, and all translations
must remain in the coefficient module.

Let $k$ have characteristic zero. Take
$B=k[\kappa]$, or a field over $k$ with
a specified level $\kappa$. Let $g=(g_{ij})$
be a nondegenerate symmetric form on a
rank-two current space. The even currents
satisfy
\[
 [a^i_\lambda a^j]=\kappa g_{ij}\lambda\mathbf1.
\]

Their state module and translation are
\[
 V_\kappa=B[x^i_m:i=1,2,\ m\geq1],\qquad
 x^i_m=a^i_{-m}\mathbf1,\qquad
 T=\sum_{i,m\geq1}m x^i_{m+1}\partial_{x^i_m}.
\]
Every state is a polynomial in finitely
many variables. There is no completion
in state weight. The right differential
module $\mathscr V_\kappa$ has action
\[
 (f\,dz\otimes v)\partial_z
 =-f'\,dz\otimes v-f\,dz\otimes Tv.
\]
All maps are relative to $B$.

Let $C(V_\kappa)$ be the word residue
coalgebra and
$A_\kappa=\Omega_{\mathrm{cut}}C(V_\kappa)$.
The equation $xy-yx=\kappa1_K$ cannot
present this algebra on positive
generators at nonzero level. Applying
the augmentation would give $0=\kappa$.
The Heisenberg bracket instead takes
values in the positive local vacuum
line inside $\mathscr V_\kappa$.
That line is distinct from the
empty-word unit.

\subsection{The full degree-one kernel}
\label{v4-arith:centers:subsec:derivations}

Put $\nabla=\partial_z+T$. For a
pair of polynomials
$P=(P_1,P_2)\in B[z,x^i_m]^2$, define
\begin{equation}\label{v4-arith:centers:eq:heis-kernel}
 \mathcal K_{ij}(P)=
 \sum_{s=1}^2\sum_{m\geq1}\frac1{(m-1)!}
 \left(
 g_{is}\lambda^m\partial_{x^s_m}P_j
 -g_{js}(-\lambda-\nabla)^m\partial_{x^s_m}P_i
 \right).
\end{equation}
The operator $\nabla$ in the second
term acts on its coefficient to the
right. All sums are finite.
A \emph{chiral derivation} is a
differential-module endomorphism
$\delta$ such that
\[
 \delta r_{12}
 =r_{12}(\delta\boxtimes1+1\boxtimes\delta)
\]
on the entire localized two-input
coefficient, for every multiplier.

\begin{theorem}\label{v4-arith:centers:thm:heis-one}
The full scalar word center satisfies
\[
\begin{split}
 H^qZ_{\mathrm{word},B}(A_\kappa)&=0\quad(q<0),\\
 H^0Z_{\mathrm{word},B}(A_\kappa)&=B,\\
 H^1Z_{\mathrm{word},B}(A_\kappa)
 &\cong
 \{P\in B[z,x^i_m]^2:
           \kappa\mathcal K_{ij}(P)=0\text{ for all }i,j\}.
\end{split}
\]
The degree-one class has
$\delta_z a^i=P_i(z)$ and acts by
the sum of $\delta$ over the state
blocks of every partition. There
is no quotient by non-scalar
inner derivations.
\end{theorem}

\begin{proof}
Use the field retract in the
normalized two-sided bar. A map
from a partition with $r$ blocks
to one with $t$ blocks vanishes
unless the source partition
refines the target. Otherwise
some coordinate difference
annihilates the source and is
inverted on the target. Open
adjunction gives zero.

For $q$ consecutive positive
factors, the source shift is
$r-q$, the target shift is
$t-1$, and Hochschild length
adds $q$. Thus a coefficient
Ext class of degree $e$ has
total degree
$e+r-t+1$. This is at least
one. For $q\geq2$, a cut source
can map to an indecomposable
target only after a merger,
so the degree is at least two.

The uniform lower bounds
survive the complete products
and coherent collision limits.
The length-zero term is $B$,
with zero scalar commutator.
This proves the first two
claims.

In degree one only ordinary
unary maps on equal partitions
remain. Their one-block
restriction is one map
$\delta:\mathscr V_\kappa\to
\mathscr V_\kappa$. The binary
equation says precisely that
it is a chiral derivation.
The simple-pole two-vacuum
input also gives
$\delta\mathbf1=0$.

Subtract its blockwise sum.
Induct on the number of
blocks. Every merger of an
indecomposable partition
remains indecomposable, so
commutation with residues
makes every merger of the
difference zero. The
Heisenberg field map is
faithful on polynomial
states, with arbitrary
spectators and Laurent
multipliers. Consequently
the difference is zero.

For completeness, this
faithfulness can be checked
before quantization. In the
associated graded by current
number, the field is
$(e^{uT}a)b$. On any finite
collection of jets, sufficiently
high translated jets give an
invertible coefficient minor:
the binomial coefficient
matrix has nonzero Vandermonde
leading determinant. Hence the
two-field map is injective.

Wick contractions lower current
number by two, so the same
filtered argument proves
injectivity at every level.
Over $k[\kappa]$, torsion-freeness
reduces to the fraction field.
Multiplying by every power
of the separation tests every
field coefficient through its
residue.

It remains to characterize
$\delta$. The differential-module
and field identities are
\[
\begin{split}
 \partial_z\delta_z+[T,\delta_z]&=0,\\
 \delta_zY(v,u)w
 &=Y(\delta_{z+u}v,u)w+Y(v,u)\delta_zw.
\end{split}
\]
The positive current modes are
$a^i_m=m\kappa\sum_sg_{is}\partial_{x^s_m}$,
while $a^i_0=0$.
Applying $\delta$ to the current
bracket gives exactly
$\kappa\mathcal K_{ij}(P)=0$.
Skew symmetry on a moving
coefficient replaces $T$
by $\nabla$, explaining the
second term of
\eqref{v4-arith:centers:eq:heis-kernel}.

Conversely, put
\begin{equation}\label{v4-arith:centers:eq:heis-reconstruct}
\begin{split}
 Q^i_n(P;z)&=\sum_{\ell\geq0}
       \frac{(\partial_z^\ell P_i(z))_{(n+\ell)}}{\ell!},\\
 \delta_P\mathbf1&=0,\\
 \delta_P(x^i_mv)&=x^i_m\delta_P(v)+Q^i_{-m}(P;z)v.
\end{split}
\end{equation}
Over $B[\epsilon]/(\epsilon^2)$,
the kernel equation says that
$Y(a^i,u)+\epsilon Y(P_i(z+u),u)$
satisfies the Heisenberg
commutators. Its nonnegative
modes kill the vacuum.
The induced vacuum module
therefore makes the recursion
independent of the order of
negative modes.
Commuting through normally
ordered differentiated fields
proves the derivation identity
on every state, including
all contractions.

Finally
\[
 [T,Q^i_n]=-nQ^i_{n-1}-\partial_zQ^i_n
\]
shows that
$\partial_z\delta_P+[T,\delta_P]$
commutes with the currents
and kills the vacuum, hence
vanishes.
The blockwise sum commutes
with cuts, residues, and the
first-cut contraction, so it
is a cocycle for all transferred
operations. The degree bounds
leave no additional degree-one
classes or inner identifications.
\end{proof}

Constant linear images
$P_i=\sum_jM_{ij}x^j_1$ give
the subspace
\[
 \kappa(Mg+gM^{\mathsf t})=0.
\]
For the hyperbolic form this
has dimension one at a
nonzero field-valued level,
and dimension four at level
zero. It is only a subspace
of the full kernel.
For every level, $P_i=c_i(z)$
gives
\[
 \delta_Px^i_m=
 \frac{c_i^{(m-1)}(z)}{(m-1)!},
\]
and $P_i=x^i_2$ gives translation.
At level zero every pair
of differential polynomials
is allowed, with
\[
 \delta_P=
 \sum_{i,m\geq1}
 \frac{\nabla^{m-1}P_i}{(m-1)!}
       \partial_{x^i_m}.
\]

At nonzero level, reconstructing
nonlinear current images
requires the quantum terms
in \eqref{v4-arith:centers:eq:heis-reconstruct}.
For one current, the image
$P=x_1x_2$ is allowed and
\[
 \delta_P(x_1^2)=2x_1^2x_2+\kappa x_4.
\]
A polynomial derivation of the
state multiplication alone
would miss the last term.

A central state also need
not extend to a center
endomorphism of the word
cobar. The regular right
module is free on $K$, so
its degree-zero right-module
endomorphisms are
$\operatorname{Hom}(K,A)=B$.
If such an endomorphism
restricts to $c_{(-1)}$
on states, evaluation on
the vacuum forces
$c=b\mathbf1$.
At level zero every state
is vertex-central, but only
these scalar states extend
in this manner.
Applying a scalar separately to each block is a different
assignment. The scalar $2$ would act by $4$ on two blocks,
whereas $2\operatorname{id}$ acts by $2$.

\subsection{The level variation}
\label{v4-arith:centers:subsec:level-variation}

In the strict word cobar, the
state modules and concatenation
are independent of $\kappa$.
Only the residue differential
$D_\kappa$ varies. Define the
$B$-linear family of cochains
$L_\kappa=\partial_\kappa D_\kappa$.
The derivative here defines a
family of maps. A cochain does
not differentiate its scalar
coefficients. Let $N$ count
current factors in every
state block.

Differentiating $D_\kappa^2=0$
and the differential derivation
identity shows that $L_\kappa$
is a degree-two Hochschild
cocycle. A Wick term with
$c$ contractions contains
$\kappa^c$ and lowers current
number by $2c$. Hence
\[
 bN=2\kappa L_\kappa.
\]
At a nonzero field-valued
level the class is zero.

At level zero, its value
on two holomorphic currents
has a nonzero double-pole
vacuum coefficient.
Every undeformed field is
regular there, so the binary
boundary of a state
endomorphism has zero such
coefficient. The degree
bounds in the preceding
proof exclude additional
derived primitives.
Thus the class is nonzero
at zero level. Over
$k[\kappa]$, reduction
modulo $\kappa$ detects
every coefficient not
divisible by $\kappa$.
Its annihilator is exactly
$(\kappa)$.

Varying the symmetric pairing
similarly gives
$(B/(\kappa))^3$ in degree
two for rank two.
Higher poles give further
classes, as follows.

\section{Degree two on words and collision coefficients}
\label{v4-arith:centers:sec:heis-two}

Let $\mathcal V_\pi$ denote the
localized partition coefficient
with its normal differential
operators and residue maps
$r_{ij}:\mathcal V_\pi\to
\mathcal V_{\pi/(i\sim j)}$.
We first give a presentation
of the entire word degree-two
group. It permits one component
for every finite word.

Let $\mathfrak G_B$ consist
of differential-module maps
$\psi_\pi:\mathcal V_\pi\to\mathcal V_\pi$
compatible with monotone
collisions, zero at the
empty word, and additive
at every admissible cut:
\[
 \psi_{\pi'\mid\pi''}
 =\psi_{\pi'}\boxtimes1+
       1\boxtimes\psi_{\pi''}.
\]
Let $\mathfrak R_B$ consist
of maps
$h_{\pi;ij}:\mathcal V_\pi\to
\mathcal V_{\pi/(i\sim j)}$
with the residue cut-coderivation
and collision-base-change
identities, and the following
linearized residue equations:
\begin{align}
 0={}&h_{kl}r_{ij}+r_{kl}h_{ij}
       -h_{ij}r_{kl}-r_{ij}h_{kl},
       \label{v4-arith:centers:eq:two-mergers}\\
 0={}&h_{(ij),l}r_{ij}+r_{(ij),l}h_{ij}
       -h_{i,(jl)}r_{jl}-r_{i,(jl)}h_{jl}
       \notag\\
    &+h_{j,(il)}r_{il}+r_{j,(il)}h_{il}.
       \label{v4-arith:centers:eq:three-mergers}
\end{align}
The first equation is for
disjoint pairs and the
second for $i<j<l$.
Each subscript after a
merger denotes the surviving
root. No skew-symmetry or
fixed-vacuum condition is
added to $h$.
Define
\[
 (\partial\psi)_{\pi;ij}
 =r_{ij}\psi_\pi-
       \psi_{\pi/(i\sim j)}r_{ij}.
\]

\begin{theorem}\label{v4-arith:centers:thm:word-two}
For the full Heisenberg word
residue algebra,
\[
 H^2Z_{\mathrm{word},B}(A_\kappa)
 \cong\mathfrak R_B/\partial\mathfrak G_B.
\]
Every class is represented
on the unchanged partition
modules by replacing
$r_{ij}$ with
$r_{ij}+\epsilon h_{\pi;ij}$.
Two such representatives
give the same class exactly
when they differ by
$\partial\psi$.
\end{theorem}

\begin{proof}
The two displayed equations
are the first-order terms
in the disjoint-merger and
three-block residue identities.
The cut and collision
identities make the new
differential a compatible
coderivation. Thus each
$h$ defines a residue
coalgebra over the dual
numbers and a Hochschild
class of its cobar.
Conjugating coefficients
by $1-\epsilon\psi$ gives
the stated boundary.

To account for derived
classes, identify the bar
of the field retract with
the cofree cut coalgebra.
Its degree-two coderivations
have only two possible
components: ordinary maps
for a single merger and
$\operatorname{Ext}^1$
between equal partitions.
Represent the latter by
flat module extensions
$E_\pi$ over the dual
numbers. The cocycle
equation says precisely
that residues lift between
these extensions and that
the lifted differential
squares to zero.

There are no extra
attaching maps between
pieces separated by two
or more mergers: their
internal degrees would
be negative. Negative
Ext vanishes, and the
opposite support direction
vanishes by open adjunction.
The same argument removes
higher attaching homotopies.
Degree-zero coproducts
allow only equal-partition
cut maps, which remain
isomorphisms over the
dual numbers.

The one-position extension
is trivialized by its
vacuum attachment. Explicitly,
if its connection perturbation
is $S$, the lifted vacuum
identity gives
$(\partial_z+T)F_0=S\mathbf1$.
This is the condition that
kills its class in
$\operatorname{coker}(\partial_z+T)$.
A splitting of the free
complement then gives the
module trivialization.
Collision compatibility
trivializes every one-block
coefficient.

For a larger indecomposable
partition, split off position
one from its block. The
remaining indecomposable
intervals have smaller word
size, so induction trivializes
their extensions. The
lifted residue that merges
the new singleton back,
restricted to its horizontal
vacuum, gives a splitting
of the target extension.
The simple-pole vacuum
residue is surjective on
principal parts, so this
splitting is an isomorphism
modulo $\epsilon$ and on
both steps of the
$\epsilon$-filtration.
It is therefore an
isomorphism.

Choose the trivializations
first on block patterns
with no adjacent positions
in the same block.
Every other pattern
contracts to one of these
by collapsing maximal
constant intervals.
Kashiwara equivalence
transports the choice, and
the terminal contraction
removes additional descent
choices. Extend across
cuts by tensor products.
All extensions are now
trivial, leaving just the
single-merger changes $h$.

Finally, a degree-one
primitive can have only
ordinary equal-partition
components. Its
coderivation and collision
conditions are exactly
those defining $\mathfrak G_B$.
Its boundary is
$\partial\psi$.
The constructions are
finite at each word size,
and their product retains
the full quotient.
\end{proof}

\subsection{All finite central-polynomial variations}
\label{v4-arith:centers:subsec:central-polynomials}

Let $\mathcal Q_B$ be the
space of polynomial
bilinear forms
\[
 Q(\lambda)=\sum_{d\geq0}C_d\lambda^d,\qquad
 Q(\lambda)=-Q(-\lambda)^{\mathsf t}.
\]
Thus $C_d^{\mathsf t}=(-1)^{d+1}C_d$:
even powers have alternating
coefficients and odd powers
have symmetric coefficients.
Only finite polynomials
are allowed.
Over $B[\epsilon]/(\epsilon^2)$,
consider
\[
 [a^i_\lambda a^j]_\epsilon
 =(\kappa g_{ij}\lambda+
       \epsilon Q_{ij}(\lambda))\mathbf1.
\]
The state module, vacuum,
and translation are fixed.
No requirement of conformal
weight preservation is
imposed.

This is a flat field
deformation on every
polynomial state. To
construct it explicitly,
write $(p)_d=p(p-1)\cdots(p-d+1)$
and set
\[
\begin{split}
 J^i_p&=x^i_{-p}\cdot\quad(p<0),\\
 J^i_p&=
 \kappa p\sum_jg_{ij}\partial_{x^j_p}
 +\epsilon\sum_{d=0}^p(p)_d
       \sum_j(C_d)_{ij}\partial_{x^j_{p-d+1}}
       \quad(p\geq0).
\end{split}
\]
The first term is zero
for $p=0$.
The commutator is
\begin{equation}\label{v4-arith:centers:eq:central-modes}
 [J^i_p,J^j_q]
 =
 \kappa p g_{ij}\delta_{p+q,0}
 +\epsilon\sum_{d\geq0}(p)_d(C_d)_{ij}
                      \delta_{p+q,d-1}.
\end{equation}

For the reversed order,
use
$(d-1-q)_d=(-1)^d(q)_d$.
When both indices are
nonnegative and $p+q=d-1$,
the factor $(p)_d$ is
zero. This checks all
mode pairs.

The commutator is central
and has finite pole order.
Moreover,
$[T,J^i_p]=-pJ^i_{p-1}$ and
$J^i(u)\mathbf1=e^{uT}a^i$.
The first identity uses
$(p)_d(p-d)=p(p-1)_d$.

Assign to a polynomial
state the normally ordered
product of its differentiated
current fields. Moving
nonnegative modes through
negative modes gives a
finite Wick expansion.
The commutators are central,
so regrouping three
collections of modes gives
the same expansion.
This proves the field
identities on all states,
with the prescribed vacuum
and translation.

Differentiate all residues
to obtain a word cocycle
$\Gamma_Q^{\mathrm w}$.
For collision coefficients,
put $L=\iota_{\mathrm{ord}}\mathscr V_\kappa$.
Its residue bracket is a
Lie bracket for $\star$,
and the collision cobar is
the envelope $U_\star L$.
Differentiating its bracket
gives the Chevalley cocycle
\[
 \Gamma_Q^{\mathrm c}:
 \Lambda^2_\star L\longrightarrow L
       \longrightarrow U_\star L.
\]
Its differential is the
linearized Jacobi identity.

Write $Q=C_0+Q_+$ with
$Q_+(0)=0$. Set
\[
 M(t)=-\frac{Q_+(-t)}{2t}g^{-1},\qquad
 P(z)=M(T)a+zC_0g^{-1}a,
\]
and define the polynomial
state endomorphism $N_Q$ by
\[
 N_Q\mathbf1=0,\qquad
 N_Qx^i_m=
 \frac{(\partial_z+T)^{m-1}P_i(z)}{(m-1)!}.
\]
Extend by the polynomial
derivation rule and then
by sums over blocks.
Sesquilinearity gives
\[
 \kappa\lambda
 \bigl(M(-\lambda)g+gM(\lambda)^{\mathsf t}\bigr)
 =\kappa Q_+(\lambda).
\]
For the constant alternating
term, the two moving points
give
\[
 \frac{\kappa}{u^2}
 \bigl((z+u)C_0+zC_0^{\mathsf t}\bigr)
 =\frac{\kappa C_0}{u}.
\]
The all-state Wick construction
then proves
\begin{equation}\label{v4-arith:centers:eq:central-boundaries}
 bN_Q=\kappa\Gamma_Q^{\mathrm w},
 \qquad d_{\mathrm{CE}}N_Q=\kappa\Gamma_Q^{\mathrm c}.
\end{equation}

At level zero, evaluating
on holomorphic currents
gives the independent
principal parts
\[
 \sum_d d!(C_d)_{ij}[u^{-d-1}]\mathbf1.
\]
All undeformed fields are
regular, so every incoming
ordinary binary boundary
vanishes on this input.

The degree bounds exclude derived alternatives.
In the collision Chevalley filtration, a total-degree-one
primitive is an ordinary map $L\to U_\star L$.
Support forces its values into $L$. The same word
bound was used above.
Thus the kernel of
$Q\mapsto[\Gamma_Q]$ is
exactly $\kappa\mathcal Q_B$
on either carrier.

For $Q=\lambda^3H$ with
$H$ symmetric, the primitive
is $P=\tfrac12Hg^{-1}T^2a$.
The mode variation is
$p(p-1)(p-2)H$ and its
value at $p=3$ is $6H$,
the coefficient of the
fourth-order pole.
The falling factorial
is required by translation.

\subsection{Products of the full derivation kernel}
\label{v4-arith:centers:subsec:common-sector}

Let $\mathfrak d_B$ be the
entire kernel in
Theorem~\ref{v4-arith:centers:thm:heis-one}.
Let $Z_{\mathrm w}$ be
the word scalar center and
$Z_{\mathrm c}$ the scalar
center of $U_\star L$.
The collision Chevalley resolution identifies its degree-one
group with the same chiral derivations. Support puts a map
$L\to U_\star L$ in the primitive part. Its cocycle equation
is the chiral derivation identity.
The empty term is $B$,
so it introduces no inner
quotient.

\begin{theorem}\label{v4-arith:centers:thm:common-two}
For $\alpha=\mathrm w,\mathrm c$
there are injections
\[
 \Lambda_B^2\mathfrak d_B
 \oplus\mathcal Q_B/\kappa\mathcal Q_B
 \longrightarrow H^2Z_\alpha,
\]
given by derivation cup
products and the classes
$\Gamma_Q^\alpha$.
Their common parameter
module identifies these
submodules canonically.
On the exterior summands,
the identification preserves
the degree-one Gerstenhaber
action.
\end{theorem}

\begin{proof}
External products of
independent differential-module
endomorphisms at two
distinct positions are
independent. Over a field,
finitely many evaluations
on states and output
coordinate functionals
detect any finite
independent family.
Apply them separately in
the two variables.
Over $k[\kappa]$, the
derivation module and its
tensor square are torsion-free,
since this ring is a
principal ideal domain.
The fraction-field argument
then proves the assertion.
Antisymmetrization embeds
the exterior square.

For words, the two-singleton
component of a cup is
$r_{12}(\delta\boxtimes\eta)$.
Its sum with the exchanged
component is
\[
 r_{12}(\delta\boxtimes\eta-\eta\boxtimes\delta).
\]
This detector vanishes on
boundaries with zero unary
component. Faithfulness of
the field map with all
Laurent multipliers, followed
by the external-product
injection, detects every
finite wedge sum.
Central-polynomial variations
have zero detector because
their local brackets remain
skew-symmetric.

For collision coefficients,
the Chevalley cup is
\[
 (\delta\smile\eta)(x,y)
 =\delta x\,\eta y-\delta y\,\eta x.
\]
Its PBW-degree-two symbol
on separated positions is
$\delta\boxtimes\eta-
\eta\boxtimes\delta$.
A degree-two boundary with
zero unary component comes
from an ordinary map
$L\to U_\star L$.
Support forces this map
into $L$, so its boundary
has zero separated
PBW-degree-two symbol.
The cocycle
$\Gamma_Q^{\mathrm c}$
also takes values in $L$.

The symbol therefore
detects the wedge
coefficient of any
relation. The pole
detector and
\eqref{v4-arith:centers:eq:central-boundaries}
then detect its central
coefficient.

The constructions retain
the current space, pairing,
vacuum, coordinate, and
state module and commute
with changes of current
basis. Finally,
\[
 [\theta,\delta\smile\eta]
 =[\theta,\delta]\smile\eta+
       \delta\smile[\theta,\eta]
\]
on both carriers. This
is the natural action
on $\Lambda^2\mathfrak d_B$.
\end{proof}

At nonzero field-valued
level the central-polynomial
summand vanishes. At
level zero it is the
entire $\mathcal Q_k$.
In particular, cubic
central polynomials lie
outside the span of
pairing variations and
all derivation cups.
The theorem does not
classify the remaining
word residue quotient or
extend the comparison
to it. Nor does it claim
that arbitrary nonlinear
derivations preserve the
central-polynomial summand.

\section{The collision center at every configuration}
\label{v4-arith:centers:sec:all-configurations}

The internal collision center admits
a stronger support theorem. Work
over a characteristic-zero field on
$X=\mathbb A^1_k$. Let $V$ be an
ordinary unital vertex algebra,
$\mathscr V$ its right differential
module, and
\[
 L=\iota_{\mathrm{ord}}\mathscr V,\qquad
 U=U_\star L,\qquad
 C=\operatorname{CE}_\star(L),\qquad
 Z=Z^\star(U).
\]
The envelope and its coefficients
lie in the heart. Its center uses
the complete Chevalley cochains,
equivalently the twisted bar Hom.

\begin{theorem}\label{v4-arith:centers:thm:empty-support}
For every such $V$,
$Z_I=0$ at every nonempty
configuration $I$. The empty
coefficient gives an equivalence
of algebras
\[
 Z_\varnothing K\xrightarrow{\sim}Z.
\]
\end{theorem}

\begin{proof}
For right differential modules
$T_1,\ldots,T_n$, use the
positive spectator object
\[
 W_n=\iota_{\mathrm{ord}}T_1
       \star\cdots\star
       \iota_{\mathrm{ord}}T_n.
\]
Distinguish its first factor.
In each partition summand,
let $B$ be that factor's
block and $a$ its first
coordinate.

Insert a new vacuum at
$z_q$, with simple pole
$u^{-1}$, $u=z_q-a$.
Retain the output block
containing $B$. In the
normal direction this uses
the quotient
\[
 k[u,u^{-1}]\longrightarrow
 \Delta_u=k[u,u^{-1}]/k[u].
\]
Writing $e_j=\partial_u^j[u^{-1}]$,
the Euler operator
$u\partial_u+1$ has eigenvalue
$-j$ on $e_j$.
Its degree-zero normal
projection selects $e_0$.

Apply this projection before
taking the remaining tangential
derived Hom. It therefore
defines an extraction operator
$R_p$ for arbitrary spectator
modules, with no compactness
assumption on them.

The Chevalley differential
has three kinds of terms.
Vacuum action on the output
block containing $B$ gives
the normal quotient above,
whose projection is the
identity. Old actions and
brackets commute with the
normal projection and
contribute $-dR_p$.
The bracket of the new vacuum with an old generator has zero
residue. The old separation $t$ from $B$ is invertible, so
$(s+t)^{-1}$ is regular in the new collision parameter $s$.
Thus
\[
 R_{p+1}d+dR_p=1.
\]

A target partition that
splits a source block
receives no map, by the
support argument used
above. A collision across
two spectator blocks
annihilates the source
summand, since it is an
open direct image and
$\delta^!j_*=0$ there.
Chart changes preserve
the internal order of $B$.
These observations make
the extraction compatible
with the full collision
diagram.

They also apply
before the coefficient
resolutions are totalized,
so the identity holds in
derived Hom.
The Chevalley filtration
has $p,q\geq0$ and only
finitely many terms in
each total degree.
Its contracted rows imply
$R\operatorname{Hom}(W_n,Z)=0$.

We now detect the coefficients
inductively. Let $F$ be
positive and vanish at
all configurations of size
less than $n$. Stratify
each ordered collision
scheme by the number of
distinct coordinates.
On a stratum with $r<n$
coordinates, collision
compatibility identifies
the exceptional restriction
with that of $F_{[r]}$,
which is zero. Hence
local cohomology on the
locus of at most $n-1$
coordinates vanishes.

On the $n$-coordinate
stratum, the labeled
partition index has the
collapse to $[n]$ as
terminal object. Its
coherent mapping limit is
therefore the identity
decomposition at $[n]$.
If $O_n$ is the distinct
configuration open, this
gives
\[
 R\operatorname{Hom}(W_n,F)
 \simeq
 R\operatorname{Hom}_{\mathcal D_{O_n}}
 ((T_1\boxtimes\cdots\boxtimes T_n)|_{O_n},
       j_n^!F_{[n]}).
\]
All overlap homotopies
remain in this equivalence.

Apply the formula to the
positive part of $Z$,
assuming the lower
coefficients already
vanish, and take
$T_i=\mathcal D_X$.
The source on $O_n$
is $\mathcal D_{O_n}$.
The open $O_n$ is affine,
so the vanishing of this
derived Hom forces
$j_n^!Z_{[n]}=0$.
The lower-coordinate
local cohomology also
vanishes. The localization
triangle gives $Z_{[n]}=0$.
Induction proves the
claim, and the isolated
empty factor gives the
algebra equivalence.
\end{proof}

The simple pole is necessary
in the contraction.
Replacing it by $u^{-2}$
gives $[u^{-2}]=-e_1$,
whose projection to $e_0$
is zero. The identity
term in the contraction
then disappears.
The theorem also uses
the envelope of an
ordinary unital vertex
algebra in the heart.
It makes no claim for
arbitrary convolution
algebras outside that
domain.

For any $W\in\mathcal C$,
the theorem gives
\[
 R\operatorname{Hom}_{\mathcal C}(W,Z)
 \simeq
 R\operatorname{Hom}_k(W_\varnothing,Z_\varnothing).
\]
Thus the scalar center
contains all the surviving
information, while its
positive evaluation is
zero. If an algebra
$\mathcal B$ has
$\mathcal B_\varnothing=k$,
every algebra map
$Z\to\mathcal B$ factors
through its unit and an
augmentation
$Z_\varnothing\to k$.
For the diagonal actor
cobar introduced below,
the augmentation matching
the convolution evaluation
therefore gives a coherent
scalar comparison.
The normalized homotopy
complex has an empty-supported
source and positive target,
so it is zero. The
corresponding homotopy is
unique up to a contractible
choice.

\subsection{Global modes in degree one}
\label{v4-arith:centers:subsec:global-modes}

The support theorem does
not make the scalar
center small. The
Chevalley filtration gives
\[
 H^iZ_\varnothing=0\ (i<0),\qquad
 H^0Z_\varnothing=k,\qquad
 H^1Z_\varnothing=\operatorname{Der}(L).
\]
Indeed, in total degree
one the term
$\operatorname{Ext}^1(K,U)$
is zero. An ordinary
map $L\to U$ has
primitive values by
support, and its
Chevalley differential
is the derivation
equation. The incoming
scalar commutator is
zero.

In coordinates a derivation
is an endomorphism $h_z$
of $k[z]\otimes V$
commuting with $\partial_z+T$.
Put $\psi=h_0$. Then
$h_z=e^{-z\operatorname{ad}_T}\psi$,
polynomial on each state,
and the residue condition
is equivalent to
\[
 [\psi,Y(a,x)]
 =Y(e^{-x\operatorname{ad}_T}\psi(a),x).
\]
All Laurent multipliers
are required to recover
this field identity
from residues.

For any state $a$ and
$n\geq0$, its global
mode is
\begin{equation}\label{v4-arith:centers:eq:global-mode}
 h_z^{a,n}
 =\sum_{m=0}^n\binom nm z^m a_{(n-m)}
 =\operatorname{Res}_y(y+z)^nY(a,y).
\end{equation}
The translation identity
$[T,a_{(j)}]=-ja_{(j-1)}$
and the mode commutator
formula prove the
criterion. Thus every
global mode is a
degree-one class.

For the rest of this
section take the rank-one
Heisenberg algebra with
$[J_\lambda J]=\kappa\lambda$,
and put $u^{(m)}=m!x_{m+1}$.
Then
$V=k[u,u',u'',\ldots]$,
$Tu^{(m)}=u^{(m+1)}$, and
\[
 J_n=\kappa n!
       \frac{\partial}{\partial u^{(n-1)}}\quad(n\geq1).
\]
At level zero a derivation
is determined freely by
its image of $x_1$, so
$H^1Z_\varnothing\cong k[z]\otimes V_0$.
At nonzero level there
is a complete global-mode
description.

\begin{theorem}\label{v4-arith:centers:thm:global-modes}
Suppose $\kappa\ne0$.
On $V[t]$ put
$\nabla=\partial_t+T$.
The map
\[
 \gamma:V[t]\longrightarrow\operatorname{Der}(L),
 \qquad
 \sum_nt^na_n\longmapsto\sum_nh^{a_n,n},
\]
is surjective with
kernel $\nabla V[t]+kx_1$.
Consequently
\[
 H^1Z_\varnothing
 \cong V[t]/(\nabla V[t]+kx_1).
\]
\end{theorem}

\begin{proof}
For a polynomial $v$ define
its linearization and
variational derivative by
\[
 \ell_v=\sum_{m\geq0}
       \frac{\partial v}{\partial u^{(m)}}\nabla^m,
 \qquad
 \frac{\delta v}{\delta u}
 =\sum_{m\geq0}(-\nabla)^m
       \frac{\partial v}{\partial u^{(m)}}.
\]
The formal adjoint of
$\sum p_m\nabla^m$ is
$\sum(-\nabla)^m\circ p_m$.
Integration by parts
gives
$\varphi P(\chi)-P^*(\varphi)\chi
\in\nabla V[t]$.
Inspection of the highest
jet gives
$\ker\nabla=k$ and
$\ell_{\nabla v}=\nabla\ell_v$.

The variational kernel
is exactly $\nabla V[t]$.
To prove the reverse
inclusion, decompose
$v$ by polynomial current
degree. For a homogeneous
piece $v_r$ with $r>0$,
Euler's identity and
integration by parts
give
$rv_r\equiv u\,\delta v_r/\delta u$
modulo $\nabla V[t]$.
The current-degree-zero
part lies in
$k[t]=\partial_tk[t]$.
Also, if $\ell_f$ is
self-adjoint, then
\[
 f=\frac{\delta}{\delta u}
       \left(\int_0^1u\,f(\lambda)\,d\lambda\right),
\]
where $f(\lambda)$ scales
every current jet by
$\lambda$. Differentiating
$\lambda f(\lambda)$
under the polynomial
integral proves this
identity.

For a derivation $h$,
write $h_zx_1=\sum_jz^jv_j$
and $v=\sum_jt^jv_j$.
The field criterion
gives
$[\psi,J(x)]=\sum_jx^jY(v_j,x)$.
This commutator determines
$h$: it determines its
action on every negative-mode
monomial and $h$ kills
the vacuum.

Since $[J(y),J(w)]$ is
scalar, its commutator
with $\psi$ vanishes.
Expand the resulting
identity in derivatives
of the formal delta
function. The coefficient
identity is
\[
 \sum_{i\geq1}\frac{\lambda^i}{i!}J_iv
 =
 \sum_{i\geq1}
       \frac{(-\lambda-\nabla)^i}{i!}J_iv.
\]
The coefficients are
unique by residue
extraction. Using the
displayed formula for
$J_i$, this says that
$\ell_v\nabla$ is
self-adjoint.

Apply both sides to $1$.
It follows that
$\delta v/\delta u=c\in k$,
so $v=cu+\nabla f$.
Self-adjointness becomes
\[
 2c\nabla
 =\nabla(\ell_f^*-\ell_f)\nabla.
\]
If the middle operator
were nonzero, the
right side would have
order at least two,
because leading coefficients
multiply in the domain
$V[t]$. Hence $c=0$
and $\ell_f=\ell_f^*$.
The polynomial integral
above gives $f=\delta a/\delta u$.

Skew symmetry and the
global-mode formula give
the current image associated
with $\gamma(a)$ as
$\kappa\nabla(\delta a/\delta u)$.
Thus $h=\gamma(\kappa^{-1}a)$.
Finally $\gamma(a)=0$
exactly when
$\delta a/\delta u\in k$.
The variational kernel
identifies these $a$
with $\nabla V[t]+ku$.
\end{proof}

The energy operator is
also visible here.
For the Heisenberg
conformal vector
$\omega=J_{-1}^2\mathbf1/(2\kappa)$,
one has $\omega_{(0)}=T$
and $\omega_{(1)}=E$.
Its first global mode is
$E+zT$. Scalar polynomial
current images give
infinitely many further
classes. Thus degree one
is infinite dimensional
at every level.

\subsection{The part of highest PBW degree}
\label{v4-arith:centers:subsec:top-degree}

The PBW splitting is
$U=\bigoplus_{r\geq0}\operatorname{Sym}^r_\star L$.
Let $Z_\varnothing^{(r)}$
denote cochains whose
values lie in the
$r$th summand.
The internal differential
acts through the inputs,
and the twisting term
is the adjoint action
of $L$ on the outputs.
Both preserve PBW degree.
Therefore
\[
 Z_\varnothing=\prod_{r\geq0}Z_\varnothing^{(r)},
 \qquad Z_\varnothing^{(0)}=k.
\]

Support forces a cochain
with $r$ output blocks
to have at least $r$
input blocks. Hence
\[
 H^iZ_\varnothing^{(r)}=0\quad(i<r).
\]
In degree $r$, the
entire group is the
space of closed ordinary
maps
\[
 \Lambda^r_\star L
 \longrightarrow\operatorname{Sym}^r_\star L.
\]
There is no boundary
to divide by in this
part.

Let
$k[O_r]=k[z_1,\ldots,z_r,
(z_l-z_m)^{-1}]$.
Such maps are equivalent
to $k[O_r]$-linear
endomorphisms $P$ of
$k[O_r]\otimes V^{\otimes r}$
that commute with
$\nabla_l=\partial_{z_l}+T^{(l)}$
and alternate under
simultaneous permutation
of positions and factors.
The Chevalley cocycle
condition is the
derivation identity
in each factor. In
the first factor it is
\begin{equation}\label{v4-arith:centers:eq:slot-derivation}
 \begin{split}
 P(Y(a,x)b\otimes w)
 ={}&Y^{(1)}(a,x)P(b\otimes w)\\
 &+Y^{(1)}(P_{z_1+x}(a\otimes w),x)b .
 \end{split}
\end{equation}
Here the subscript
$z_1+x$ denotes Taylor
expansion, and $Y^{(1)}$
acts on the first
state factor.

The detection lemma
in the support theorem
identifies the maps on
$O_r$. In the Chevalley
differential, choose the
component in which two
specified inputs merge.
Its three terms are
exactly
\eqref{v4-arith:centers:eq:slot-derivation}.
All components arise
this way. This proves
the criterion with all
collision maps retained.

\begin{proposition}\label{v4-arith:centers:prop:all-exterior}
For every $r\geq1$,
cup product gives an
injection
\[
 \Lambda^r_k\operatorname{Der}(L)
 \lhook\joinrel\longrightarrow H^rZ_\varnothing.
\]
Its highest-PBW-degree
component is the
alternating external
product of the
derivations.
\end{proposition}

\begin{proof}
The alternating product
satisfies
\eqref{v4-arith:centers:eq:slot-derivation}
and commutes with the
connections. It is
therefore closed.

Evaluate at distinct
points $c_1,\ldots,c_r$.
A derivation is determined
by its value at one
point, since
$h_{c+x}=e^{-x\operatorname{ad}_T}h_c$.
Tensor products of
independent state
endomorphisms remain
independent, as in
the preceding cup
calculation. Thus
the alternating
external product
is injective.
The highest-PBW-degree
part has no incoming
boundaries.

The cup uses the
unshuffle coproduct.
Under the suspension
convention
$\widetilde P(sy_1\cdots sy_r)
=(-1)^{r(r-1)/2}P(y_1,\ldots,y_r)$,
its symbol is precisely
this alternating
product. This
identifies the cup
class asserted.
\end{proof}

For rank-one Heisenberg,
every positive degree
is consequently infinite
dimensional. At level
zero the highest-PBW-degree
part is all alternating
values
$F\in k[O_r]\otimes V_0^{\otimes r}$
on $x_1^{\otimes r}$,
extended as differential
derivations in each
factor. At nonzero
level the analogous
description has a
variational condition.

\subsection{Mode sequences and rational classes}
\label{v4-arith:centers:subsec:mode-sequences}

A \emph{mode sequence}
is a sequence $(A_m)_{m\geq0}$
of endomorphisms of $V$,
eventually zero on each
state, satisfying
\[
 [T,A_m]=-mA_{m-1},\qquad
 [A_m,Y(b,x)]
 =\sum_{j\geq0}\binom mjx^{m-j}Y(A_jb,x),
\]
where $A_{-1}=0$.
For sequences $A^1,\ldots,A^r$
and $f\in k[O_r]$, define
\begin{equation}\label{v4-arith:centers:eq:multi-mode}
 P^f_{A^1,\ldots,A^r}
 =
 \sum_{m_1,\ldots,m_r\geq0}
 \partial_{z_1}^{(m_1)}\cdots
 \partial_{z_r}^{(m_r)}f\,
 A^1_{m_1}\otimes\cdots\otimes A^r_{m_r}.
\end{equation}
The sum is finite on
each tensor of states.
The first sequence
identity cancels the
two terms in
$[\nabla_l,P]$.
The second, together
with Taylor's formula,
gives
\eqref{v4-arith:centers:eq:slot-derivation}.
Alternating this operator
therefore gives a closed
map of highest PBW
degree.

At nonzero level, the
current sequence and
$f=(z_1-z_2)^{-1}$
give an already
alternating operator
with
\[
 P^f_{J,J}(x_1\otimes x_1)
 =-2\kappa^2(z_1-z_2)^{-3}
       \mathbf1\otimes\mathbf1.
\]
Every finite sum of
derivation cups has
polynomial value on
this input. Thus
this rational class
lies outside their
span.

There is one additional
type of sequence:
\[
 \varphi_m=-\frac{J_{m+1}}{m+1}.
\]
It satisfies the
sequence identities
because $J_0=0$.
For the same $f$,
\[
 P^f_{\varphi,\varphi}(x_1\otimes x_1)
 =\frac{\kappa^2}{z_1-z_2}
       \mathbf1\otimes\mathbf1.
\]
This class cannot be
a sum of operators
using only state
mode sequences.
Their vacuum coefficient
on this input lies in
$\partial_{z_1}\partial_{z_2}k[O_2]$,
by skew symmetry and
$J_0=0$.
Writing $w=z_1-z_2$,
no $\partial_w$-derivative
of a Laurent polynomial
has a $w^{-1}$ term.

\begin{proposition}\label{v4-arith:centers:prop:sequences}
At nonzero rank-one
Heisenberg level, every
mode sequence is
\[
 A_m=a_{(m)}+c\varphi_m
\]
for a state $a$ and
$c\in k$. The scalar
$c$ is unique and
$a$ is unique modulo
the vacuum line.
\end{proposition}

\begin{proof}
Each $A_m$ gives a
global derivation by
the sequence identities.
A translation-commuting
derivation has constant
$h_z$.
The global-mode
classification and
its variational kernel
then express its
value as $a_{(0)}+cJ_1$:
the variational derivative
of a representing
polynomial has the
form $f_0+ct$.

Grade states by energy,
$\deg x_n=n$.
Take a homogeneous
sequence whose $A_m$
has degree $\delta-m$.
If it vanishes before
$m=p$, then $A_p$
commutes with translation
and is $b_{(0)}+cJ_1$.
For $p\geq1$, the states
\[
 b'=\frac{(-T)^pb}{p!},
 \qquad e=\frac{(-T)^{p-1}x_1}{p!}
\]
have mode sequences
vanishing before $p$
and taking values
$b_{(0)},J_1$ at $p$.
Subtract them.
For $p=0$, subtract
the modes of $b$
and the appropriate
multiple of $\varphi$.

Energy forces this
procedure to terminate
once $p\geq\delta+2$.
Only the component
$\delta=-1$ can retain
$\varphi$.

There are only finitely many nonzero homogeneous sequences.
Their values on $x_1$ have finitely many energy components
and eventually vanish in $m$. These values determine all
the endomorphisms.
For uniqueness, energy
$-1$ separates the
$\varphi$ term.
If all nonnegative
modes of $a$ vanish,
skew symmetry on
$x_1$ and the largest
nonzero $J_na$ force
every $J_na$ to vanish.
Then $a$ is scalar.
\end{proof}

\begin{theorem}\label{v4-arith:centers:thm:top-classification}
At nonzero rank-one
Heisenberg level, every
closed map
$\Lambda^r_\star L\to
\operatorname{Sym}^r_\star L$
is a finite sum of
alternations of
\eqref{v4-arith:centers:eq:multi-mode},
using state modes or
$\varphi$ in each
factor.

Equivalently, evaluation
on $x_1^{\otimes r}$
identifies
$H^rZ_\varnothing^{(r)}$
with alternating
$F\in k[O_r]\otimes V^{\otimes r}$
such that
$\ell_F^{(1)}\nabla_1$
is self-adjoint.
The superscript means
linearization in the
first state factor,
with all other factors
as coefficients.
\end{theorem}

\begin{proof}
The slot identity
first makes $P$ vanish
when an input is the
vacuum. Commuting a
negative current mode
through it then reduces
its value to its value
on $x_1$ in that
factor. Induction over
the state energies
shows that $F$ determines
$P$. The formal-delta
commutator calculation
used for global modes,
now over the coefficient
ring of the other
factors, proves
self-adjointness of
$\ell_F^{(l)}\nabla_l$
for every $l$.

The variational argument
works over any differential
domain $(B,\partial)$.
Its kernel statements
become
\[
 \ker(\partial+T)=B^\partial,\qquad
 \ker(\delta/\delta u)
 =(\partial+T)(B\otimes V)+B.
\]
If $\ell_v(\partial+T)$
is self-adjoint, then
$v=(\partial+T)f+b$
with $b\in B$ and
$\ell_f$ self-adjoint.
The proof is the
same integration by
parts and leading-order
calculation as above.
The current-degree-zero
term now remains in
$B$.

Apply this to a homogeneous
current-multidegree
component of $F$.
In every factor with
positive current degree,
$F$ lies in the image
of $\nabla_l$.
These images intersect
as their product:
if
$F=\nabla_1\cdots\nabla_{s-1}g$
and $\delta_sF=0$,
injectivity of the
other $\nabla_l$ on
positive current degree
gives $\delta_sg=0$.
Thus $g$ also lies
in the image of
$\nabla_s$.

Write
$F=(\prod\nabla_l)g$
over those factors.
Their injectivity and
commutation with the
other adjoints imply
that $\ell_g^{(l)}$
is self-adjoint.
Apply the polynomial
variational integral
in each factor to
write
$g=(\prod\delta_l)h$.
It preserves the
self-adjointness conditions
in the other factors.

Expand $h$ as a finite
sum of pure state
tensors with coefficients
in $k[O_r]$.
State-mode sequences
evaluate on the current
as $\kappa\nabla_l\delta_l$.
$\varphi$ evaluates as
$-\kappa$ in factors
of current degree zero.
After dividing these
nonzero scalar factors,
\eqref{v4-arith:centers:eq:multi-mode}
constructs a closed
operator with value
$F$. Alternation and
division by $r!$
give the prescribed
symmetry. Since $F$
determines the operator,
this proves both
descriptions.
\end{proof}

\section{Degree-two rigidity and the degree-three remainder}
\label{v4-arith:centers:sec:low-degree-rigidity}

The highest-PBW-degree calculation does not by
itself determine the lower-PBW-degree groups.
For rank-one Heisenberg coefficients, the
remaining degree-two group vanishes at
nonzero level. In degree three, derived
extension terms can be removed at every
level, leaving two ordinary cohomology
groups to compute.

\subsection{Removing the derived term in degree two}
\label{v4-arith:centers:subsec:derived-two}

At every level the state differential
module splits as
\[
 \mathscr V=\omega_X\oplus\mathscr F,
 \qquad \mathscr F=W\otimes_k\mathcal D_X.
\]
To obtain the splitting, take the
vacuum projection along positive
energy. Translation is injective
on the positive-energy part.
A graded complement $W$ to its
image gives
$V_+=\bigoplus_{n\geq0}T^nW$,
which is the free differential-module
summand.
Consequently
\[
 \operatorname{Ext}^1_{\mathcal C}(L,L)
 \cong
 \operatorname{coker}(\partial_z+T:k[z]\otimes V\to k[z]\otimes V),
 \qquad
 \operatorname{Ext}^{q}_{\mathcal C}(L,L)=0\quad(q\geq2).
\]
The free summand has no higher
Ext, and the vacuum summand
has its two-term de Rham
resolution.

The first Chevalley differential
is injective on this
$\operatorname{Ext}^1$.
Represent a class by
$[v]\circ p$, with $p$
the vacuum projection.
Restrict its differential
to two vacuum inputs.
The three terms become
three signed copies of
$\Delta_*[v]\circ\nu$,
where
\[
 0\to\omega_{X^2}\to
 j_*j^*\omega_{X^2}
 \xrightarrow{\nu}\Delta_*\omega_X\to0.
\]
Their coefficient is odd,
hence nonzero in characteristic
zero.

Composition with $\nu$ is injective on the relevant Ext group.
The regular term has no map to $\Delta_*\mathscr V$.
A normal derivative annihilates its vacuum generator and acts
injectively on the diagonal target.
Kashiwara equivalence
then detects $[v]$.

It follows that
$H^2Z_\varnothing^{(1)}$
is the second cohomology
of the ordinary complex
\[
 B^p=\operatorname{Hom}_{\mathcal C}
       (\Lambda^p_\star L,L),
 \qquad d=d_\beta.
\]
The possible derived
term dies on the second
page, and no later
differential enters
total degree two.

\subsection{From all words to chiral operations}
\label{v4-arith:centers:subsec:word-reduction}

Let $C^p_{\mathrm{ch}}$
be the space of alternating
differential-module operations
\[
 j_*j^*(\mathscr V^{\boxtimes p})
 \longrightarrow\Delta_*\mathscr V,
\]
with the Chevalley differential
of the chiral bracket.
Here $p\geq1$ and every
localized multiplier is
retained. A binary
operation is
\emph{vacuum normalized}
when it vanishes if
either input lies in
the local vacuum line.

An ordinary cocycle in
$B^2$ has a two-block
operation $\rho_n$ for
each reduced word of
length $n$. Its linearized
Jacobi equation is
\eqref{v4-arith:centers:eq:three-mergers},
with the lengths induced
by the corresponding
three-label word.
The vacuum attachment
\[
 j_*j^*(\omega_X\boxtimes\mathscr V)
 \longrightarrow\Delta_*\mathscr V
\]
has regular kernel.
Every map from the
regular kernel to the
diagonal target is zero.
Thus restriction of
$\rho_n$ to either
vacuum axis is uniquely
the attachment followed
by an endomorphism
$\lambda_n$ or $\Psi_n$.

The two families
$(ij)^kl$ and $(ji)^kjl$
have six merger lengths
\[
 (2k,2,2;2,2k,2k+1),\qquad
 (2k+1,2,2;2,2k+1,2k+2).
\]
Insert the vacuum in
each of their three
blocks, using a simple
pole on just one
attachment. The
linearized Jacobi equation
first makes
$\lambda_n=\Psi_n=\ell$
one scalar, since the
only horizontal state
is the vacuum. Subtract
$\ell\beta=d_\beta(\ell\operatorname{id})$.
The remaining operations
vanish on both vacuum
axes. The same two
families then give
$\rho_n=\rho_2$ for
every $n\geq2$.

This common operation
is alternating. With
the vacuum in the
third block, write
the three iterated
principal-part routes
as $L,R,N$, in their
residue orientations.
They satisfy
\[
 L-R+N=0.
\]

The cocycle identity
compares $L(G)-R(G)+N(H)$,
where $G,H$ are the
two exchanged field
series. Subtract the
three-channel identity.
All Laurent multipliers
then give $N(H-G)=0$,
which detects every
series coefficient.
Thus $H=G$, the
chiral skew-symmetry
identity. The word
$123$ gives its
chiral cocycle
equation.

Restriction to the
word $12$ also detects
boundaries. Therefore
\begin{equation}\label{v4-arith:centers:eq:ordinary-two}
 H^2C^\bullet_{\mathrm{ch}}
 \simeq H^2B^\bullet.
\end{equation}
This reduction uses
the horizontal-state
hypothesis. For the
commutative vertex
algebra $k[y]$ with
$T=\partial_y$, the
section $y-z$ is
horizontal, so that
hypothesis fails.

\subsection{Vanishing at nonzero level}
\label{v4-arith:centers:subsec:heis-rigidity}

\begin{theorem}\label{v4-arith:centers:thm:heis-two}
At nonzero rank-one
Heisenberg level,
\[
 H^2C^\bullet_{\mathrm{ch}}=0,\qquad
 H^2Z_\varnothing^{(1)}=0.
\]
Hence
$H^2Z_\varnothing=H^2Z_\varnothing^{(2)}$,
and the complete
degree-two collision
group is the space
of closed maps
described in
Theorem~\ref{v4-arith:centers:thm:top-classification}
for $r=2$.
\end{theorem}

\begin{proof}
Insert a global mode
$\alpha\in V[t]$
as a first input
and take its residue.
This defines
$\iota_\alpha:C^p_{\mathrm{ch}}\to
C^{p-1}_{\mathrm{ch}}$
for $p\geq2$.
Expanding the Chevalley
differential with this
input gives the
Cartan identity
\[
 \mathcal L_\alpha
 =d\iota_\alpha+\iota_\alpha d.
\]
The output action and
the actions on each
old input give
$\mathcal L_\alpha$.
The remaining terms
are the differential
of the insertion.
Residue commutes with
the old operations
by base change.

Give a separation
and a position weight
$-1$, and state
$x_n$ energy $n$.
The mode $t\omega$
acts on a homogeneous
binary cochain of
degree $e$ by
multiplication by $e$.
Every nonzero-degree
cocycle component
is therefore a
boundary, with
primitive
$e^{-1}\iota_{t\omega}A_e$.
The sum of these
primitives is finite
on each input and
coefficient: the
residue uses only
two coefficients
of the field series.

Use the word reduction
to start with a
vacuum-normalized
degree-zero cocycle.
The mode $\omega$
acts as minus the
position derivative.
On endomorphisms of
the free complement
$\mathscr F$, integration
in $z$ solves
$\mathcal L_\omega B=\iota_\omega A$
with $B\mathbf1=0$.
Then $A-dB$ is
translation invariant
and remains vacuum
normalized.

The corresponding
first-order fields
$Y_\epsilon=Y+\epsilon\Phi_A$
satisfy the Borcherds
identities: these
are the coefficients
of the residue
Jacobi identity with
all Laurent multipliers.
Energy and skew
symmetry force the
current bracket to
be
$[J^\epsilon_m,J^\epsilon_n]
=(\kappa+\epsilon\eta)m\delta_{m+n,0}$.
Higher singular terms
have negative state
energy, and the
simple-pole coefficient
is its own negative.

Rescale the current
modes by
$1-\epsilon\eta/(2\kappa)$.
They now satisfy
the original Heisenberg
relations and kill
the vacuum in
nonnegative degree.
The induced Fock
module gives an
isomorphism
$\theta=1+\epsilon\psi$
that fixes the
vacuum and commutes
with translation.

Conjugate $Y_\epsilon$
by $\theta$.
Its current modes
are the original
ones. The Borcherds
identity reconstructs
every negative-mode
state field from
these currents, so
the conjugated fields
are exactly $Y$.
The coefficient of
$\epsilon$ says
$A=d(-\psi)$.

This proves
$H^2C^\bullet_{\mathrm{ch}}=0$.
Equation
\eqref{v4-arith:centers:eq:ordinary-two}
and the removal of
the derived term
give the second
vanishing. The
PBW degree bound
leaves only the
highest-degree part.
\end{proof}

At level zero the
central-polynomial
classes already
show that the
lower-PBW-degree
group need not
vanish.

\subsection{Formal brackets on the fixed coefficient module}
\label{v4-arith:centers:subsec:formal-rigidity}

Let $R=k[[t_1,\ldots,t_d]]$
for finitely many
degree-zero parameters,
with maximal ideal
$\mathfrak m$.
Complete only in these
parameters. The
polynomial state
module and all word
sizes remain as above.
For an endomorphism
$h\in B^1$, fix
the sign convention
\[
 (d_\beta h)(x,y)
 =\beta(hx,y)+\beta(x,hy)-h\beta(x,y).
\]

\begin{theorem}\label{v4-arith:centers:thm:formal-rigidity}
At $\kappa\ne0$,
every formal Lie
bracket
\[
 \beta_R=\beta+
 \sum_{|\alpha|\geq1}\mathbf t^\alpha b_\alpha,
 \qquad b_\alpha\in B^2,
\]
on the fixed
$L[[\mathbf t]]$
has a continuous
automorphism
$\Phi=1\bmod\mathfrak m$
such that
\[
 \Phi\beta_R=\beta(\Phi\star\Phi).
\]
If every correction
is vacuum normalized,
$\Phi$ fixes the
local vacuum line.
\end{theorem}

\begin{proof}
Suppose the bracket
has already been
made equal to
$\beta$ modulo
$\mathfrak m^n$.
Its order-$n$ error
$b_n$ is a
$d_\beta$-cocycle:
two correction terms
have order at
least $2n\geq n+1$.
The space
$\mathfrak m^n/\mathfrak m^{n+1}$
is finite dimensional,
so the vanishing
of $H^2B^\bullet$
gives $h_n$ with
$d_\beta h_n=b_n$.
For $g_n=1+h_n$,
\begin{equation}\label{v4-arith:centers:eq:gauge-step}
 g_n(\beta+b_n)(g_n^{-1}\star g_n^{-1})
 =\beta+b_n-d_\beta h_n
 =\beta\pmod{\mathfrak m^{n+1}}.
\end{equation}
The sign follows
from $g_n^{-1}=1-h_n$
to this order.

The products
$g_N\cdots g_1$
stabilize modulo
every power of
$\mathfrak m$.
Their limit is
$\Phi$, with
inverse given by
the compatible
reverse products.
Each parameter
coefficient uses
finitely many
compositions of
maps in $\mathcal C$,
so it preserves
all collisions.

For a normalized
error, every
primitive annihilates
the vacuum.
Indeed,
$h\mathbf1=c\mathbf1$
by horizontality.
Evaluating $d_\beta h$
on the simple-pole
two-vacuum input
gives $c+c-c=c$.
Normalization forces
$c=0$.
Conjugation by
$1+h_n$ therefore
preserves both
vacuum attachments,
by their naturality.
Induction gives
the last assertion.
\end{proof}

Tensor powers of
$\Phi[1]$ give an
isomorphism of
the completed
Chevalley coalgebras,
and applying $\Phi$
in every cobar
factor gives an
isomorphism of
their cobars.
The unshuffle
coproduct and all
cut maps commute
with these tensor
powers.
At a fixed
configuration, only
finitely many
nonempty factors
occur, so the
maps commute with
parameter completion.
In the normalized
case they are
under the vacuum
coalgebra and its
cobar.

This is rigidity
of strict Lie
brackets on the
fixed differential-module
family. Deformations
of the underlying
module, or arbitrary
higher algebra
operations, are
separate questions.

\subsection{The remaining degree-three groups}
\label{v4-arith:centers:subsec:degree-three}

Keep any level
$\kappa\in k$.
The vacuum splitting
gives
$L=L_0\oplus L_F$,
with
$L_0=\iota_{\mathrm{ord}}\omega_X$
and
$L_F=\iota_{\mathrm{ord}}\mathscr F$.
On the first page
for PBW degree $r$,
write
\[
 E^{a,b,q}_{(r)}
 =
 \operatorname{Ext}^q_{\mathcal C}
 \left(\Lambda^a_\star L_0\star
       \Lambda^b_\star L_F,\operatorname{Sym}^r_\star L\right).
\]
The differential has
parts adding a
vacuum input,
adding a free
input, and replacing
a vacuum input by
two free inputs.

\begin{theorem}\label{v4-arith:centers:thm:degree-three}
Every derived term
of total degree
three vanishes on
the second Chevalley
page. Thus
\[
\begin{split}
 H^3Z_\varnothing
 \cong{}&
 H^3\operatorname{Hom}_{\mathcal C}
       (\Lambda^\bullet_\star L,L)\\
 &\oplus
 H^3\operatorname{Hom}_{\mathcal C}
       (\Lambda^\bullet_\star L,\operatorname{Sym}^2_\star L)\\
 &\oplus H^3Z_\varnothing^{(3)}.
\end{split}
\]
The first two
summands use
ordinary maps and
the Chevalley
differential.
The third consists
of closed maps
$\Lambda^3_\star L\to
\operatorname{Sym}^3_\star L$.
\end{theorem}

\begin{proof}
First, two free
inputs have no
positive Ext into
PBW degree one
or two. For
degree one, the
normal source is
the localization
$\mathcal D_{X^2}[u^{-1}]$,
the colimit of
free modules under
left multiplication
by $u$.
Its Hom into
$\Delta_*\mathscr V$
is the derived
inverse limit of
right multiplication
by $u$. On
$\mathscr V\otimes k[\partial_u]$
this is
$m\partial_u^n\mapsto n\,m\partial_u^{n-1}$,
which is surjective.

The countable tower
has no derived
limit in positive
degrees. Products
over the free
basis are exact.
For PBW degree
two, restriction
to $O_2$ has
free source on
an affine scheme,
giving the same
vanishing.

To treat vacuum
inputs, use the
natural exact
sequence
\[
 0\to\omega_X\boxtimes N\to
 j_*j^*(\omega_X\boxtimes N)
 \xrightarrow{\nu_N}\Delta_*N\to0.
\]
Restriction in
Ext has kernel
the image of
$\operatorname{Ext}^1_X(N,\mathscr V)$
under composition
with $\nu_N$.
For free $N$,
restriction is
therefore injective.

Replace every occurrence
of a label $a$
in a reduced
two-label word
$w$ by two
vacuum labels
$y_0y_1$.
Localize only along
$z_{y_0}=z_{y_1}$.
Of the Chevalley
terms adding a
vacuum, exactly
three survive.
Naturality of the
attachment identifies
each with a
signed copy of
the restricted
class followed
by $\nu$.
Their sum is
an odd multiple.

The same regular-to-diagonal
Hom vanishing
makes attachment
injective on
Ext.
Thus a class
with one vacuum
and one free
input in the
kernel must
vanish.

For two vacuum
inputs and PBW
degree one,
the restricted
class is likewise
zero. It therefore
has the form
$\Delta_*\eta_n\circ\beta_n$
at a reduced
word of length
$n$, with
$\eta_n\in
\operatorname{Ext}^1_X(\omega_X,\mathscr V)$.
Let $K_L,K_R,K_N$
be the three
vacuum merger
routes. They obey
$K_L-K_R+K_N=0$.
The first two
are independent:
on $1/(s(s-t))$
their values
are $s^{-1}t^{-1},0$,
and on $1/(st)$
they are
$0,s^{-1}t^{-1}$.

This independence
persists with
the Ext coefficients.
Separate the two
normal coordinates
from the tangential
one. The normal
localization is
holonomic and has
a bounded resolution
by finite projectives.
The tangential
vacuum has its
Spencer resolution.
The K\"unneth map
therefore embeds
normal degree zero
tensored with
tangential degree
one into total
degree one.

The cocycle
equation forces
the three sums
of the relevant
$\eta_n$ to agree.
The two word
families displayed
above give
$\eta_{2k+1}=\eta_{2k}$
and
$\eta_{2k+2}=\eta_{2k+1}$.
All $\eta_n$
are one class,
and its family
is the boundary
of that class
in arity one.

For PBW degree
two, restriction
from the two-block
open-image target
loses no class.
Apply the same
three surviving
terms to the
component in
which the two
vacua merge.
They force the
restriction to
vanish and hence
the class itself
to vanish.

Now take a
cocycle in
$E_1^{2,1}$.
The two-free-input
part is zero.
In PBW degree
one, subtract
the boundary
just constructed
from its two-vacuum
part. The
one-vacuum part
then vanishes
by injectivity.
In PBW degree
two, both
vacuum parts
are already
detected injectively.
Thus
$E_2^{2,1}=0$.

The terms
$E_1^{1,2}$
and
$E_1^{0,3}$
also vanish.
The only possible
later source
into $E^{3,0}$
is the degree-two
derived term
already removed.
For PBW degree
three only the
highest-degree
ordinary maps
occur.
First-quadrant
connectivity
excludes extra
limit terms.
\end{proof}

The proof uses
the vacuum splitting
with free complement,
not the numerical
value of the
level. At nonzero
level the final
summand has the
explicit description
of
Theorem~\ref{v4-arith:centers:thm:top-classification}.
The two ordinary
degree-three groups
remain to be
computed. Their
reduction to
ordinary cochains
is not a
vanishing theorem.

\section{The ordinary transverse endomorphism monoid}
\label{v4-arith:centers:sec:transverse}

The diagonal action problem uses another
tensor product. Write
\[
 (F\otimes^!G)_I=\Delta_I^!(F_I\boxtimes G_I)
\]
on the ordered collision category.
Its unit is the vacuum coefficient
system $J$, whose empty coefficient
is $k$. The convolution unit $K$
is instead supported only at the
empty configuration.

Both products come from the same
external product at different
relative positions. Convolution
extends from the disjoint open.
Diagonal tensor restricts to the
ordered branch where each input
position is stacked with its copy.
For one position, the two ordered
diagonal branches give the
localization triangle
\[
 \delta_{12,*}(M\otimes^!N)
 \oplus\delta_{21,*}(M\otimes^!N)
 \longrightarrow q^!(M\boxtimes N)
 \longrightarrow j_*j^!q^!(M\boxtimes N).
\]
The connecting maps are their
residues. Thus diagonal tensor
cannot be replaced by disjoint
convolution.

Let $\operatorname{Coalg}$
denote the enhanced coaugmented
convolution coalgebras in the
bar--cobar domain, with fixed
empty coefficient $k$.
The forward interchange
\[
 \zeta:(F\star G)\otimes^!(H\star L)
 \longrightarrow
 (F\otimes^!H)\star(G\otimes^!L)
\]
keeps common ordered decompositions.
It tensors coproducts and makes
$(\operatorname{Coalg},\otimes^!,J)$
symmetric monoidal.

There is also a reverse
binary map $\chi$ including
the common-decomposition
summands. It is compatible
with collision pullback
and all positive-arity
associativity diagrams.
Indeed, a final decomposition
occurs precisely when the
decompositions in every row
agree, independently of
parenthesization. On that
summand the maps are the
same external tensor
identifications. Hence
\[
 \zeta\chi=1,\qquad
 \chi\zeta=\text{projection to common decompositions}.
\]

For a cut open $j$, consider
$j_*M\otimes^!j_*N\to j_*(M\otimes^!N)$.
The source has zero exceptional restriction to the complement
and restricts to $M\otimes^!N$ on the open.
The comparison is therefore an equivalence. This supplies
the comparison on
each summand.
The binary reverse
map alone gives
no inverse to the
full interchange
or to its unit
comparisons.

\begin{definition}\label{v4-arith:centers:def:transverse}
A \emph{transverse monoid}
is an associative unital
algebra object $F$ of
$(\operatorname{Coalg},\otimes^!,J)$.
It has coalgebra maps
\[
 m:F\otimes^!F\to F,\qquad u:J\to F
\]
with coherent associativity
and unit constraints.
A transverse action on
$C$ is a left module
structure
$F\otimes^!C\to C$
in this category.
\end{definition}

The coalgebra condition
on $m$ is
\[
 \Delta_Fm
 =(m\star m)\zeta
       (\Delta_F\otimes^!\Delta_F),
\]
including its coherent
homotopies. At one
position, after converting
right modules to left
modules with the prescribed
density shift, this is
an ordinary horizontal
associative multiplication
over $\mathcal O_X$.
It is not a product
at distinct positions.

\subsection{The closed category and its universal action}
\label{v4-arith:centers:subsec:closed-transverse}

\begin{theorem}\label{v4-arith:centers:thm:transverse-closed}
The category
$\operatorname{Coalg}$
is presentable, and
$-\otimes^!C$ preserves
colimits. It has a
right adjoint
$\underline{\operatorname{Hom}}_{\operatorname{Coalg}}(C,-)$.
The object
\[
 \mathcal E(C)=
 \underline{\operatorname{Hom}}_{\operatorname{Coalg}}(C,C)
\]
is a transverse monoid
with evaluation action.
For every transverse
monoid $F$, evaluation
induces an equivalence
between monoid maps
$F\to\mathcal E(C)$
and transverse actions
of $F$ on $C$.
\end{theorem}

\begin{proof}
Bar--cobar identifies
$\operatorname{Coalg}$
with augmented algebras
having fixed empty
coefficient. These
are nonunital algebras
on the positive ideal
in a presentable
category whose
convolution preserves
colimits separately.
Their category is
presentable.

Colimits of coalgebras
are computed on reduced
parts: in the opposite
category this is the
fact that limits of
algebras are computed
on underlying objects.
The diagonal tensor
has reduced part
$c\otimes^!d$, and
extraordinary inverse
image and external
product preserve
colimits. Thus
$-\otimes^!C$ preserves
colimits and has the
stated right adjoint.

Composition of the
mapping object gives
an associative unital
endomorphism object.
The unit corresponds
to $\operatorname{id}_C$.
Under the adjunction,
the associative and
unital diagrams for
a map into this
object are exactly
the module diagrams
for its adjoint
evaluation. Applying
the same adjunction
to all iterated
products identifies
their coherent
homotopies. This
proves the equivalence
of action spaces.
\end{proof}

At one position,
\begin{equation}\label{v4-arith:centers:eq:actor-singleton}
 \mathcal E(C)_{\{1\}}
 \simeq
 \underline{\operatorname{Hom}}_X
       (C_{\{1\}},C_{\{1\}}).
\end{equation}
To check this, a
bar word with two
positive factors
vanishes at one
position. The
single remaining
term is the
pointwise mapping
object, and the
bar and cobar
shifts cancel.
Diagonal insertion
and Kashiwara
adjunction identify
its evaluation
with internal Hom
on $X$.

\subsection{Its untwisted cochains}
\label{v4-arith:centers:subsec:pointwise-cochains}

For positive coefficient
objects $c,a$, let
$\underline{\operatorname{Hom}}_!(c,a)$
represent maps
$V\otimes^!c\to a$.
Take $c$ to be the reduced part of a coalgebra $C$
and $a$ the augmentation ideal of an algebra $A$.
The mapping object then has a convolution cup product.
Its adjoint is
the composite
\[
\begin{split}
 (V\star W)\otimes^!c
 &\longrightarrow
 (V\star W)\otimes^!(c\star c)\\
 &\xrightarrow{\zeta}
 (V\otimes^!c)\star(W\otimes^!c)
 \longrightarrow a\star a
 \xrightarrow{\mu_A}a,
\end{split}
\]
using the reduced
coproduct and
the two evaluations.

\begin{proposition}\label{v4-arith:centers:prop:untwisted}
There is a natural
algebra equivalence
\[
 \Omega
 \underline{\operatorname{Hom}}_{\operatorname{Coalg}}(C,BA)
 \simeq
 K\oplus\underline{\operatorname{Hom}}_!(c,a).
\]
On strict representatives,
the differential on
the right is the
ordinary mapping-object
differential and
the product is
the cup just
defined.
\end{proposition}

\begin{proof}
For positive $V$,
the bar of the
free convolution
algebra $T_\star V$
is $K\oplus V[1]$
with primitive
positive part.
Tensoring it diagonally with $C$ still gives zero reduced
coproduct. The interchange retains only common decompositions.
The primitive factor has no decomposition into two positive parts.
Its cobar is
therefore
$T_\star(V\otimes^!c)$.

The free adjunctions
identify maps from
$V$ into the
augmentation ideal
on the left with
maps
$V\otimes^!c\to a$.
This gives the
underlying mapping
object for every
test $V$.
The cup is
associative because
the coproduct,
interchange, and
product are
coherent.

Let $\pi$ be
the bar twisting
cochain of this
cup algebra.
Compose it with
evaluation to get
$\tau=\operatorname{ev}(\pi\otimes^!1)$.
The defining cup
identity turns
the Maurer--Cartan
equation of $\pi$
into the twisting
equation of
$\tau$.
Thus evaluation
defines the
coalgebra map
to $BA$.
Its adjoint
is the identity
on the augmentation
mapping object,
so it is the
required algebra
equivalence.
\end{proof}

Taking $A=\Omega C$
gives a concrete
model for
$\Omega\mathcal E(C)$.
Its unit point
$J\to\mathcal E(C)$
corresponds to
the identity
of $C$.
The twisting
cochain is the
element expressing
this point.
It is not added
to the differential
of the pointwise
cochain algebra.
The convolution
Hochschild center
instead has its
twisted differential
$\partial+[\tau,-]$.

\subsection{Constant symmetries and exchange}
\label{v4-arith:centers:subsec:constant-symmetries}

For a coassociative
coalgebra $H$ over
$k$, put
$H^J_\varnothing=k$
and
$H^J_I=J_I\otimes_kH$
for nonempty $I$.
For a decomposition
into two nonempty
parts, its coproduct
is that of $J$
tensored with
$\Delta_H$.
For an empty part,
it is the identity.
This construction
is symmetric
monoidal.
Thus a bialgebra
$H$ gives a
transverse monoid
$H^J$, with
unit induced
by $1_H$.

An action of
$H^J$ is a
family of horizontal
actions
$H\otimes C_I\to C_I$
compatible with
collisions and
\[
 \Delta_{A,B}(h\cdot c)
 =
 \sum(h_{(1)}\otimes h_{(2)})
       \cdot\Delta_{A,B}(c).
\]
The universal
property gives
its map to
$\mathcal E(C)$.
Over the disjoint
open, factorization
of $H^J$ would
make $\Delta_H$
an equivalence.
For an ordinary
bialgebra this
forces $H=k$:
the two counit
maps would be
the same inverse,
so $h\otimes1=1\otimes h$
and $h=\epsilon(h)1$.

If $C$ is
factorizing, its
two ordered
coproduct components
define an exchange
$R_C$. The
coalgebra action
identity gives
\[
 R_C\Delta_H(h)
 =\Delta_H^{\mathrm{op}}(h)R_C.
\]
For a factorizing
transverse actor
$F$, the analogous
identity is
\[
 R_C(\alpha_1\boxtimes\alpha_1)
 =
 (\alpha_1\boxtimes\alpha_1)
       (R_F\otimes^!R_C)
\]
on the disjoint
two-point locus.
Both follow by
commuting the
action with the
two coproduct
components and
inverting one
factorization map.
The first identity
is a constraint
on a given
bialgebra action.
It does not
reconstruct a
Hopf algebra
from the exchange.

\subsection{A dual for diagonal tensor}
\label{v4-arith:centers:subsec:diagonal-dual}

Transport diagonal
tensor through
bar--cobar and
write the resulting
product on algebras
as $\diamond$,
with unit
$\mathbb U=\Omega J$.
Define
\[
 A^\dagger=[A,\mathbb U]_\diamond.
\]
Then
\[
 \operatorname{Map}(Y,A^\dagger)
 \simeq
 \operatorname{Map}(Y\diamond A,\mathbb U)
 \simeq
 \operatorname{Map}_{\operatorname{Coalg}}
       (BY\otimes^!BA,J).
\]
The last expression
is a pairing
of bar coalgebras
with values in
the vacuum
coalgebra.
Its augmentation
ideal is
$\underline{\operatorname{Hom}}_!
(\overline{BA},\overline{\mathbb U})$.
At one position,
\[
 (A^\dagger)_{\{1\}}
 \simeq
 \underline{\operatorname{Hom}}_X
       (a_{\{1\}}[1],\omega_X).
\]

Evaluation and
symmetry give
a natural map
$A\to A^{\dagger\dagger}$,
and
$\mathbb U^\dagger\simeq\mathbb U$.
The bidual map
requires its own
hypotheses to
be an equivalence.
This internal
diagonal dual
retains positive
configurations,
whereas internal
convolution Hom
into $K$ sees
only the empty
coefficient.

\section{Factorization and the pointed action problem}
\label{v4-arith:centers:sec:pointed}

Let $\operatorname{Coalg}_{\mathrm{fact}}$
be the full subcategory whose ordered
coproduct maps are equivalences on
every disjoint open. It is closed
under diagonal tensor and contains
$J$. A \emph{local center} of a
factorizing coalgebra $C$ is a
final transverse monoid acting
on $C$ within this subcategory.
It is final in the category of
actors together with their actions,
not merely among the underlying
coalgebras.

For a pointed factorizing coalgebra
$u:J\to F$, define
\[
 \Phi_C(F,u)=
 \operatorname{Map}_{C/}(F\otimes^!C,C).
\]
The source lies under $C$ by
$u\otimes^!1_C$, and the target
by its identity. Thus a point
is a coalgebra map whose
restriction along the specified
vacuum is identified with the
identity of $C$.
Representability of this
pointed functor supplies a
local center and its monoid
structure. This is the
centralizer construction of
Lurie's \emph{Higher Algebra},
Proposition 5.3.1.8.

If $L_C$ is a local center,
its action gives
$\beta:L_C\to\mathcal E(C)$.
For every factorizing
transverse monoid $F$,
\[
 \operatorname{Map}_{\mathrm{Alg}}(F,L_C)
 \xrightarrow{\sim}
 \operatorname{Map}_{\mathrm{Alg}}(F,\mathcal E(C)).
\]
Indeed, the fiber over any
action map is the space
of maps to the final
action, hence contractible.
The ordinary endomorphism
monoid $\mathcal E(C)$
is itself a local center
exactly when it factorizes.

\subsection{Obstructions to a factorization right adjoint}
\label{v4-arith:centers:subsec:factorization-obstructions}

Factorizing coalgebras are
closed under sifted colimits.
For each ordered decomposition,
the coproduct comparison
commutes with colimits,
and the diagonal of a
sifted indexing category
is cofinal in every
finite power. Thus the
factorization equivalences
survive.
Arbitrary coproducts do
not have this property.

The reduced part of
$J\sqcup J$ is the sum
of two copies of that
of $J$. At two distinct
positions its factorization
map has two summands
in the source and
four in the target.
Its image consists
of the two diagonal
summands. The two
cross terms remain
in the cofiber.
Therefore
$J\sqcup J$ is not
factorizing.

It follows that the
full inclusion
$\operatorname{Coalg}_{\mathrm{fact}}
\subset\operatorname{Coalg}$
has no right adjoint.
If it had one,
the colimit of any
diagram of objects
in its image would
again lie in that
image. To see this,
factor the colimit
cocone through the
counit. The universal
property gives a
section of the
counit, and full
faithfulness makes
the reverse composite
the identity.
The coproduct just
computed contradicts
this consequence.

There is an independent
obstruction already
at the singleton
algebra.

\begin{proposition}\label{v4-arith:centers:prop:matrix-actor}
On a connected nonempty
curve, a factorizing
transverse monoid
cannot have singleton
algebra
$M_n(k)\otimes J_{\{1\}}$
for $n\geq2$.
Consequently, if $B$
is a constant
commutative algebra
with $2\leq\dim_kB<\infty$,
then $\mathcal E(C_B)$
does not factorize.
\end{proposition}

\begin{proof}
For a constant
singleton algebra
$R$, the two-position
localization boundary
is the normal residue
tensored with a
linear map
$\tau:R\otimes R\to R$.
We check that
diagonal tensor
tensors these
coefficient maps,
including derived
terms.

In a normal coordinate
$u$, put
$P=\mathcal O[u^{-1}]$
and
$D=P/\mathcal O$.
After the density
and cohomological
shifts, a coefficient
map $\tau:V\to R$
has complex
$[V\otimes P\to R\otimes D]$
in degrees zero
and one.
Split the linear
map into the
elementary arrows
$k\xrightarrow1k$,
$k\to0$, and
$0\to k$.
Their complexes
are respectively
$\mathcal O,P,Q$,
where
$Q=[\mathcal O\to P]\simeq D[-1]$.

Localization gives
\[
 P\otimes^{\mathbb L}_{\mathcal O}P=P,\qquad
 P\otimes^{\mathbb L}_{\mathcal O}Q=0,\qquad
 Q\otimes^{\mathbb L}_{\mathcal O}Q\simeq Q.
\]
The last identity
follows by tensoring
$Q\to\mathcal O\to P$
with $Q$.
These identities
give the tensor
of arbitrary
coefficient arrows.

Naturality for
the transverse
product and its
unit therefore
makes $\tau$ a
unital algebra
homomorphism
$R\otimes R\to R$.
For $R=M_n(k)$,
its source is
$M_{n^2}(k)$,
which is simple.
A unital map
must be injective,
but $n^4>n^2$.
This is impossible.

For $C_B$, the
singleton of the
ordinary endomorphism
monoid is
$\operatorname{End}_k(B)$,
with composition
and identity.
The first assertion
applies.
\end{proof}

Linearity alone would
not suffice for this
obstruction. When
$n$ is invertible,
the map
$r\otimes s\mapsto
n^{-1}\operatorname{tr}(s)r$
is unital and linear.
It kills
$1\otimes E_{12}$
and
$1\otimes E_{21}$,
but sends their
product to
$n^{-1}I$.
It fails precisely
the required
multiplicativity.

The inclusion of
factorizing transverse
monoids also has
no right adjoint.
Take $B=k\times k$
and the coproduct
of two copies
of $C_B$ as
transverse monoids.
Its singleton
algebra is $B*B$,
and its two-position
open algebra is
$(B\otimes B)*(B\otimes B)$.
These free products
also compute the
derived coproducts:
the augmentations
split, and the
alternating tensor
description of the
free product is
over the field
$k$.

Let $x,y$ name
the two inclusions
and put $e=(1,0)$.
The factorization
map to
$(B*B)\otimes(B*B)$
kills
\[
 [x(e\otimes1),y(1\otimes e)],
\]
because the images
belong to different
tensor factors.

The source element
is nonzero. Represent
its two idempotents
by
\[
 P=\begin{pmatrix}1&0\\0&0\end{pmatrix},
 \qquad
 Q=\begin{pmatrix}1&1\\0&0\end{pmatrix}.
\]
Use a character
on each unused
tensor factor.
The commutator
maps to
$PQ-QP=E_{12}\ne0$.
Thus the coproduct
fails factorization,
and the same
counit argument
excludes a right
adjoint.

\subsection{The positive result in the heart}
\label{v4-arith:centers:subsec:heart}

Let $\mathcal A$
be an
$\mathcal O_X$-flat
chiral algebra.
Its chiral center
$Z_{\mathrm{ch}}(\mathcal A)$
consists of sections
$z$ for which the
chiral bracket
vanishes on every
regular tensor
$z\boxtimes a$.
This definition
tests regular
inputs. It does
not require all
meromorphic
residues to vanish.
Flatness makes
the center a
commutative chiral
subalgebra.
After changing to
left modules,
$Z_{\mathrm{ch}}(\mathcal A)^\ell$
is a commutative
differential algebra
and acts on
$\mathcal A$.

For a factorization
algebra $\mathcal B$,
let $C_{\mathcal B}$
be its shifted
right-module
coalgebra on the
ordered collision
schemes. Its vacuum
is $J\to C_{\mathcal B}$.
On a chart with
$|I|$ positions,
the module lies
in degree $-|I|$.
This is the
factorization heart.
Diagonal tensor
of such objects
lies in degrees
at most $-|I|$,
since in left-module
notation it is
derived tensor
over functions.

Mapping spaces
from those tensor
powers into a
heart object
are discrete.
They are determined
by top cohomology
on each chart.
The maps agree
on the disjoint
open, and their
difference would
have diagonal-supported
image. Factorization
modules have no
such sections,
so the maps
glue uniquely.
Coalgebra maps
are therefore
maps of factorization
structures, and
being under $J$
is exactly
preserving units.

\begin{theorem}\label{v4-arith:centers:thm:heart}
For every factorization
algebra $\mathcal B$,
the space
\[
 \operatorname{Map}_{C_{\mathcal A}/}
       (C_{\mathcal B}\otimes^!C_{\mathcal A},C_{\mathcal A})
\]
is discrete. Its
points are the
chiral algebra
maps from the
chiral algebra
of $\mathcal B$
to
$Z_{\mathrm{ch}}(\mathcal A)$.

For a commutative
unital differential
algebra $N$,
transverse actions
of $C_N$ on
$C_{\mathcal A}$
are exactly
the differential-algebra
maps
$N\to Z_{\mathrm{ch}}(\mathcal A)^\ell$.
The center coalgebra
with its central
action is final
among these actors.
\end{theorem}

\begin{proof}
The factorization
dictionary just
described makes
the mapping
space discrete.
The equivalence
between factorization
and chiral algebras,
and the tensor
universal property,
are Beilinson--Drinfeld,
\emph{Chiral Algebras},
Sections 3.4.9
and 3.4.15.
A map out of
the tensor product
corresponds to
two commuting
chiral maps.
The condition
under $C_{\mathcal A}$
fixes the second
map as the
identity.
The first
therefore takes
values in the
chiral center.

For commutative
$N$, the two
composites expressing
associativity have
the same restrictions
to the three
factors
$N,N,\mathcal A$.
Uniqueness in
the tensor
universal property
makes them equal.
Thus every
central map
is an action,
with singleton
formula
$b\otimes a\mapsto
\varphi(b)\cdot a$.
Commutative chiral
maps are
commutative
differential-algebra
maps, by
Sections 3.3.3
and 3.3.9 of
the same reference.

An equivariant
map of two
such actors
is a map
$g:N\to N'$
with
$\varphi'g=\varphi$.
For
$N'=Z_{\mathrm{ch}}(\mathcal A)^\ell$
and
$\varphi'=\operatorname{id}$,
there is exactly
one, namely
$g=\varphi$.
\end{proof}

In the heart,
a factorization
coalgebra admits
a transverse
monoid structure
exactly when
its factorization
algebra is
commutative.
The structure
is then unique.
Indeed, let
$e$ be a
transverse unit.
The sections
$e\boxtimes\mathbf1$
and
$\mathbf1\boxtimes e$
restrict to
$e$ on the
diagonal.
Their product
away from the
diagonal is
$\mathbf1\boxtimes\mathbf1$.
No diagonal-supported
section can
alter this
equality.

Restriction gives
$m(e,e)=\mathbf1$,
while the unit
law gives
$m(e,e)=e$.
Thus $e=\mathbf1$.
The product is a factorization map restricting to the identity
on each unit axis. This is equivalent to commutativity and
determines the product uniquely (Beilinson--Drinfeld,
Section 3.4.20).

The map from
this center
coalgebra to
$\mathcal E(C_{\mathcal A})$
is injective on
singleton cohomology.
Evaluate its
action on the
vacuum:
$z\cdot\mathbf1=z$.
The composite
is the inclusion
of the chiral
center into
$\mathcal A$.

\subsection{Every actor bounded above by the heart}
\label{v4-arith:centers:subsec:below-heart}

The same central
object represents
more than the
ordinary heart
actors. Say
that $F$ is
\emph{bounded above
by the factorization
heart} if, on
every chart,
$H^j(F_I)=0$
for $j>-|I|$.
No uniform lower
bound is required.

\begin{theorem}\label{v4-arith:centers:thm:below-heart}
For every pointed
factorizing $F$
with this bound,
central action
gives an equivalence
of discrete spaces
\[
 \operatorname{Map}_{J/}
       (F,C_{Z_{\mathrm{ch}}(\mathcal A)})
 \xrightarrow{\sim}\Phi_{C_{\mathcal A}}(F,u).
\]
Its inverse is
restriction to
the vacuum of
$C_{\mathcal A}$.
Consequently the
center coalgebra
is final among
all factorizing
transverse monoids
with this bound
that act on
$C_{\mathcal A}$.
\end{theorem}

\begin{proof}
The degree bound
makes all mapping
spaces into the
heart target
discrete and
reduces maps to
top cohomology.
In left-module
notation, an
action gives
$a_I:M_I\otimes A_I\to A_I$.
Let $e$ be
the image of
the pointed unit
and set
$\zeta(f)=a_1(f\otimes\mathbf1)$.
The pointed
condition says
$a_1(e\otimes a)=a$.

The two-position
localization boundary
of $F$ gives
a chiral operation
$\mu_F$ on
$M_1$.
Tensor the
localization triangle
with $C_{\mathcal A,2}$
and use naturality
of its connecting
map under the
action.
Apply the resulting
identity to
$(f\boxtimes e)\varphi$
and to the
two sections
$\mathbf1\boxtimes a$,
$a\boxtimes\mathbf1$
of $A_2$.
Their diagonal
restrictions agree,
so
\[
 \mu_{\mathcal A}
       ((\zeta(f)\boxtimes a)\varphi)
 =
 \mu_{\mathcal A}
       ((a_1(f\otimes a)\boxtimes\mathbf1)\varphi).
\]

For $\varphi=1$,
the right side
is zero, making
$\zeta(f)$ central.
For
$\varphi=(z_1-z_2)^{-1}$,
the unit residue
gives
$a_1(f\otimes a)=\zeta(f)\cdot a$.
Both sides are
horizontal maps,
so the result
is coordinate
independent.

Restriction to
the vacuum gives
a coalgebra map
$g:F\to C_{\mathcal A}$.
On distinct
positions it
is the external
product of the
central singleton
maps.
Inside the
localized external
state module,
the center
factorization
module is
\[
 Z(\mathcal A)^\ell_{X^I}
 =
 A_I\cap j_{I,*}j_I^*
       (Z(\mathcal A)^{\ell\boxtimes I}).
\]
This follows by
restricting the
first Chevalley--Cousin
differential to
the central
subalgebra
(Beilinson--Drinfeld,
Sections 3.4.11--3.4.12).
Thus $g$
factors through
the center
coalgebra.
Its collision
and coproduct
identities follow
from those in
$C_{\mathcal A}$
by the injective
inclusion.

The reconstructed
central action
agrees with the
original on the
distinct locus.
Their difference
would have
diagonal-supported
image in $A_I$,
which is zero.
Evaluation on
the vacuum
recovers $g$,
giving the
inverse.

For a monoid
$F$, apply the
same equivalence
to $F\otimes^!F$.
Associativity of
the action makes
$g$ preserve
multiplication.
The pointing
makes it preserve
the unit.
All iterated
mapping spaces
remain discrete,
so the coherences
are unique.
This proves
finality.
\end{proof}

If a local center on all factorizing actors exists,
it receives this ordinary center. Their mapping functors
agree on actors with the stated bound.
Above the heart,
the reduction to
top cohomology
fails. Already
$\pi_1\operatorname{Map}(k[-1],k)=k$
although
$H^0(k[-1])=0$.
The following
Heisenberg actions
retain precisely
such derived
information.

\section{Derived actions of the Heisenberg boundary}
\label{v4-arith:centers:sec:derived-actions}

Fix a field $k$ of characteristic zero. The degree-one operator
in the following calculation comes from varying a charge. It is
invisible in the ordinary center, although it acts on a resolution
of the vacuum.

Let $\mathfrak h_\kappa$ have generators $J_n$, $n\in\mathbb Z$,
and a central element $\mathbf c$, with
\[
 [J_m,J_n]=\kappa m\delta_{m+n,0}\mathbf c.
\]
A \emph{smooth module} has $\mathbf c=1$, and each vector is
annihilated by all sufficiently large $J_n$. Write
$\mathsf{Sm}_\kappa$ for their abelian category. These are precisely
the modules for the Heisenberg vertex algebra $V_\kappa$:
normal products reconstruct the fields from the smooth current
action. On the vacuum,
\[
 V_\kappa=k[x_1,x_2,\ldots],\qquad
 J_{-n}=x_n,\quad J_n=\kappa n\partial_{x_n}\quad(n\geq1),
 \qquad J_0=0.
\]
Smoothness does not require local nilpotence of each positive mode.

\begin{proposition}[Continuous vacuum cochains]
\label{v4-arith:centers:prop:smooth-cochains}
For $M\in\mathsf{Sm}_\kappa$, put
\[
 C^\bullet_{\mathrm{sm}}(M)
 =M\otimes\Lambda(\xi_0,\xi_1,\ldots),\qquad
 |\xi_n|=1,\qquad
 d(m\otimes\omega)=\sum_{n\geq0}J_nm\otimes\xi_n\wedge\omega.
\]
The exterior algebra consists of finite sums of monomials.
Then
\[
 \operatorname{Ext}^i_{\mathsf{Sm}_\kappa}(V_\kappa,M)
 \simeq H^i C^\bullet_{\mathrm{sm}}(M).
\]
\end{proposition}

\begin{proof}
The sum defining $d$ is finite on every element. The nonnegative
modes commute, so $d^2=0$. Let $\mathfrak b$ be their span and let
$\mathfrak a_p$ be the span of $J_n$ for $n\geq p$. The displayed
complex is the filtered union of
\[
 C^\bullet(\mathfrak b/\mathfrak a_p,M^{\mathfrak a_p}).
\]
Its degree-zero cohomology is $M^{\mathfrak b}$, which is
$\operatorname{Hom}(V_\kappa,M)$ by the induced presentation
of the vacuum.

It remains to check acyclicity on injectives in the smooth category.
For $S_p=k[J_0,\ldots,J_{p-1}]$, induction from $\mathfrak b$,
with central charge one, is exact and takes $S_p$-modules
to smooth modules. It is left adjoint to
$M\mapsto M^{\mathfrak a_p}$. Thus an injective smooth module
$Y$ has $Y^{\mathfrak a_p}$ injective over $S_p$.

The finite Koszul complex has no positive cohomology on this module.
Filtered unions are exact, so $C^\bullet_{\mathrm{sm}}(Y)$
has no positive cohomology. The exact cochain functors and
dimension shifting now identify their cohomology with Ext.
\end{proof}

\begin{theorem}[Charge and vacuum extensions]
\label{v4-arith:centers:thm:charge-extensions}
At nonzero level,
\[
 \operatorname{Ext}^*_{\mathsf{Sm}_\kappa}(V_\kappa,V_\kappa)
 =\Lambda(\xi_0)
\]
as a graded algebra. Its nonzero degree-one class is the extension
$V_\kappa[\varepsilon]/(\varepsilon^2)$ with $J_0=\varepsilon$.
For the Fock module $F_\mu$ with nonzero charge $J_0=\mu$,
all groups $\operatorname{Ext}^i(V_\kappa,F_\mu)$ vanish.
At zero level the vacuum cochain differential vanishes, giving
\[
 \operatorname{Ext}^*_{\mathsf{Sm}_0}(V_0,V_0)
 \simeq V_0\otimes\Lambda(\xi_0,\xi_1,\ldots)
\]
as graded vector spaces.
\end{theorem}

\begin{proof}
For $\kappa\ne0$, the vacuum differential is
$\sum_{n\geq1}\kappa n\xi_n\partial_{x_n}$.
The degree-minus-one derivation
\[
 \sigma(\xi_n)=\frac{x_n}{\kappa n}\quad(n\geq1),
 \qquad \sigma(\xi_0)=\sigma(x_n)=0
\]
satisfies $d\sigma+\sigma d=N$, where $N$ counts the $x_n$
and the $\xi_n$ with $n\geq1$. Dividing by $N$ contracts
each positive-degree summand. The remaining complex is
$\Lambda(\xi_0)$ with zero differential.

In the extension displayed in the statement, $J_0$ is nonzero.
It would vanish on a split extension of two vacuum modules.
The extension therefore spans the one-dimensional Ext group.
Its square is zero because Ext in degree two vanishes.

For $F_\mu$, the extra differential is $\mu\xi_0\wedge$.
$\mu^{-1}\partial_{\xi_0}$ contracts the complex.
At $\kappa=0$ every nonnegative mode kills the vacuum, proving
the final assertion.
\end{proof}

The degree-one class differentiates the family $F_\mu$ at $\mu=0$.
Requiring $J_0=0$, or requiring it to act semisimply, excludes
this extension. The ambient module category is therefore part of
the answer. The smooth module
\[
 k[x_1,x_1^{-1},x_2,x_3,\ldots]
\]
also shows why smoothness cannot be replaced silently by local
nilpotence: $J_1=\kappa\partial_{x_1}$ is not locally nilpotent
when $\kappa\ne0$.

\subsection{The vacuum with a moving root}

Let $X=\mathbb A^1_k$ with coordinate $y$.
A \emph{module with a moving root} is a left $D_X$-module
with an $\mathcal O_X$-linear smooth mode action satisfying
\[
 [\nabla_{\partial_y},J_n]=-nJ_{n-1}.
\]
This equation makes the current horizontal when its argument is
measured from the moving root. The vacuum is
$\mathcal V=\mathcal O_X\otimes V_\kappa$, with connection
$\partial_y+T$. On its cochains define
\[
 Tx_n=nx_{n+1},\qquad T\xi_n=(n+1)\xi_{n+1}.
\]
This derivation commutes with $d$.

\begin{theorem}[Internal endomorphisms of the moving vacuum]
\label{v4-arith:centers:thm:moving-vacuum}
The differential $D_X$-module
\[
 \mathcal Z^\ell
 =\mathcal O_X\otimes C^\bullet_{\mathrm{sm}}(V_\kappa),
 \qquad \nabla_{\partial_y}=\partial_y+T
\]
represents the derived functor
\[
 R\operatorname{Hom}_{\mathrm{mov}}
 (N\otimes_{\mathcal O_X}\mathcal V,\mathcal V)
 \simeq
 R\operatorname{Hom}_{D_X}(N,\mathcal Z^\ell)
\]
for every left $D_X$-module $N$. For $\kappa\ne0$, translation
is null-homotopic:
\[
 K(\xi_n)=\kappa^{-1}x_{n+1},\qquad K(x_n)=0,
 \qquad dK+Kd=T.
\]
Consequently the cohomology is
$\mathcal O_X\otimes\Lambda(\xi_0)$ with trivial connection.
\end{theorem}

\begin{proof}
The functor $N\mapsto N\otimes\mathcal V$ is exact.
Its right adjoint takes nonnegative-mode invariants, with the
induced connection. Thus the desired representing object is
the derived invariant complex.

To compute it in the moving category, let
$A_\kappa=U(\mathfrak h_\kappa)/(\mathbf c-1)$ and adjoin
$y,D$ with relations
\[
 [D,y]=1,\qquad [y,J_n]=0,\qquad [D,J_n]=-nJ_{n-1}.
\]
The last derivation preserves the Heisenberg relations:
its value on the central commutator relation is
$-\kappa m(m+n-1)\delta_{m+n-1,0}=0$.
The Ore normal form is free as a right $A_\kappa$-module,
with basis $y^aD^r$. Induction is therefore exact.

It also preserves smoothness. Indeed,
\[
 J_nD^r=\sum_{j=0}^r
 \binom rj n(n-1)\cdots(n-j+1)D^{r-j}J_{n-j}.
\]
A vector killed by all modes of index at least $p$ induces
a vector $y^aD^r\otimes m$ killed by all modes of index
at least $p+r$. Hence the forgetful functor preserves injectives.

Proposition~\ref{v4-arith:centers:prop:smooth-cochains}
computes their derived invariants by the displayed cochains.
The computation respects the connection. Applying an injective
resolution proves the adjunction.

Finally, both $dK+Kd$ and $T$ are derivations. On generators,
\[
 (dK+Kd)x_m=mx_{m+1},\qquad
 (dK+Kd)\xi_m=(m+1)\xi_{m+1}.
\]
Thus $T$ acts trivially on cohomology, as asserted.
\end{proof}

This theorem gives an internal Hom in the moving module category.
To obtain a transverse factorizing action, commutation must hold
on a boundary resolution and across all configurations. We now
construct that additional structure.

\subsection{A quartet resolution}

Assume $\kappa\ne0$. Let $\Lambda[c]$ be the commutative
vertex superalgebra on an odd degree-one field and its translates.
Set
\[
 \mathcal Z_\kappa=V_\kappa\otimes\Lambda[c],
 \qquad \xi_n=\frac{T^nc}{n!},\qquad
 d=(:cJ:)_{(0)},\qquad dJ=\kappa Tc,\quad dc=0.
\]
A zero mode is a derivation of all vertex operations.
These formulas therefore define a differential vertex algebra
whose complex is $C^\bullet_{\mathrm{sm}}(V_\kappa)$.
Its bracket $[J_\lambda J]=\kappa\lambda$ is retained.

Let $E$ have even fields $t,p$ of degree zero and odd fields
$\eta,\iota$ of degrees $1,-1$. Its nonzero generating brackets
are
\[
 [p_\lambda t]=1,\qquad [\iota_\lambda\eta]=1,
\]
together with those forced by graded skewsymmetry.
Give this Weyl--Clifford vertex algebra the differential
\[
 Q=(:\eta p:)_{(0)},\qquad
 Qt=\eta,\quad Q\iota=p,\quad Q\eta=Qp=0.
\]
All state spaces here are algebraic polynomial or exterior spaces.
Each state involves finitely many jets.

\begin{theorem}[Commuting currents on a resolution]
\label{v4-arith:centers:thm:quartet}
In $W=V_\kappa\otimes E$, with differential $Q$, the maps
\[
 \begin{array}{c|cc}
 &J&c\\ \hline
 i:V_\kappa\longrightarrow W&J+p&\\
 j:\mathcal Z_\kappa\longrightarrow W&J-\kappa Tt&-\eta
 \end{array}
\]
preserve the differential vertex operations and have commuting
images. The map $i$ is a quasi-isomorphism. If
\[
 W^{\mathfrak b}
 =\{w:(J+p)_nw=0\text{ for all }n\geq0\},
\]
then $j:\mathcal Z_\kappa\to W^{\mathfrak b}$ is also a
quasi-isomorphism. The resulting derived endomorphism comparison
takes $[c]$ to $[\xi_0]$ and preserves the Yoneda product.
\end{theorem}

\begin{proof}
Write $U=J-\kappa Tt$. The decisive brackets are
\[
 [(J+p)_\lambda(J+p)]=[U_\lambda U]=\kappa\lambda,
 \qquad
 [(J+p)_\lambda U]=\kappa\lambda-\kappa\lambda=0.
\]
Here $[p_\lambda Tt]=\lambda$. Also $Q(J+p)=0$ and
$QU=-\kappa T\eta$, proving compatibility with the differentials.

Use state coordinates $x_n,\pi_n,\tau_m,\bar\eta_m,\bar\iota_n$
with $n\geq1$, $m\geq0$. They are respectively the creation
modes of $J,p,t,\eta,\iota$. The last two are odd.
The differential on $W$ is
\[
 Q=\sum_{m\geq0}\bar\eta_m\partial_{\tau_m}
   +\sum_{n\geq1}\pi_n\partial_{\bar\iota_n}.
\]

The operator
$\sum_m\tau_m\partial_{\bar\eta_m}
+\sum_n\bar\iota_n\partial_{\pi_n}$ contracts the nonconstant
quartet monomials after division by their total number of
generators. Thus $H(E)=k$. Since $p=Q\iota$, the map $i$
induces the identity on $H(W)=V_\kappa$.

The nonnegative modes of $J+p$ are
\[
 (J+p)_0=\partial_{\tau_0},\qquad
 (J+p)_n=\kappa n\partial_{x_n}+\partial_{\tau_n}.
\]
In the coordinates $v_n=x_n-\kappa n\tau_n$ they become
the derivatives $\partial_{\tau_m}$. Their common kernel is
\[
 k[v_n,\pi_n:n\geq1]\otimes
 \Lambda(\bar\eta_m:m\geq0,\bar\iota_n:n\geq1).
\]
Its differential has
$Qv_n=-\kappa n\bar\eta_n$ and $Q\bar\iota_n=\pi_n$.

The image of $j$ consists of the $v_n,\bar\eta_m$ factors.
The remaining pairs are contractible. This is a decomposition
of complexes, without an assertion that the two vertex
subalgebras commute.

The continuous invariant complex in the $\tau$ coordinates
is a polynomial de Rham complex. Hence
\[
 \mathcal Z_\kappa\longrightarrow W^{\mathfrak b}
 \longrightarrow C^\bullet_{\mathrm{sm}}(W)
 \longleftarrow C^\bullet_{\mathrm{sm}}(V_\kappa)
\]
is a zigzag of quasi-isomorphisms. To justify the totalizations,
give each coordinate and each cochain generator weight equal
to its mode index. At fixed weight $r$, the quartet degree
lies between $-r$ and $r+1$, and the cochain degree lies
between $0$ and $r+1$. Both filtrations are finite.

Finally, in the total complex,
\[
 (Q+d_{\mathrm{CE}})\tau_0=\bar\eta_0+\xi_0.
\]
The class of $j(c)=-\bar\eta_0$ is therefore $[\xi_0]$.
Both cohomology algebras are exterior on this class, with
the same unit, which proves the product assertion.
\end{proof}

The correction to the current is necessary.
On the original cochain complex,
$[d,J_{-1}]\mathbf1=\kappa\xi_1$.
The field $p$ makes the boundary current closed, and
$-\kappa Tt$ cancels its bracket with the cochain current.

There is also a binary product in the homotopy category.
In $\mathcal Z_\kappa\otimes E$, with differential $d+Q$,
the two commuting maps are
\[
 \begin{array}{c|cc}
 &J&c\\ \hline
 b_1&J+p&c\\
 b_2&J-\kappa Tt&c-\eta.
 \end{array}
\]
Their product, followed by the inverse of the quartet
quasi-isomorphism, is unital in the homotopy category.
The two unit homotopies are explicit over $k[s,ds]$:
\[
 \begin{array}{c|cc}
 &J&c\\ \hline
 b_1(s)&J+sp+ds\,\iota&c\\
 b_2(s)&J-s\kappa Tt&c-s\eta-ds\,t.
 \end{array}
\]
The signs follow from the tensor differential:
$d(ds\,\iota)=-ds\,p$ and $d(ds\,t)=-ds\,\eta$.
On cohomology this is the exterior product.

A binary operation and its unit homotopies do not yet give
an associative object with all coherences. The construction
below supplies them and specifies its order: its multiplication
is the opposite of the binary operation just displayed.

\subsection{A factorizing ghost action}

For a differential vertex algebra $B$, write
$C_B=C^{\mathrm{CE}}_\star(L_B)$ for its Chevalley collision
coalgebra. Let $R=\Lambda[c]$, $F=C_R$, and $C=C_{V_\kappa}$.
The multiplication of the commutative differential algebra $R$
makes $F$ a transverse monoid. Its vacuum is the pointing
$J\to F$.

We first justify the tensor comparison used in the action.
For a finite configuration $I$, put
\[
 P_I(t)=\prod_{i\in I}(t-z_i),\qquad
 H_I=\varinjlim_m
 \bigl(\mathcal O_{X^I}[t]/P_I(t)^m\bigr)^\vee.
\]
The factorization lattice of $R$ is
$\Lambda_{\mathcal O_{X^I}}(H_I)$, with generators in degree one.
The finite jet modules are free because $P_I^m$ is monic.
Exterior powers and filtered colimits give flat graded
components. With zero differential this is a $K$-flat complex.

Chinese remainders give factorization on disjoint opens.
On collision strata, the positive multiples of the changed
root multiplicities are cofinal, giving the collision maps.
The finite Cousin resolution identifies the lattice with
$C_R$ in the normalization used here.
The same construction gives, naturally,
\[
 C_{R^{\otimes a}}\otimes^!C_B
 \simeq C_{R^{\otimes a}\otimes B}.
\]
It respects multiplication, units, and polynomial homotopies.

The finite Cousin direction has at most $|I|$ terms,
so its passage to the internal differential is valid.
External tensor over $k$, localization, and finite
symmetric-group coinvariants preserve quasi-isomorphisms.
In particular, the maps $C_{i_b}$ below are equivalences.

Let $W_b=V_\kappa\otimes E^{\otimes b}$ and
$i_b(J)=K_b=J+\sum_{j=1}^bp_j$. On the cube of polynomial forms
\[
 \Omega_b=k[s_1,\ldots,s_{b-1}]
 \otimes\Lambda(ds_1,\ldots,ds_{b-1})
\]
define rows recursively by
\[
 M_1=e_1,\qquad
 M_i=s_{i-1}e_i+(1-s_{i-1})M_{i-1}.
\]
Every row has sum one. Put forms to the left of fields and set
\[
 \gamma_i=-\sum_j\bigl(M_{ij}\eta_j+(dM_{ij})t_j\bigr).
\]

\begin{proposition}[The coherent ghost action]
\label{v4-arith:centers:prop:ghost-action}
The assignments $c_i\mapsto\gamma_i$ and $J\mapsto K_b$
give maps
\[
 R^{\otimes b}\otimes V_\kappa\longrightarrow W_b\otimes\Omega_b
\]
whose faces define a unital transverse action of $F$ on $C$.
Its action map is the roof
\[
 C_R\otimes^!C_{V_\kappa}
 \longrightarrow C_{W_1}\xleftarrow{\sim}C_{V_\kappa},
 \qquad c\mapsto-\eta_1,\quad J\mapsto J+p_1.
\]
\end{proposition}

\begin{proof}
For the total differential $D=Q+d$,
$D(M\eta+dM\,t)=0$: the two $dM\,\eta$ terms cancel.
The $\gamma_i$ commute, and
\[
 [K_b{}_\lambda\gamma_i]=-\sum_jdM_{ij}=0.
\]
Also $DK_b=0$ and $[K_b{}_\lambda K_b]=\kappa\lambda$.
Thus the assignments preserve all generating operations.

At $s_i=0$ the adjacent rows coincide. This face multiplies
the corresponding ghost factors. The target comparison replaces
one quartet by
\[
 p\mapsto p_i+p_{i+1},\quad
 \iota\mapsto\iota_i+\iota_{i+1},\quad
 t\mapsto t_i,\quad\eta\mapsto\eta_i.
\]
It preserves brackets and differentials, lies under the
boundary, and is a quasi-isomorphism.

At $s_i=1$ the recurrence restarts at $e_{i+1}$.
This face composes the actions of the two consecutive blocks.
In a common resolution their boundary current is $K_b$.

These face identities agree on intersections.
Repeated merges add $p,\iota$ and retain the first $t,\eta$,
independently of parentheses. Repeated splits use the same
separate quartets, and disjoint merges and splits commute.
Thus the cubes give all associative coherences.

More explicitly, the cubical resolution of a nonunital monoid
has terms $[0,1]^{b-1}\odot F^{\otimes^!b}$.
Zero faces multiply adjacent inputs. Multiplication concatenates
strings with intervening parameter one. Shrinking every parameter
to zero contracts its augmentation to $F$.
Our zero and one faces give a multiplicative map from this
resolution.

Polynomial forms define actual mapping-space cubes as follows.
Pull them back along every simplex of a cube.
Each polynomial simplex algebra has cohomology $k$ by its
Euler contraction. The resulting diagrams
$C_{W_b\otimes\Omega_n}$ receive objectwise equivalences
from the constant boundary diagram. Invert this natural
transformation in the functor category. Naturality in faces,
products, and target comparisons preserves all the identities
already proved.

When every actor input is the vacuum, the roof is the identity
under $C$, with its coherences. The quasi-unit theorem
in \emph{Higher Algebra}, Theorem 5.4.3.5, supplies the
unital refinement with this pointed restriction.
\end{proof}

\begin{proposition}[An action outside the ordinary center]
\label{v4-arith:centers:prop:ghost-detection}
The ghost action does not factor through the diagonal tensor
unit $J$. In particular, the ordinary center of $V_\kappa$
does not represent the pointed action functor on all derived
factorizing actors.
\end{proposition}

\begin{proof}
Put $\mathfrak g=L_{V_\kappa}$ and
$\mathfrak h=L_{R\otimes V_\kappa}$. The one-ghost summand
is $N[-1]$, where $N$ has left state module
\[
 \bigoplus_{m\geq0}ke_m\otimes V_\kappa,\qquad
 T(e_m\otimes v)=(m+1)e_{m+1}\otimes v+e_m\otimes Tv.
\]
For a pointed map, bar--cobar gives a map
$U_\star\mathfrak h\to U_\star\mathfrak g$
under $U_\star\mathfrak g$.
Restrict to $\mathfrak h$, project by PBW symmetrization
to $\mathfrak g$, and restrict to $N[-1]$.
The adjoint action preserves symmetric degree, so this
defines a homotopy invariant
\[
 \chi:\pi_0\Phi_C(F,u)\longrightarrow
 \operatorname{Ext}^1_{\mathfrak g\text{-}\mathrm{Mod}}
 (N,\mathfrak g).
\]
Negating the ghost negates $\chi$.

For our roof, the extension is represented by
\[
 \nu(e_m\otimes v)=-\frac1{m!}:(T^m\eta)i(v):.
\]
A middle module is
\[
 \{(n,w)\in N\oplus W^0:Qw=\nu(n)\}
/\{(0,Qv):v\in W^{-1}\}.
\]
With kernel inclusion $v\mapsto[(0,-i(v))]$, it represents
the indicated connecting class for the usual cochain cone
$d(b,a)=(d_Bb+f(a),-d_Aa)$.
Changing this kernel marking changes the sign, without
changing whether it splits.

Lift $e_0\otimes\mathbf1$ by $-t$. Then
\[
 Q(-t)=-\eta,\qquad (J+p)_0(-t)=-\mathbf1.
\]
The same zero mode vanishes on the quotient and on the kernel.
Changing the lift cannot remove the displayed vacuum class.
The vacuum is not a $Q$-boundary. Hence the extension is
nonsplit.

Main-diagonal differential modules inherit their
usual heart, and convolution with $\mathfrak g$ is exact
there. This heart extension therefore gives a nonzero
derived Ext class.

The automorphism negating the boundary current conjugates
the action to the action with $c\mapsto-c$.
It is realized on the resolution by negating
$J,t,p,\eta,\iota$. Thus conjugation negates the nonzero
invariant $\chi$. An action factoring through $J$ is
unchanged by boundary conjugation. In characteristic zero
these conclusions are incompatible.

Finally, the positive modes on $V_\kappa$ are the nonzero
multiples of $\partial_{x_n}$, so its ordinary center is
$k\mathbf1$ and its center coalgebra is $J$.
This proves the asserted failure of representability.
\end{proof}

\subsection{Tensor products and parameter algebras}

The preceding tensor comparison extends to free fields over $k$.
Suppose $B$ has finitely many homogeneous generators, a PBW basis
of their jets, central polynomial generating brackets, and
a differential linear in the jets. Its factorization lattice
has a PBW filtration whose graded pieces are symmetric powers
of the free principal-parts modules. Each fixed PBW degree
is a bounded complex of free modules. Extensions and filtered
unions preserve $K$-flatness, so every lattice $B_I$ is $K$-flat.

For two such algebras, the generating Lie algebras have a common
constant and otherwise commute. Their vacuum modules give
\[
 (A\otimes_k B)_I=A_I\otimes_{\mathcal O_{X^I}}B_I.
\]
The identities respect collisions, disjoint factorization,
units, and symmetry. Realization contracts one diagonal tensor
unit, giving a symmetric monoidal functor
$B\mapsto C_B$. Polynomial homotopies of the form
$B\otimes_k\Omega_b$ are also preserved.

The coefficient field in this assertion matters. A family over
a differential parameter algebra need not be flat for derived
tensor, even when its underlying graded algebra is free.

\begin{proposition}[A parameter tensor with extra cohomology]
\label{v4-arith:centers:prop:parameter-tensor}
Let $\Omega=k[s,ds]$, with $d s=ds$, and define
\[
 A_+=\Omega[a_n:n\geq0],\qquad
 A_-=\Omega[b_n:n\geq0],
 \qquad Da_n=ds\,a_n,\quad Db_n=-ds\,b_n.
\]
All generating vertex brackets vanish, and translation shifts
each subscript by one and kills $\Omega$.
The augmentations $A_\pm\to\Omega$ are quasi-isomorphisms.
Nevertheless,
\[
 H^0(A_+\otimes_\Omega A_-)=B_0,\qquad
 H^{i\ne0}(A_+\otimes_\Omega A_-)=0,
\]
where $B_0$ is the charge-zero part of $k[a_i,b_j]$ for
$\operatorname{ch}(a_i)=1$, $\operatorname{ch}(b_j)=-1$.
It has the presentation
\[
 B_0=
 \frac{k[z_{ij}:i,j\geq0]}
 {(z_{ij}z_{kl}-z_{il}z_{kj})},\qquad z_{ij}=a_ib_j,
 \qquad Tz_{ij}=z_{i+1,j}+z_{i,j+1}.
\]

For $r$ row indices and $s$ column indices, its polynomial
degree-$q$ dimension is
\[
 \binom{r+q-1}{q}\binom{s+q-1}{q}.
\]
The derived-to-ordinary comparison on cohomology is $k\to B_0$.
Neither $A_+$ nor $A_-$ is $K$-flat over $\Omega$.
\end{proposition}

\begin{proof}
On a monomial of charge $c$, the coefficient complex is
\[
 k[s]\xrightarrow{\partial_s+c}k[s]\,ds.
\]
For $c\ne0$ its inverse on a polynomial $f$ is the finite sum
\[
 (\partial_s+c)^{-1}f
 =\sum_{n\geq0}\frac{(-1)^n}{c^{n+1}}\partial_s^nf.
\]
For $c=0$, differentiation is surjective with kernel $k$.
In each $A_\pm$, only the vacuum has charge zero.
In the tensor, every balanced monomial does.
This proves the cohomology calculation and, more precisely,
the quasi-isomorphism
$\Omega\otimes_kB_0\to A_+\otimes_\Omega A_-$.

Every balanced monomial is a product of $a_ib_j$.
The minor relations exchange the paired column indices of two
factors. They put each monomial into the pairing of its sorted
row and column multisets. Distinct pairs of multisets have
distinct images in $k[a_i,b_j]$, proving that these relations
generate the kernel. This argument applies to infinitely many
variables because each polynomial uses finitely many.
Counting the two multisets gives the dimension formula,
and the Leibniz rule gives translation.

The inclusions $\Omega\to A_\pm$ compute their derived tensor,
whose cohomology is $k$. The comparison is the inclusion of
constants into $B_0$. Tensoring either augmentation with the
other algebra kills the nonzero class $a_0b_0$.
It therefore fails to preserve that quasi-isomorphism,
proving both failures of $K$-flatness.
\end{proof}

\subsection{The coherent cochain monoid}

In the boundary resolution $W_b$, define
\[
 \mathcal Z^{(b)}=W_b^{\mathfrak b},
 \qquad K_b=J+\sum_{j=1}^bp_j.
\]
For an order-preserving surjection
$\phi:\{1,\ldots,b'\}\to\{1,\ldots,b\}$, let
$\phi_*(j)=\min\phi^{-1}(j)$. The merge map fixes $J$ and sends
\[
 p_j\mapsto\sum_{\phi(l)=j}p_l,\quad
 \iota_j\mapsto\sum_{\phi(l)=j}\iota_l,\quad
 t_j\mapsto t_{\phi_*(j)},\quad
 \eta_j\mapsto\eta_{\phi_*(j)}.
\]
It maps $W_b$ to $W_{b'}$ and preserves the boundary current.
The shift $r_a:W_b\to W_{a+b}$ sends $J$ to $K_a$ and
quartet $j$ to quartet $a+j$. Define
\[
 \mu_{a,b}:\mathcal Z^{(a)}\otimes W_b\longrightarrow W_{a+b},
 \qquad z\otimes w\longmapsto :z\,r_a(w):.
\]

\begin{lemma}[Products through boundary resolutions]
\label{v4-arith:centers:lem:tower}
The merge maps are quasi-isomorphisms on both $W_b$ and
$\mathcal Z^{(b)}$. The latter is freely generated by
\[
 J-\kappa Tt_1,\quad t_j-t_1\ (j\geq2),\quad
 p_j,\eta_j,\iota_j.
\]
The products $\mu_{a,b}$ preserve differentials and vertex
operations, are strictly associative, and commute with merges.
Their restrictions map
$\mathcal Z^{(a)}\otimes\mathcal Z^{(b)}$
to $\mathcal Z^{(a+b)}$.
The $b$ cochain inputs and one boundary input have images
\[
 J^{(i)}\mapsto U_i=J+\sum_{l<i}p_l-\kappa Tt_i,\qquad
 c^{(i)}\mapsto-\eta_i,\qquad J\mapsto K_b.
\]
The currents $U_1,\ldots,U_b,K_b$ commute pairwise,
and each has level $\kappa$.
\end{lemma}

\begin{proof}
A merged sum of $p_l$ pairs with the retained $t_{\phi_*(j)}$
through exactly one term. The same holds for $\iota,\eta$.
This checks brackets. The formulas also commute with $Q$.

Every quartet index occurs once, so $K_b$ maps to $K_{b'}$.
The least-element rule is compatible with composition.
Since all $i_b$ are quasi-isomorphisms, so are the merges.

The nonnegative boundary modes are
$\sum_j\partial_{\tau_0^{(j)}}$ and
$\kappa n\partial_{x_n}+\sum_j\partial_{\tau_n^{(j)}}$.
Use coordinates
$x_n-\kappa n\tau_n^{(1)}$ and
$\tau_n^{(j)}-\tau_n^{(1)}$ for $j\geq2$.
The modes become derivatives in $\tau_n^{(1)}$,
so their common kernel has the stated free generators.
Continuous polynomial de Rham cochains in these variables
compare each commutant to the derived invariants of $W_b$.
The finite weight filtrations used in
Theorem~\ref{v4-arith:centers:thm:quartet}
apply for every finite $b$. Naturality and the equivalences
$i_b$ then show that merges are equivalences on commutants.

The image of $r_a$ commutes with $\mathcal Z^{(a)}$:
its current is $K_a$, and its other fields use later quartets.
Thus normal multiplication defines a vertex algebra map.
The identity $r_ar_b=r_{a+b}$ proves associativity.
Merges preserve these shifts and inclusions, giving naturality.

Finally, shift the two images of $j$ successively.
This gives the displayed $U_i$ and ghosts.
For $i<j$, the $\kappa\lambda$ contribution from their
common $J$ cancels the pairing of $-\kappa Tt_i$ with $p_i$.
The same cancellation proves $[U_i{}_\lambda K_b]=0$.
Their self-brackets retain $\kappa\lambda$.
\end{proof}

\begin{theorem}[A factorizing cochain monoid and its action]
\label{v4-arith:centers:thm:cochain-monoid}
The coalgebra $C_{\mathcal Z_\kappa}$ has a unital transverse
monoid structure and a transverse action on $C_{V_\kappa}$.
The $b$-fold product is represented by
\[
 \mathcal Z_\kappa^{\otimes b}\longrightarrow\mathcal Z^{(b)}
 \xleftarrow{\sim}\mathcal Z^{(1)}
 \xleftarrow[\ j\ ]{\sim}\mathcal Z_\kappa.
\]
The action is the quartet roof of
Theorem~\ref{v4-arith:centers:thm:quartet}.
It restricts along $C_R\to C_{\mathcal Z_\kappa}$ to the
pointed ghost action, and does not factor through $J$.
\end{theorem}

\begin{proof}
The free-field tensor comparison transports the preceding
strict products and action to factorizing coalgebras.
Let $\mathcal P$ be the category of nonempty finite ordered
sets and order-preserving surjections.
The two towers are functors on $\mathcal P^{\mathrm{op}}$,
with lax multiplication given by concatenation.
They take every arrow to an equivalence.

Their coherent descent can be seen in the operad for a
nonunital algebra and a left module.
The symmetric monoidal envelope of the pair of towers has
fibers that are finite products of $\mathcal P^{\mathrm{op}}$.
Each is contractible, since the singleton is final in
$\mathcal P$. Localizing at the fiber arrows therefore gives
the envelope of the algebra-module operad.
The strict tower data descend through this localization.
This proves associativity with all coherences, and gives
the product roofs in the statement.

The vacuum is a quasi-unit. Right multiplication is the
inclusion into the two-quartet term. Left multiplication is
the shift $r_1$. On the image of $j$, the inclusion agrees
with the merge. Two polynomial homotopies join the merge
to the shift:
\[
 \begin{array}{c|cc}
 &J&c\\ \hline
 &J-\kappa T((1-s)t_1+st_2)
 &-(1-s)\eta_1-s\eta_2-ds(t_2-t_1)\\
 &J+sp_1+ds\,\iota_1-\kappa Tt_2&-\eta_2.
 \end{array}
\]
They preserve brackets, differential, and commutation with
$K_2$.

The vacuum action on the module is also an equivalence,
since $r_1i_1=i_2$.
The quasi-unit theorems
(\emph{Higher Algebra}, Theorem 5.4.3.5 and
Proposition 5.4.3.14) refine both structures to unital ones,
uniquely up to contractible choice.

Transport along $j$ and $i_1$ gives the stated objects.
Let $m:W_1\to W_2$ be the merge map.
The identity
\[
 \mu_{1,1}(j(z)\otimes i_1(v))
 =:j(z)i_2(v):
 =m\bigl(:j(z)i_1(v):\bigr)
\]
identifies the action roof. Its restriction has
$c\mapsto-\eta_1$ and $J\mapsto J+p_1$.
Proposition~\ref{v4-arith:centers:prop:ghost-detection}
excludes factorization through $J$.
\end{proof}

The action at one position is the full evaluation of the
moving-vacuum internal Hom. Indeed, its adjoint is
$z\mapsto j(z)$ into the resolved invariants.
In their continuous cochains, compare
\[
 f_0:(J,c)\mapsto(J+p,c),\qquad
 f_1:(J,c)\mapsto(J-\kappa Tt,-\eta).
\]
They are joined through $(J,c)\mapsto(J,c)$ by
\[
 (J+sp+ds\,\iota,c),\qquad
 (J-s\kappa Tt,(1-s)c-s\eta-ds\,t).
\]
Here the total cochain differential has $Dt=\eta+c$.

Integrating the $ds$ coefficient gives a chain homotopy
commuting with translation. Thus the adjoint is the identity
of $\mathcal Z^\ell$, and evaluation induces
\[
 R\operatorname{Hom}_{D_X}(N,\mathcal Z^\ell)
 \simeq
 R\operatorname{Hom}_{\mathrm{mov}}
 (N\otimes_{\mathcal O_X}\mathcal V,\mathcal V)
\]
for every complex $N$. This universal property at one
position does not establish finality of the factorizing actor
on all configurations.

The order of the coherent product can also be checked on a
common resolution. Put
$\mathcal T=C^\bullet_{\mathrm{sm}}(W_2)=W_2\otimes\Lambda[c]$.
Its differential satisfies
\[
 DJ=\kappa Tc,\quad Dt_l=\eta_l+c,\quad
 D\iota_l=p_l,\quad Dp_l=D\eta_l=Dc=0.
\]
The tower roof and the exchanged binary roof map into
$\mathcal T$ through quasi-isomorphisms. Their unit
maps are joined by the two homotopies used above.

A family of two-input maps is given by
\[
 \begin{split}
 J^{(i)}&\longmapsto
 J+\sum_l\bigl(\beta_l^{(i)}p_l+
 \kappa\epsilon_l^{(i)}Tt_l+
 ds\,\partial_s\beta_l^{(i)}\iota_l\bigr),\\
 c^{(i)}&\longmapsto
 (1+\epsilon_1^{(i)}+\epsilon_2^{(i)})c+
 \sum_l\bigl(\epsilon_l^{(i)}\eta_l+
 ds\,\partial_s\epsilon_l^{(i)}t_l\bigr).
 \end{split}
\]
It preserves brackets precisely when
\[
 \beta^{(i)}\cdot\epsilon^{(i)}=0,\qquad
 \beta^{(1)}\cdot\epsilon^{(2)}
 +\beta^{(2)}\cdot\epsilon^{(1)}=-1,
\]
and every cross pairing is constant in $s$.

The tower starts at
\[
 (\beta^{(1)},\epsilon^{(1)},\beta^{(2)},\epsilon^{(2)})
 =((0,0),(-1,0),(1,0),(0,-1)).
\]
Change successively
\[
 \epsilon^{(2)}=(0,s-1),\quad
 \beta^{(2)}=(1,2s),\quad
 \epsilon^{(1)}=(-1+2s,-s),\quad
 \beta^{(1)}=(s,s),
\]
keeping the other three vectors at their current values.
Each segment satisfies the bracket conditions.
The endpoint is the exchanged binary map.
Thus the tower product equals the opposite binary product
in the localization at quasi-isomorphisms. Its induced
cohomology product remains the exterior product.

\subsection{A coherent loop that prevents finality}

The cochain monoid has the expected moving-vacuum evaluation.
A local center must also represent the homotopies of actions
on separated configurations. An actor in degree two detects
a homotopy that the cochain monoid misses.

Let
\[
 N=k[t_0,t_1,\ldots],\qquad |t_n|=2,\qquad
 Tt_n=t_{n+1},\qquad d=0,
\]
with zero vertex brackets. Its augmentation $\epsilon$ kills
the $t_n$. The coalgebra $F_2=C_N$ is a commutative transverse
monoid. Its lattice on $X^I$ is
$\operatorname{Sym}_{\mathcal O_{X^I}}(H_I[-2])$.
It is $K$-flat by the same finite-jet argument as for the ghost,
and gives natural equivalences
\[
 F_2^{\otimes^!p}\otimes^!C
 \simeq C_{N^{\otimes p}\otimes V_\kappa}.
\]
Let $\alpha_0$ be the augmentation action and
$g_0:F_2\to J\to C_{\mathcal Z_\kappa}$ its factorization
through the unit.

Put $\mathfrak h_N=L_{N\otimes V_\kappa}$ and
$\mathfrak g=L_{V_\kappa}$. The derived deformation complex
at $\alpha_0$ is the relative convolution complex
\[
 \mathcal T_{\mathrm{rel}}=
 \operatorname{Hom}^{\mathrm{der}}
 \bigl(\overline{C^{\mathrm{CE}}_\star(\mathfrak h_N)},
       U_\star\mathfrak g\bigr)_{\mathrm{rel}},
 \qquad d_{\tau_0}=d+[\tau_0,-].
\]
Here \emph{relative} means vanishing on
$C^{\mathrm{CE}}_\star(\mathfrak g)$, and all coefficient
Homs remain derived. Completion is by products over the
total polynomial degree $G$ in the $t_n$. The base twisting
cochain $\tau_0$ is the augmentation on singleton inputs.

This filtration is complete, starts at $G=1$, and is added
by convolution. Thus the complex is a pronilpotent differential
graded Lie algebra. Derived bar--cobar identifies
\[
 \pi_1\Phi_C(F_2,u)=
 \bigl(H^0\mathcal T_{\mathrm{rel}},\mathrm{BCH}\bigr),
 \qquad
 \pi_n\Phi_C(F_2,u)=H^{1-n}\mathcal T_{\mathrm{rel}}
 \quad(n\geq2)
\]
at the augmentation. A closed degree-zero cochain $h$ gives
the loop $\tau_0+ds\,h$, with class $\exp[h]$.

Because $\tau_0$ is Lie-valued, PBW output degree is preserved.
In output degree $a$, a component has a partition into $a$
blocks, determined on their distinct locus. Its differential
acts inside each block. This assertion uses support and
localization, rather than a factorization of Hom for
infinite-dimensional coefficient modules.

\begin{lemma}[The degree bound]
\label{v4-arith:centers:lem:loop-degree}
In total $t$-degree $G$ and output degree $a$, the complex
is concentrated in degrees at least $a-2G$.
For $G=1,a=2$, closed degree-zero cochains have ordinary
binary coefficients, and no degree-zero boundaries occur.
\end{lemma}

\begin{proof}
With $n$ inputs and coefficient extension degree $e$,
suspension changes the input degree from $2G$ to $2G-n$.
The cochain degree is therefore $e+n-2G$.
The underlying coefficient modules admit derived Hom in
degrees $e\geq0$. Support forces $n\geq a$: a source on
at most $n$ blocks restricts to zero on an $a$-point locus
when $n<a$. Derived adjunction gives the same vanishing.

At each finite configuration the number of blocks is finite.
Products and the configuration limit preserve the bound.
For $G=1,a=2$, degree zero forces $n=2,e=0$, and degree
minus one does not occur.
\end{proof}

Define the augmentation derivation
$\ell:N\to k$ by $\ell(t_n)=\delta_{n0}$, vanishing on
other polynomial degrees. It obeys
\[
 \ell(Ta)=0,\qquad
 \ell(ab)=\ell(a)\epsilon(b)+\epsilon(a)\ell(b).
\]
The corresponding singleton cochain is
\[
 \kappa_1\bigl(s(t_n\otimes v)\bigr)=\delta_{n0}v.
\]
It has degree $-1$ and is closed.

Indeed, $t_0\mapsto u$, $t_{n>0}\mapsto0$ defines the
first-order vertex map with $|u|=2$, $u^2=0$.
The coefficient of $u$ in its Maurer--Cartan equation is
$d_{\tau_0}\kappa_1=0$.
The degree bound excludes boundaries below it, and its
value on $s(t_0\otimes\mathbf1)$ is nonzero.

On the vacuum, $J_1=\kappa\partial_{x_1}$ commutes with $T$
and is a derivation of every vertex operation.
It kills $\mathbf1$ and sends $J$ to $\kappa\mathbf1$.
Define the two-output cochain
\[
 h=\kappa_1\boxtimes J_1
   =-\bigl(\kappa_1\ast(J_1\tau_0)\bigr)^{(2)}.
\]
On two separated singleton blocks its value is
\[
 h\bigl(s(t_n\otimes v),s(\mathbf1\otimes w)\bigr)
 =\delta_{n0}\,v\boxtimes J_1w.
\]
The negative convolution sign cancels the sign from passing
the degree-one second cochain across the degree-one first
input. Both factors are closed, so $h$ is closed of degree zero.

\begin{lemma}[The separated vacuum coefficient]
\label{v4-arith:centers:lem:vacuum-probe}
There is a degree-minus-one chain map
$\mathsf P:\mathcal T_{\mathrm{rel}}^{(a)}
\to\mathcal T_{\mathrm{rel}}^{(a-1)}$ with
\[
 d_{\tau_0}\mathsf P+\mathsf P d_{\tau_0}=0,\qquad
 \mathsf P(h)=\kappa\kappa_1.
\]
Its value on the degree-zero, $G=1$, two-output part of
every class induced by a pointed map to
$C_{\mathcal Z_\kappa}$ is zero.
\end{lemma}

\begin{proof}
Insert a separated singleton at $z$ carrying
$\mathbf1\otimes J$, use test function one there, and apply
the vacuum coefficient to its output. Expand the resulting
rational function at $z=\infty$ and take its constant term.
Place the suspended current first, including its permutation
sign. On a degree-$r$ cochain multiply this operation by
$(-1)^{r+1}$. This defines $\mathsf P$.

The input $J$ is a free translation generator, since it has
weight one and $TV_+$ has no weight-one state.
Hence insertion is defined on derived coefficient resolutions.
Extraction commutes with the coefficient differential and
with derivatives and residues in the remaining variables.
A bracket inside the singleton block has no second input.
The two products with $\tau_0(J)=J$ cancel on the separated
locus. The remaining differential is the one on the other
blocks.

Insertion has degree $-1$. The displayed sign gives
the asserted anticommutation. These operations preserve
finite normal order and $G$, so extend to the completion.

For $h$, the inserted current gives $J_1J=\kappa\mathbf1$,
constant in $z$. The permutation sign cancels the sign in
$\mathsf P$, leaving $\kappa\kappa_1$.

For the image assertion, use the quartet commutant model
of the cochain action. Its transferred twisting cochain is
Lie-valued and supported on small diagonals. Its linear term
is the weight-preserving contraction of
\[
 m(\zeta,v)=\zeta_{(-1)}i(v).
\]

Here is a transfer retaining these properties.
Take the quartet contraction $(i,\pi,H)$ in coordinates
that keep $i(V_\kappa)$ fixed. Then
$QH+HQ=1-i\pi$, $Hi=\pi H=H^2=0$, and $H\mathbf1=0$.

On $n$ suspended factors let $K_n$ be the signed sum
of the shifted $H$, divided by the number of factors
of positive quartet degree. Set it to zero when that
number is zero. It satisfies
\[
 d_QK_n+K_nd_Q=1-\operatorname{Sym}^n(i\pi).
\]
With $P_1=s\pi s^{-1}$, define recursively
\[
 P_n=\left(
 \frac12\sum_{n_1+n_2=n}
 \beta_V(P_{n_1}\otimes P_{n_2})\Delta_{n_1,n_2}
 -P_{n-1}\delta_W\right)K_n.
\]
Here $\beta_V$ is the suspended bracket, $\delta_W$
the Chevalley bracket differential, and $\Delta$ the
unshuffle coproduct.

The expression in parentheses is
a cycle and vanishes on $\operatorname{Sym}^n(i)$.
The contraction identity therefore proves the coalgebra-map
equation inductively. It gives $P\operatorname{CE}(i)=1$.
Each component is a finite sum of bracket and contraction
composites, with values in the Lie object.

The higher components vanish on a vacuum input.
Indeed, the vacuum has quartet degree zero, and its
bracket with the output equals the sum of its brackets
with the inputs. This is the partial-fraction residue
identity around those inputs. The two remaining terms
in the recursion cancel, with the suspended vacuum sign.
Thus $\sigma=\vartheta_VP\operatorname{CE}(m)$ is the
required normalized twisting cochain, where $\vartheta_V$
is the canonical Lie-valued twisting cochain.

Free localized normal resolutions have surjective normal
derivative. Their inverse limits and products preserve
this comparison. Its derivatives retain the normalization.

Differentiate composition at the augmentation.
At a separated input carrying $\mathbf1\otimes J$, either
the base point supplies $J$, or the variable block supplies
$\zeta_{(-1)}(J+p)$ before contraction.
Both have positive conformal weight. Their vacuum coefficient
vanishes, since the resolution has nonnegative weights and
the contraction preserves weight.

This is also a derived vanishing: the map factors through
the subcomplex of positive weights before applying
coefficient resolutions. By the degree bound, its closed
degree-zero two-output representative is ordinary.
Its vacuum coefficient is already zero before taking
the constant term at infinity. This proves the assertion.
\end{proof}

\begin{theorem}[Failure of finality of the cochain actor]
\label{v4-arith:centers:thm:nonfinal}
The mapping space from the augmentation action
$(F_2,\alpha_0)$ to the cochain action
$(C_{\mathcal Z_\kappa},\alpha_{\mathcal Z})$
has at least two connected components.
Consequently this cochain actor is not a local center
on the full category of factorizing coalgebras.
\end{theorem}

\begin{proof}
The class $[h]$ is missed by the map on pointed-action
fundamental groups. Its probe is the nonzero class
$\kappa[\kappa_1]$, whereas every image has zero probe.
A source loop is represented by $ds\,K'$ with $K'$ closed
of degree zero. Only the derivative of composition contributes,
since $(ds)^2=0$. The exponential identification for the
complete Lie algebra therefore shows that $\exp[h]$
is not an image loop.

We must also retain action coherences. For $p$ actor factors put
\[
 \epsilon_p(a_1,\ldots,a_p)=\prod_i\epsilon(a_i),\qquad
 \ell_p(a_1,\ldots,a_p)
 =\sum_i\ell(a_i)\prod_{j\ne i}\epsilon(a_j).
\]
Let $\kappa_p$ be the resulting singleton cochain and
$\tau_0^{(p)}$ the augmentation twisting cochain.
Then
\[
 h_p=-\bigl(\kappa_p\ast(J_1\tau_0^{(p)})\bigr)^{(2)}
\]
is closed. Over a square-zero parameter $\delta$ of degree one,
$\tau_0^{(p)}+\delta h_p$ is a twisting cochain.

Multiplication faces use the derivation rule for $\ell$.
Unit faces use $\ell(1)=0$. Composition uses
\[
 \ell_{p+q}=\ell_p\otimes\epsilon_q+
            \epsilon_p\otimes\ell_q.
\]
Terms with two copies of $\delta$ vanish.
On every parenthesized word, the linear term is therefore
the same sum over the factor on which $\ell$ acts.

The free associative bar resolution has constant compatible
fillers on all simplices and its usual extra unit degeneracy.
It gives all action coherences. Extending $\delta\mapsto ds$
produces a loop of transverse actions based at $\alpha_0$,
whose underlying pointed loop is $\exp[h]$.

The mapping space in the statement is the homotopy fiber
over $\alpha_0$ of
\[
 \operatorname{Map}_{\mathrm{Alg}}
 (F_2,C_{\mathcal Z_\kappa})
 \longrightarrow
 \operatorname{Act}(F_2,C).
\]
It contains $g_0$. The loop just constructed is not the
image of a loop on the left, since forgetting multiplication
would contradict the pointed obstruction. The homotopy-fiber
sequence therefore gives at least two components.
A final action would make this fiber contractible.
\end{proof}

This obstruction concerns the specified cochain actor.
It does not exclude every possible local center, or claim
that individual actions cannot lift. The missing homotopy
changes the external-product identifications at a separated
boundary input. Fixing those identifications coherently
defines a further representability problem. So does
constructing a larger actor that retains their deformations.

\section{The marked center of the affine \texorpdfstring{$E_6$}{E6} category}
\label{v4-arith:centers:sec:e6}

A finite tensor category supplies a different central action.
Its sectors may change the boundary object. The affine
$E_6$ realization makes that action explicit, including
the algebra that recovers the boundary and the fields
that implement its junctions.

Work over $\mathbb C$ at the unitary embedding of
Theorem~\ref{v4-arith:asi:e6-analytic-comparison}.
Put
\[
 V=L_{\widehat{\mathfrak{sl}}_2}(10,0),\qquad
 \mathcal D=\operatorname{Rep}(V).
\]
Its simple modules $H_j$, $0\leq j\leq10$, have lowest weights
$h_j=j(j+2)/48$. Retain the connected commutative separable
algebra $A=H_0\oplus H_6$ with its specified multiplication.

Let $\mathcal M=\mathcal D_A$ be the category of right
$A$-modules. Our balanced tensor gives a right module $N$
the left action
\[
 \ell_N=r_Nc_{A,N}.
\]
This convention is part of the center computation.

The even component $\mathcal C$ has simples $1,g,\rho$ and
\[
 g^2=1,\qquad g\rho=\rho g=\rho,\qquad
 \rho^2=1+g+2\rho.
\]
With $d=1+\sqrt3$, their dimensions are $1,1,d$, and
$\operatorname{Dim}\mathcal C=6+2\sqrt3$.
The local modules form
$\mathcal E=\operatorname{Rep}(W)$ for
$W=L_{\widehat{\mathfrak{so}}_5}(1,0)$:
\[
 W_0=H_0\oplus H_6,\quad
 W_v=H_4\oplus H_{10},\quad
 W_s=H_3\oplus H_7.
\]
Their dimensions are $1,1,\sqrt2$ and their lowest weights
are $0,1/2,5/16$.

Let $S=L_{\widehat{\mathfrak{so}}_{11}}(1,0)$.
Its three simples $S_0,S_v,S_s$ have the same dimensions and
weights $0,1/2,11/16$. The affine realization identifies
$\operatorname{Rep}(S)$ with $\mathcal E^{\mathrm{rev}}$,
where reversal means
$c^{\mathrm{rev}}_{X,Y}=c_{Y,X}^{-1}$.
This is the reversed Ising category in Bischoff's
center realization, Corollary 4.7 and Section 5.

Set $T=V\otimes S$ and
$\mathcal T=\mathcal D\boxtimes\operatorname{Rep}(S)$.
The invertible object $J_*=H_{10}S_v$ has order two and
integral weight three. Choose the unitary commutative algebra
$L=1\oplus J_*$ with its vacuum unit.
Its vertex extension is
\[
 U=H_0S_0\oplus H_{10}S_v,\qquad c_U=8.
\]

Here $H_jS_a$ abbreviates the external tensor module
$H_j\otimes S_a$. The category tensor is the specified
vertex tensor product, not the algebraic tensor of states.
The tensor-product and extension theorems apply to these
strongly rational affine algebras.

An object of the \emph{Drinfeld center} $Z(\mathcal C)$
is an object $M$ with a natural half-braiding
$\gamma_N:M\otimes_A N\to N\otimes_A M$
satisfying the hexagon. Its forgetful functor retains
the object while forgetting this family.

\begin{theorem}[Ten central sectors]
\label{v4-arith:centers:thm:e6-center}
There is a braided equivalence
\[
 \Psi:\operatorname{Rep}(U)^{\mathrm{rev}}
 \xrightarrow{\sim}Z(\mathcal C).
\]
The central forgetful functor $F$ is obtained before descent
from
\[
 \widetilde F(X,Y)=(X\otimes A)\otimes_A Y,
 \qquad X\in\mathcal D,\quad Y\in\mathcal E.
\]
Its half-braiding with a right $A$-module $N$ is induced by
\[
 X\otimes Y\otimes N
 \xrightarrow{\,1\otimes c_{Y,N}\,}X\otimes N\otimes Y
 \xrightarrow{\,c_{N,X}^{-1}\otimes1\,}N\otimes X\otimes Y.
\]
Associators are inserted between these displayed tensors.
Restrict to $\mathcal C$ and descend by the trivial character
of $\widetilde F(L)\simeq\mathbb C[\mathbb Z/2]\otimes1$.

The ten simples and their forgetful images are
\[
\begin{array}{c|c|c|c|c}
 &\text{restriction to }T&F&\dim&h\bmod1\\ \hline
 x_0&H_0S_0\oplus H_{10}S_v&1&1&0\\
 x_2&H_2S_0\oplus H_8S_v&\rho&d&1/6\\
 x_4&H_4S_0\oplus H_6S_v&g+\rho&d+1&1/2\\
 x_6&H_6S_0\oplus H_4S_v&1+\rho&d+1&0\\
 x_8&H_8S_0\oplus H_2S_v&\rho&d&2/3\\
 x_{10}&H_{10}S_0\oplus H_0S_v&g&1&1/2\\
 y_1&(H_1\oplus H_9)S_s&\rho&d&3/4\\
 y_3&(H_3\oplus H_7)S_s&1+g+\rho&d+2&0\\
 z_+&H_5S_s&\rho&d&5/12\\
 z_-&H_5S_s&\rho&d&5/12 .
\end{array}
\]
Sums in this table are direct sums.
The twist in $Z(\mathcal C)$ is $e^{-2\pi i h}$.
The universal categorical central action is
$(Z,M)\mapsto F(Z)\otimes_A M$ with this half-braiding.
\end{theorem}

\begin{proof}
Apply the center theorem for a connected étale algebra to
$\mathcal D^{\mathrm{rev}}$.
With our module tensor convention, the theorem of
Davydov--Müger--Nikshych--Ostrik gives
\[
 \mathcal D^{\mathrm{rev}}\boxtimes\mathcal E
 \simeq Z(\mathcal M).
\]
Its two crossings give the displayed formula.
Locality of $Y$, expressed by
$r_Yc_{A,Y}c_{Y,A}=r_Y$, preserves the balancing relations.
The braiding hexagons then give the half-braiding hexagon.

The unit checks the orientation. For $Y=N=A$ the inverse
crossing cancels the crossing in the left unit constraint.
A positive crossing would leave the double braiding
$c_{A,X}c_{X,A}$. On the $H_6$ channel of
$H_2\boxtimes H_6$ its scalar is
\[
 e^{2\pi i(h_6-h_2-h_6)}=e^{-\pi i/3}\ne1.
\]
It would therefore fail the unit constraint.

The central object $H_{10}S_v$ forgets to $g\otimes_A g=1$.
Its central structure is the nontrivial parity character.
The fusion rules force $g,\rho$ to have neutral degree in
any grading, while $H_1\otimes A$ generates the odd part.
Thus the universal grading of $\mathcal M$ is $\mathbb Z/2$.
The center theorem for graded fusion categories
(Gelaki--Naidu--Nikshych, Corollary 3.7)
identifies the neutral de-equivariantized center with
$Z(\mathcal C)$.

Concretely, an $L$-action becomes an involution $t$ after
applying $\widetilde F$. The idempotent $(1+t)/2$ selects
its vacuum character. The half-braiding of $J_*$ is trivial
on $\mathcal C$, so it descends to this image.
The vertex extension theorem identifies
$\mathcal T_L^0$ with $\operatorname{Rep}(U)$ as a braided
tensor category. Reversing the ambient braiding reverses
the local-module braiding, proving the claimed orientation.

The monodromy of $J_*$ on $H_jS_a$ is $(-1)^j$ for
$a=0,v$, and $(-1)^{j+1}$ for $a=s$.
There are seventeen local simples before extension.
Eight orbits have size two. The fixed object $H_5S_s$
has two equivariant structures, producing $z_+,z_-$.
The multiplication of $L$ fixes $t^2=1$. Changing the
sign of its normalized fixed-point intertwiner exchanges
these two labels.

To compute $F$, let $\widetilde I$ be right adjoint to
$\widetilde F$. Rigidity and induction give
\[
 \operatorname{Hom}_{\mathcal M}(\widetilde F(X,Y),1)
 \simeq
 \operatorname{Hom}_{\mathcal D}(X,\operatorname{Res}Y^*).
\]
Consequently
\[
 \widetilde I(1)=
 (H_0\oplus H_6)S_0\oplus
 (H_4\oplus H_{10})S_v\oplus
 (H_3\oplus H_7)S_s.
\]

The images of $H_0S_0,H_{10}S_0,H_2S_0$ are $1,g,\rho$.
By the projection formula, the multiplicity of one of these
images in $\widetilde F(Q)$ is the coefficient of its
corresponding $T$-module in $\widetilde I(1)\boxtimes Q$.
Use
\[
 H_a\boxtimes H_b=
 \bigoplus_{\substack{|a-b|\leq j\leq\min(a+b,20-a-b)\\
                      j\equiv a+b\ (2)}}H_j
\]
and $S_v^2=S_0$, $S_vS_s=S_s$, $S_s^2=S_0+S_v$.
This gives the forgetful column. At the fixed point the
image is $2\rho$. Each resolved summand has dimension $d$
and therefore image $\rho$.

Dimensions and weights follow by induction.
The sum of their squared dimensions is
$(6+2\sqrt3)^2$. Finally, a central tensor functor is
exactly a tensor functor with the specified half-braidings.
Its factorization through $Z(\mathcal C)$ gives the
universal categorical action.
\end{proof}

This action has actual affine fields.
For $p:F(Z)\otimes_A M\to N$, take the intertwining operator
associated with
\[
 \operatorname{Res}F(Z)\boxtimes\operatorname{Res}M
 \longrightarrow\operatorname{Res}(F(Z)\otimes_A M)
 \xrightarrow{\operatorname{Res}p}\operatorname{Res}N.
\]
The first arrow is the balanced quotient defined by the
separability idempotent.
Vertex tensor associativity transports its identities
to every descendant coefficient.
The $X$ constituent crosses clockwise and the local $Y$
constituent counterclockwise.

The parameter states are those of $\operatorname{Res}F(Z)$.
A field with the states of the $U$-module $Z$ has another
domain.

\subsection{The distinguished algebra and the boundary sectors}

Let $I$ be right adjoint to $F$ and put $B=I(1)$.
Its multiplication is part of the central data:
if $\epsilon:F(B)\to1$ is the counit, then
$\mu_B:B\boxtimes B\to B$ is adjoint to
\[
 F(B\boxtimes B)\simeq F(B)\otimes_A F(B)
 \xrightarrow{\epsilon\otimes\epsilon}1.
\]
The unit is adjoint to $1_1$. A direct-sum decomposition
alone would not specify this algebra.

\begin{proposition}[The marked algebra]
\label{v4-arith:centers:prop:e6-marking}
With the preceding multiplication,
\[
 B=x_0\oplus x_6\oplus y_3,\qquad
 F(B)=3\,1\oplus g\oplus2\rho.
\]
It is connected, commutative, separable, of dimension
$6+2\sqrt3$, and has trivial twist.
Using the same positive balancing convention on right
$B$-modules,
\[
 \operatorname{Rep}(U)_B\simeq\mathcal C,\qquad
 \operatorname{Rep}(U)_B^0\simeq\operatorname{Vec}.
\]
Its vertex extension is the holomorphic algebra
$L_{\widehat{\mathfrak e}_8}(1,0)$.
\end{proposition}

\begin{proof}
Adjunction identifies the multiplicity of $Z$ in $I(1)$
with the multiplicity of $1$ in $F(Z)$.
The table gives the displayed sum.
The adjoint of the central forgetful functor carries
its canonical connected commutative separable algebra,
with the multiplication just defined.
The canonical-algebra theorem recovers $\mathcal C$
as its module category.

Commutativity for a braiding also implies commutativity
for its reverse. The canonical module tensor uses
$c_{N,B}^{-1}$ in the central category.
Since that category is $\operatorname{Rep}(U)^{\mathrm{rev}}$,
this is the positive crossing $c_{B,N}$ of $U$-modules.
The identity
$\dim(B)^2=\operatorname{Dim}\operatorname{Rep}(U)$
makes the local module category trivial.
All three summands of $B$ have integral weight.

The vertex extension theorem constructs a simple CFT-type
algebra $\mathbb V_B$ with multiplication $\mu_B$.
Its restriction to $T$ is $\widetilde I(1)$.
It has a one-dimensional weight-zero space. The other
two $U$-summands begin in weight one.

Apply Creutzig--Kanade--McRae's gluing theorem,
Theorem 5.12 and Corollary 5.13, with $S$ as the first factor.
Its distinct simple modules pair with the $V$-modules
$\operatorname{Res}W_a$, and
\[
 \dim\operatorname{Hom}_V(V,\operatorname{Res}W_a)=\delta_{a0}.
\]
It proves strong rationality of this extension.

Every irreducible extension module is grading-restricted.
Its restriction to the regular algebra $T$ has finite
multiplicities, since each conformal-weight space is finite.
The local-module equivalence therefore places it in
$\operatorname{Rep}(U)_B^0=\operatorname{Vec}$.
Thus the vacuum is the only irreducible module.

The central charge is eight, so Dong--Mason's
classification, Theorem 1, identifies it with
$L_{\widehat{\mathfrak e}_8}(1,0)$.
This constructs the embedding from $\mu_B$.
It requires no uniqueness assertion about another
embedding into the same holomorphic algebra.
\end{proof}

The other simple right $B$-modules are determined by
the same adjunction:
\[
 \begin{split}
 I(g)&=x_{10}\oplus x_4\oplus y_3,\\
 I(\rho)&=x_2\oplus x_4\oplus x_6\oplus x_8
 \oplus y_1\oplus y_3\oplus z_+\oplus z_-.
 \end{split}
\]
Their complete restrictions are
\[
 \begin{aligned}
 I(1)|_T={}&(H_0\oplus H_6)S_0
 +(H_4\oplus H_{10})S_v+(H_3\oplus H_7)S_s,\\
 I(g)|_T={}&(H_4\oplus H_{10})S_0
 +(H_0\oplus H_6)S_v+(H_3\oplus H_7)S_s,\\
 I(\rho)|_T={}&
 (H_2\oplus H_4\oplus H_6\oplus H_8)(S_0\oplus S_v)\\
 &+(H_1\oplus H_3\oplus2H_5\oplus H_7\oplus H_9)S_s .
 \end{aligned}
\]
These are full module identities, including descendants.
Their $S_0$-multiplicity spaces are the original three
affine sectors. Only $I(1)$ is local.
The other two are ordered boundary sectors, rather than
ordinary modules of the holomorphic extension.

The extension also realizes fields with full $U$-module
inputs. For $K_M=I(M)$, the projection formula is
\[
 \kappa_{Z,M}:Z\boxtimes K_M
 \xrightarrow{\sim}I(F(Z)\otimes_A M).
\]
It is adjoint to $1\otimes\epsilon_M$.
Rigidity and adjunction identify, for every $Q$,
\[
 \begin{split}
 \operatorname{Hom}(Q,Z\boxtimes I(M))
 &\simeq
 \operatorname{Hom}_{\mathcal C}
 (F(Z)^*\otimes_A F(Q),M)\\
 &\simeq
 \operatorname{Hom}(Q,I(F(Z)\otimes_A M)).
 \end{split}
\]
Hence $\kappa_{Z,M}$ is an isomorphism.
The counit construction makes it $B$-linear and compatible
with iterated tensor products.

For a junction $p:F(Z)\otimes_A M\to N$, the field is
\[
 \mathcal J_p=\mathcal Y_{\,I(p)\kappa_{Z,M}}:
 Z\otimes K_M\longrightarrow\overline{K_N}\{u\}.
\]
The bar completes the output by conformal weight.
The braces allow the prescribed fractional powers.
The vertex tensor correspondence gives all descendant
coefficients, associativity, and the central crossing.
At $Z=U$ this is the ordinary module field $Y_{K_M}$.
It can have singular products.

The mirror states cannot be removed by an operation-preserving
vacuum projection. Already on $V\otimes S$, let
$P=1_V\otimes\epsilon_S$ project to the mirror vacuum.
For normalized mirror currents $a,b$,
\[
 P(a_{(1)}b)=1_V,\qquad P(a)_{(1)}P(b)=0.
\]
The extension has weight-one dimension
$3+55+7+5\cdot11+4\cdot32=248$. The mirror currents alone
contribute $55$ states. The multiplicity-space identification
does not delete their vertex operations.

\subsection{The pointed field action and its obstruction}

For a vertex algebra $Q$, define
\[
 Z_{\mathrm{vert}}(Q)=
 \{q:q_{(n)}v=0\text{ for every }v\in Q,\ n\geq0\}.
\]
A pointed action of $Q$ on $W$ means a unital vertex map
$Q\otimes W\to W$ restricting to the identity on
$1_Q\otimes W$. It need not preserve a conformal vector.

\begin{proposition}[The current double-pole obstruction]
\label{v4-arith:centers:prop:e6-pointed}
Pointed vertex maps $Q\otimes W\to W$ are in bijection
with vertex maps $Q\to\mathbb C1$.
In particular, there is no pointed map $U\otimes W\to W$.
The derived pointed mapping space
\[
 \operatorname{Map}_{C_W/}(C_U\otimes^!C_W,C_W)
\]
is empty.
\end{proposition}

\begin{proof}
If $w$ is central, skewsymmetry makes $Y(v,z)w$ regular.
Taking $v=\omega_W$ gives $L_0w=0$.
Since $W$ is of CFT type, $Z_{\mathrm{vert}}(W)=\mathbb C1$.

The image of $q\otimes1$ under a pointed map must commute
with $1\otimes W$ in every nonnegative product.
It therefore lies in this center.
Conversely, a vertex map $\varphi:Q\to\mathbb C1$ defines
the pointed action $q\otimes w\mapsto\varphi(q)_{(-1)}w$.
These constructions are inverse.

Take standard currents $e,f$ in
$L_{\widehat{\mathfrak{sl}}_2}(10,0)\subset U$, with
$(e,f)=1$. They satisfy
\[
 e_{(1)}f=10\,1_U.
\]
All nonnegative products of scalar vacuum states vanish.
A unital map $U\to\mathbb C1$ would send this identity
to $10=0$, which is impossible.

The affine vertex algebras give flat chiral algebras.
The heart theorem identifies their pointed factorization
mapping space with the discrete space of chiral maps
from $U$ to the chiral center of $W$.
The same current relation excludes such a map locally.
Thus no higher coherence supplies a missing point.
Replacing the entire pointed diagram by equivalent
objects preserves this empty mapping space.
\end{proof}

A regular action of a commutative vertex algebra on all
three boundary sectors, compatible with their ordered
junctions, is consequently scalar.
Its restriction to $W$ is a map of commutative differential
algebras $\varphi:R^\ell\to\mathcal O_X$.
Inserting the vacuum in the junction $W\otimes_A M\to M$
forces the action on every $M$ to be
$r\otimes m\mapsto\varphi(r)m$.
If it commutes with translations, $\varphi$ comes from
a vertex map $R\to\mathbb C$.

The categorical central action remains valid.
It takes a central sector to a tensor functor on boundary
objects, often changing the object. Its universal property
does not give a pointed field map on the fixed vacuum $W$.
The ordinary transverse endomorphism monoid of $C_W$ also
exists by the earlier construction. The extra factorization
and pointing conditions determine the obstruction just proved.

\subsection{Derived current classes}

Scalar regular actions do not exhaust derived central
cohomology. Let $A_W=\Omega_{\mathrm{cut}}C(W)$ be the word
cobar of the ordered residue coalgebra. Its scalar word
center uses product totalization in word size and
Hochschild length, with algebraic input state spaces.

\begin{proposition}[Polynomial current derivations]
\label{v4-arith:centers:prop:e6-derived}
The scalar word center satisfies
\[
 H^{i<0}Z_{\mathrm{word}}(A_W)=0,\qquad
 H^0Z_{\mathrm{word}}(A_W)=\mathbb C.
\]
There is a natural injection of Lie algebras
\[
 \mathfrak{so}_5[z]\hookrightarrow
 H^1Z_{\mathrm{word}}(A_W),
 \qquad
 x\otimes f\longmapsto
 \delta_{x,f}(z)=
 \sum_{n\geq0}\frac{f^{(n)}(z)}{n!}x_{(n)}.
\]
On a current state $a$,
\[
 \delta_{x,f}(a)=f(z)[x,a]+f'(z)(x,a)\mathbf1.
\]
The invariant form gives long roots squared length two.
Constant conformal-weight-preserving vertex derivations
form the ten-dimensional subspace of constant currents.
\end{proposition}

\begin{proof}
The negative and degree-zero assertions are the general
word-center calculation above.
The translation relation gives
\[
 \partial_z\delta_{x,f}+[T,\delta_{x,f}]=0.
\]
The current commutator formula and Taylor expansion give
\[
 [\delta_{x,f}(z),Y(a,u)]
 =Y(\delta_{x,f}(z+u)a,u).
\]
These are precisely the chiral derivation equations.
Summing the operator over state blocks commutes with
all mergers, cuts, and normal derivatives.
The all-word derivation identification therefore gives
the class, with no inner-derivation quotient.

The affine relation at level one gives its value on $a$.
If a polynomial combination vanishes, its weight-one
part commutes with every element of $\mathfrak{so}_5$.
The Lie center of $\mathfrak{so}_5[z]$ is zero, proving
injectivity. For nonnegative modes the central term of
their commutator vanishes. The derivative Leibniz rule yields
\[
 [\delta_{x,f},\delta_{y,g}]=\delta_{[x,y],fg}.
\]

A constant weight-preserving derivation restricts to a
derivation of the simple Lie algebra $W_1$.
It is inner there and extends uniquely, because the
currents generate $W$. Thus it is $a_{(0)}$ for a unique
$a\in\mathfrak{so}_5$.
\end{proof}

The same construction with $\mathfrak{sl}_2[z]$ acts on
every affine boundary summand $H_j$.
For a junction intertwiner $\mathcal Y_p$, its Jacobi
identity reads
\[
 \delta_{x,f}^{N}(z)\mathcal Y_p(m,u)
 -\mathcal Y_p(m,u)\delta_{x,f}^{M}(z)
 =
 \mathcal Y_p(\delta_{x,f}^{F(Z)}(z+u)m,u).
\]
The balancing maps are affine module maps, so these
identities respect their composites and central crossing.
On the level-ten currents the scalar term is
$10f'(z)(x,a)\mathbf1$.
It retains the level in the field operations.
It does not prove that these classes survive a different
endomorphism complex for the full marked transverse center.

The remaining comparison must retain the root, full collision
atlas, and distinguished algebra $I(1)$ with its multiplication.
It also requires a passage to rational boundary sectors. The center equivalence supplies the
categorical action and its intertwiners.
The double-pole calculation excludes replacing this
passage by a pointed action of $C_U$ on the original vacuum.

\section{Representation tensors from specified fields}
\label{v4-arith:centers:sec:representation-tensors}

A field exchange becomes a representation intertwiner only
after the tensor action has been specified. Central contractions
give a model in which both the positive construction and the
possible mismatch can be computed on all modes.

Let $P=k[x_0,x_1,\ldots]$, with $Tx_i=x_{i+1}$ and
$\operatorname{wt}(x_i)=i+1$. Introduce a formal parameter
$\hbar$ and a central state $\kappa$ with
$T\kappa=d\kappa=0$. For a Laurent kernel $K(z)$, prescribe
\[
 Y_K(x_0,z)x_0=(e^{zT}x_0)x_0+\hbar\kappa K(z).
\]
The simple pole $K=z^{-1}$ is nonlocal. The symmetric
double pole $K=z^{-2}$ gives the usual central local
current bracket.

We retain a resolution of this polynomial algebra.
Let $A_{\mathrm{ext}}=\Lambda(\theta_i:i\geq0)$, where
$|\theta_i|=1$ and $\operatorname{wt}(\theta_i)=i+1$.
The weight-restricted dual
\[
 E=(B_{\mathrm{alg}}A_{\mathrm{ext}})^\circ
\]
is freely associative on $y_I$ for nonempty increasing
finite subsets $I$ of $\mathbb N$. Its differential is
\[
 |y_I|=1-|I|,\qquad
 dy_I=\sum_{I=A\sqcup B}
 (-1)^{|A|+1}\epsilon(A,B)y_Ay_B.
\]
The sum uses ordered nonempty parts, and $\epsilon(A,B)$
is their exterior permutation sign.
Thus $x_i=y_{\{i\}}$ and $dy_{ij}=x_ix_j-x_jx_i$.

\begin{lemma}[The finite-weight resolution]
\label{v4-arith:centers:lem:field-resolution}
Sending singleton generators to $x_i$ and higher generators
to zero gives a quasi-isomorphism $\pi:E\to P$.
It uses a direct sum of finite weight duals.
\end{lemma}

\begin{proof}
Let $\operatorname{Sym}^c(\bigoplus_i kt_i)$ have primitive
degree-zero generators of weights $i+1$. The map to the bar is
\[
 j(t_{i_1}\cdots t_{i_r})
 =\sum_{\sigma\in S_r}
 [s\theta_{i_{\sigma(1)}}|\cdots|s\theta_{i_{\sigma(r)}}].
\]
Adjacent interior products cancel in pairs, since the
$\theta_i$ are odd. Deconcatenation matches the unshuffle
coproduct, including repeated-index multiplicities.

The free $A_{\mathrm{ext}}$-module complex
$A_{\mathrm{ext}}\otimes k[t_i]$, with $dt_i=\theta_i$,
resolves $k$. Its contraction on positive total generator
degree $N$ is $N^{-1}\sum_i t_i\partial_{\theta_i}$.
The normalized bar is another semifree resolution.
The map $1\otimes j$ intertwines their differentials,
including the endpoint term $\sum_i\theta_i\partial_{t_i}$,
and lifts the identity on $k$. Lifting through the free
filtrations gives a homotopy inverse, weight by weight.

Every weight is finite-dimensional. Dualize there with
$\langle x^\alpha,t^\beta\rangle=\alpha!\delta_{\alpha\beta}$,
which converts the unshuffle coproduct into polynomial
multiplication. The dual is $\pi$. Taking the direct sum
of the weight comparisons proves the assertion.
\end{proof}

Extend translation to $E$ by raising each index of $y_I$
in turn, with a repeated index giving zero and exterior
reordering carrying its sign. Let
$\partial_m(y_I)=1$ for $I=\{m\}$ and zero otherwise.
These graded derivations satisfy
\[
 [d,T]=[d,\partial_m]=0,\qquad
 [T,\partial_m]=-\partial_{m-1},\qquad
 [\partial_m,\partial_n]=0,\qquad\partial_{-1}=0.
\]
On the bar, the transposed operators lower an exterior
index and insert $s\theta_m$ in every gap.
Those operations commute with its differential, proving
the identities. The map $\pi$ intertwines them.

Put
\[
 K_{mn}(z)=\partial_z^m(-\partial_z)^nK(z),\qquad
 C_K(z)=\hbar\kappa\sum_{m,n\geq0}
 K_{mn}(z)\partial_m\otimes\partial_n.
\]
With the ordered coefficient multiplication $\mu_E$, define
\[
 Y_K(a,z)b=\mu_E(e^{zT}\otimes1)e^{C_K(z)}(a\otimes b).
\]
Every contraction lowers polynomial degree in both inputs,
so the exponential is finite on polynomial states.
The same formula with $\mu_P$ defines the polynomial fields.

All parenthesized ordered products expand the common expression
\[
 \mu_r\left(\bigotimes_i e^{z_iT}\right)
 \exp\left(\sum_{i<j}C_K^{ij}(z_i-z_j)\right).
\]
The pair contractions commute, and their derivation rule
gives the product and iterate identities.
Thus both carriers are differential nonlocal vertex algebras,
and $\pi$ preserves every field mode.
The completion is
$(E\otimes k[\kappa])[[\hbar]]$: each parameter coefficient
has finite polynomial and energy support.
There is no unrestricted energy product.

For inputs of energies $a,b$, a term with $q$ contractions
and a kernel of pole order $p$ has pole order at most
$a+b+(p-2)q$. Modulo $\hbar^N$, also
$q\leq\min(a,b,N-1)$.
This proves the required Laurent bounds at each precision.

The regular action is faithful, since
$Y_E(a,z)1=e^{zT}a$ has constant term $a$.
Its operator-field multiplication uses regularized
substitution $w_1=w+u$ after multiplying by a sufficient
power of $w_1-w$. The common Wick expression proves
\[
 \mathcal Y(Y_E(a,w),u)Y_E(b,w)=Y_E(Y_E(a,u)b,w).
\]
This retains the full substitution, not only a principal part.

\subsection{The universal level and additive module tensor}

The coproduct is over $k[[\hbar]]$, with $\kappa$ a state:
\[
 \Delta\kappa=\kappa\otimes1+1\otimes\kappa,\qquad
 \Delta y_I=y_I\otimes1+1\otimes y_I.
\]
Its target has independent level states $\kappa_L,\kappa_R$.
The counit kills them and all $y_I$.

\begin{proposition}[The universal vertex bialgebra]
\label{v4-arith:centers:prop:universal-level}
These formulas make $(E\otimes k[\kappa])[[\hbar]]$
a counital differential nonlocal vertex bialgebra.
The map $\pi$ is a coalgebra map and an all-mode field
quasi-isomorphism. The induced module tensor adds levels.
\end{proposition}

\begin{proof}
Each $dy_I$ is a sum of graded commutators of primitive
generators, so $\Delta$ is a chain map.
Coassociativity and the coefficient counit hold on generators.
For each $m$,
\[
 (\partial_m\otimes1)\Delta
 =\Delta\partial_m=(1\otimes\partial_m)\Delta.
\]
Each outer derivative separately gives this identity.
Their sum would double it.

The target contractions therefore contribute
$\hbar(\kappa_L+\kappa_R)K_{mn}$, exactly the image
of the source contraction. Exponentiation proves
\[
 \Delta Y_K
 =(Y_K\otimes Y_K)(1\otimes\tau\otimes1)(\Delta\otimes\Delta).
\]
Every contracted term contains $\kappa$, so the counit
kills it. The uncontracted term is its ordinary augmentation.
Equivariance of $\pi$ proves the remaining claims.
\end{proof}

A nonzero fixed-level ideal $(\kappa-c)$ is not a coideal:
its image is $\kappa_L+\kappa_R-c$, while setting both
target levels to $c$ gives $c$.
This does not prevent maps from level $c_1+c_2$ to the
tensor of levels $c_1,c_2$.

An admissible module here is topologically free,
$M=M_0[[\hbar]]$, with continuous fields lower truncated
on each input modulo every $\hbar^N$.
For two such modules define
$M\widehat\boxtimes N=(M_0\otimes_kN_0)[[\hbar]]$ and
\[
 Y_{M\widehat\boxtimes N}(a,z)(m\otimes n)
 =\sum_{(a)}(-1)^{|a_{(2)}||m|}
 Y_M(a_{(1)},z)m\otimes Y_N(a_{(2)},z)n.
\]
At each precision only finitely many coproduct terms occur,
and two lower-bounded Laurent series have finite
coefficient convolution. The bialgebra identity proves
module associativity. Coassociativity proves the pentagon.

A derived version uses $K$-flat models and the specified
completed derived tensor. This is an additive-level
module tensor, with its stated carrier.

\subsection{Normal-ordering transport}

Let $H(z)\in k[[z]]$ be even and put
\[
 H_{mn}=(-1)^nH^{(m+n)}(0),\qquad
 D_H=\frac12\sum_{m,n}H_{mn}\partial_m\partial_n,\qquad
 W_H=e^{\hbar\kappa D_H}.
\]
The matrix is symmetric and
$H_{m+1,n}+H_{m,n+1}=0$.

\begin{proposition}[Fields, tensors, and formal states]
\label{v4-arith:centers:prop:normal-ordering}
The map $W_H$ is invertible, preserves the vacuum,
differential, and translation, and satisfies
\[
 \begin{split}
 W_HY_K(a,z)b&=Y_{K+H}(W_Ha,z)W_Hb,\\
 \Delta W_H&=(W_H\otimes W_H)\Delta,\qquad
 \pi W_H=W_H\pi,\\
 W_{H_2}W_{H_1}&=W_{H_1+H_2},\qquad W_H^{-1}=W_{-H}.
 \end{split}
\]
Define a coefficient-valued functional independently by
zero odd moments and
\[
 \lambda_H(x_{i_1}\cdots x_{i_{2r}})
 =(-\hbar\kappa)^r
 \sum_{\text{pairings}}\prod_{\{a,b\}}H_{i_ai_b},
 \qquad \lambda_H(1)=1.
\]
Then $\lambda_H=\lambda_0W_H^{-1}$, where $\lambda_0$
sets currents to zero while retaining $\kappa$.
\end{proposition}

\begin{proof}
All contractions commute with $d$.
The recurrence for $H_{mn}$ makes $[T,D_H]=0$.
The second-order Leibniz rule gives
\[
 D_H\mu_E=\mu_E\left(D_H\otimes1+1\otimes D_H+
             \sum_{m,n}H_{mn}\partial_m\otimes\partial_n\right).
\]
Exponentiate and move the cross term past $e^{zT}$.
Its Taylor series is $K_{mn}(z)+
\partial_z^m(-\partial_z)^nH(z)$, proving the field identity.

For the coproduct, each outer $D_H$ equals
$\Delta D_H$ after $\Delta$. The source level splits as
$\kappa_L+\kappa_R$, giving precisely the two target
normal-ordering operators. Their commutation proves
the composition and inverse formulas.

Finally, the constant term of $D_H^r/r!$ counts complete
pairings. The factor $1/2$ removes their two orientations
and $r!$ removes their order. This proves the stated
moment formula, including repeated indices.
\end{proof}

The formal state has a separate bar realization.
Let $I_m$ insert $s\theta_m$ into every bar gap.
Adjacent exterior products cancel, so $I_m$ commutes
with the bar differential. Insertion operators commute.
The vector
\[
 \nu_H=\exp\left(-\frac{\hbar\kappa}{2}
             \sum_{m,n}H_{mn}I_mI_n\right)1
\]
lies in the product over output weights, with a finite
formula at each weight and parameter order.

Its pairing against weight-restricted dual observables
is a continuous chain functional.
The identity $I_mj(f)=j(t_mf)$ and the factorial pairing
give $\langle-,\nu_H\rangle=\lambda_H\pi$.
Thus this comparison follows from two specified constructions
of the state.

These functionals are formal and coefficient-valued.
They retain $\kappa$ and detect contractions, unlike the
bialgebra counit. For example,
\[
 \lambda_H(x_0^2)=-\hbar\kappa H(0),\qquad
 \lambda_H(Y_{K+H}(x_0,z)x_0)=\hbar\kappa K(z).
\]
The second identity cancels the new regular contraction.

On a curve chart, a regular symmetric kernel $H(z,w)$
has diagonal jets $H_{mn}(s)$ satisfying
$\partial_sH_{mn}=H_{m+1,n}+H_{m,n+1}$.
Hence $[\partial_s+T,D_H(s)]=0$.
For supplied kernels whose differences are
$H_{\alpha\beta}$, the identity
\[
 W_{H_{\beta\gamma}}W_{H_{\alpha\beta}}
 =W_{H_{\alpha\gamma}}
\]
gives horizontal gluing of fields and comparisons.
It assumes the kernels and their coordinate identifications.
A projective connection alone does not supply a global
bidifferential or state. A regular symmetric change leaves
the exchange-antisymmetric part of the kernel unchanged.

\subsection{Changing the tensor action}

At unit level and $K(z)=z^{-1}$, the polynomial
state-exchange operator is
\[
 \mathcal S(z)=\exp\left(-2\hbar
 \sum_{m,n}K_{mn}(z)\partial_m\otimes\partial_n\right).
\]
The additive current acts by
$J_{-1}^{\mathrm{tot}}=M_{x_0}\otimes1+1\otimes M_{x_0}$.
They do not commute:
\[
 [\mathcal S(z),J_{-1}^{\mathrm{tot}}](x_0\otimes1)
 =-\frac{2\hbar}{z}\,1\otimes1.
\]
The ordinary flip does intertwine the cocommutative additive
action. Thus field exchange and the displayed module tensor
have different intertwiners.

A change of action gives a positive comparison.
Decorate each admissible module by commuting even chain
operators $D_m$ lowering an exhaustive nonnegative energy
filtration by $m+1$. Each algebraic state has bounded energy.
Morphisms preserve these operators and the field action.

For a finite set of spectral slots $I$, set
\[
 R_I=k[z_i:i\in I][(z_i-z_j)^{-1}:i\ne j],\qquad
 M_I=(R_I\otimes_k\bigotimes_iM_{i,0})[[\hbar]].
\]
The field variable $w$ is independent of the $z_i$.
For an order $\sigma$ define
\[
 C_{ij}=\hbar\sum_{m,n}
 K_{mn}(z_i-z_j)D_m^{(i)}D_n^{(j)},\qquad
 F_\sigma=\exp\left(\sum_{i\prec_\sigma j}C_{ij}\right),
 \qquad
 \rho_\sigma(a,w)=F_\sigma\rho_0(a,w)F_\sigma^{-1}.
\]
The coefficient $\hbar$ is specified tensor data.
It is not forced by an arbitrary coproduct.

A term with $q$ contractions consumes at least
$\sum(m+n+2)$ units of energy and has pole order
$\sum(m+n+1)$. Its pole order is at most
$\sum_i a_i-q$ on inputs of energies $a_i$.
Also $q<N$ modulo $\hbar^N$.
These bounds make $F_\sigma^{\pm1}$ continuous on the
stated carrier and preserve each lower Laurent bound in $w$.

\begin{theorem}[Coherent tensor transport on separated slots]
\label{v4-arith:centers:thm:spectral-transport}
For ordered blocks $B,D$, the maps
$F_{B,D}=\exp(\sum_{i\in B,j\in D}C_{ij})$
give a monoidal equivalence between the original decorated
tensor actions and $\rho_\sigma$.
If $P_{\tau\sigma}$ is the graded permutation of carriers,
then
\[
 c_{\tau\sigma}=F_\tau P_{\tau\sigma}F_\sigma^{-1}
\]
intertwines every field mode and satisfies all block
hexagons and braid identities. On two slots it is
\[
 c^F_{M,N}(z)=
 \tau\exp\left(-2\hbar\sum_{m,n}
 K_{mn}(z)D_m\otimes D_n\right).
\]
\end{theorem}

\begin{proof}
All pair contractions commute, including those sharing
a slot. For three blocks, both parenthesizations give
\[
 \exp(C_{B,D}+C_{B,H}+C_{D,H}).
\]
Every higher parenthesization counts each pair once.
This proves the pentagon. Empty pair sets give the unit.

The original permutation intertwines $\rho_0$.
Conjugating this identity proves the new intertwining
relation on every mode. Consecutive order changes
telescope, while splitting blocks partitions the same
pair set. These facts prove the braid and hexagon identities.
Finally $K_{nm}(-z)=-K_{mn}(z)$ gives the two-slot formula.
The energy and parameter bounds justify every exponential.
\end{proof}

This transported category is monoidally equivalent to the
original symmetric category. Its nonidentity formula
does not by itself construct an inequivalent representation
category.

For two regular polynomial factors put
$M=M_{x_0^{(1)}+x_0^{(2)}}$ and $C=C_{12}$.
Then
\[
 [C,M]=\hbar A(z),\qquad [C,A(z)]=0,
\]
where
$A(z)=\sum_mK_{m0}(z)D_m^{(1)}
+\sum_nK_{0n}(z)D_n^{(2)}$.
The transported creation modes are
\[
 J_{-1}^{\pm}=M\pm\hbar A(z),\qquad
 e^{-2C}J_{-1}^+=J_{-1}^-e^{-2C}.
\]
On polynomials in the two lowest generators this reads
$x+y\pm\hbar z^{-1}(\partial_x+\partial_y)$.

For the full current,
\[
 J^\pm(w)=J^0(w)\pm\hbar\sum_{n\geq0}
 \bigl(K_{n0}(z-w)D_n^{(1)}
       +K_{0n}(z+w)D_n^{(2)}\bigr),
\]
expanded in nonnegative powers of $w$.
The first commutator gives the correction and the second
vanishes. The fields of the higher $y_I$ are unchanged:
the decorations kill those generators and their translates.
The negative-degree resolution cells are retained.

Equivariance of $\pi$ gives a quasi-isomorphism
$\Pi_I=\pi^{\widehat\otimes I}$ intertwining all transported
fields, tensorators, and exchanges.
The finite-weight contraction extends coefficientwise
to the parameter completion.

There are compatible formal states as well.
Define $\lambda_\sigma$ by complete matchings between
distinct slots, assigning
$-\hbar K_{mn}(z_i-z_j)$ to a paired occurrence
$x_m^{(i)},x_n^{(j)}$ with $i\prec_\sigma j$.
The bar-side definition uses
\[
 \nu_\sigma=
 \exp\left(-\hbar\sum_{i\prec_\sigma j}\sum_{m,n}
 K_{mn}(z_i-z_j)I_m^{(i)}I_n^{(j)}\right)1^{\otimes I}.
\]
The pairing-count argument proves
\[
 \lambda_\sigma=\lambda_0^{\otimes I}F_\sigma^{-1},
 \qquad
 \lambda_\sigma^E=\lambda_\sigma\Pi_I,\qquad
 \lambda_\tau c_{\tau\sigma}=\lambda_\sigma.
\]

For four lowest currents,
\[
 \lambda_{1234}(x_0^{\otimes4})
 =\hbar^2\left(
 \frac1{z_{12}z_{34}}+\frac1{z_{13}z_{24}}+
 \frac1{z_{14}z_{23}}\right).
\]
At $(1,3,7,13)$ it is $29\hbar^2/240$.
These are normalized formal functionals on the separated
tensor spaces, with no positivity assertion.

On the actual order-chart overlap, in one fixed tensor order,
\[
 G_{\tau\sigma}=F_\tau F_\sigma^{-1}.
\]
Only pairs whose orders differ occur in its exponent.
Their differences are precisely the ones inverted
on that overlap. On regular $E,P$ coefficients the connection
is $\nabla_i=\partial_{z_i}+T_i$.
The kernel recurrence gives
$[T_i,C_{ij}]=-\partial_{z_i}C_{ij}$, so $G_{\tau\sigma}$
is horizontal. Cancellation gives its cocycle identity.

These maps glue differential modules, but they do not
supply the unchanged tensor-power collision equivalence.
Modulo $\hbar$, diagonal restriction of the free
two-point module leaves $P\otimes P$.
Its energy-one piece has dimension two, whereas that of
$P$ has dimension one. The normal Koszul complex adds
no differential that changes these ranks.
Thus an energy-preserving collision equivalence with
singleton $P$ is impossible on this carrier.
Additional collision relations or coefficients are needed.

\subsection{An actual spectral symmetry Hopf algebra}

Let $\mathfrak t$ have generators $T,p_m$, $m\geq0$, with
\[
 [p_m,p_n]=0,\qquad [T,p_m]=-p_{m-1},\qquad p_{-1}=0.
\]
Give $\mathscr H=U(\mathfrak t)[[\hbar]]$ the primitive
coproduct. On modules with $T$ raising energy and $p_m$
lowering it by $m+1$, define
\[
 \mathcal R(z)=
 \exp\left(-2\hbar\sum_{m,n}K_{mn}(z)p_m\otimes p_n\right).
\]
This is an operator in the completion defined by these
modules. It need not be a finite tensor in $\mathscr H^{\otimes2}$.

The exponent commutes with each $\Delta p_m$.
Its two commutators with $T$ are opposite $z$-derivatives,
so it also commutes with $\Delta T$.
Therefore $\tau\mathcal R(z)$ is a module intertwiner.
Applying the primitive coproduct splits its exponent,
proving both hexagons. The pair exponents commute, giving
Yang--Baxter at arguments $u,u+v,v$ and
\[
 \mathcal R^{21}(-z)\mathcal R(z)=1.
\]

For $E,P$, take $p_m=\partial_m$.
The field has shifted equivariance
\[
 hY_E(a,z)b
 =\sum Y_E\bigl((e^{-z\operatorname{ad}T}h_{(1)})a,z\bigr)
             h_{(2)}b.
\]
This follows by moving $h_{(1)}$ through $e^{zT}$.
The shifted coproduct commutes with the contractions.
It identifies the symmetry acting in this model.
The earlier creation-mode defect prevents identifying
that symmetry with an arbitrary vertex-module tensor action.

\subsection{Noncommuting contractions and the singlet order}

Permutation operators give a precise limit of the
commuting-contraction argument.
On $(k^2)^{\otimes3}$, scalar normalization of
$e^{\hbar P/u}$ gives $1+f(u)P$ with
$f(u)=\tanh(\hbar/u)$.
Direct multiplication yields
\[
 \begin{split}
 &(1+aP_{12})(1+bP_{13})(1+cP_{23})\\
 &\quad-(1+cP_{23})(1+bP_{13})(1+aP_{12})\\
 &=(ac-b(a+c))(P_{12}P_{23}-P_{23}P_{12}).
 \end{split}
\]

For $a=f(u)$, $b=f(u+v)$, $c=f(v)$, the first scalar defect is
\[
 -\frac{\hbar^4(u^2+uv+v^2)}
 {3u^2v^2(u+v)^2}+O(\hbar^6).
\]
The commutator matrix has rank four.
At $u=v=1$ its order-four coefficient is $-1/4$.
A scalar prefactor multiplies both Yang--Baxter sides
by the same factor and cannot remove this defect.

By contrast, $f(u)=\hbar/u$ satisfies $ac=b(a+c)$.
Hence
\[
 R_{\mathrm{rat}}(u)=\frac{u\,1+\hbar P}{u+\hbar}
\]
satisfies Yang--Baxter and
$R_{\mathrm{rat}}^{21}(-u)R_{\mathrm{rat}}(u)=1$.
The integral RTT Hopf algebra and the specified
Bethe-algebra representation of
Section~\ref{v4-arith:ic:formal-operators}
retain this positive algebraic domain.
Neither existence nor nonexistence of an arithmetic
realization follows from the absence of a subtraction
operation between local fields.

Likewise, the quadratic Hecke relation
$(\check R-q)(\check R+q^{-1})=0$ gives
\[
 \check R-\check R^{-1}=(q-q^{-1})1.
\]
For $C(z)=z\check R-z^{-1}\check R^{-1}$ it implies
\[
 C(z^{-1})C(z)=(q^2+q^{-2}-z^2-z^{-2})1.
\]
This calculation remains valid at repeated roots.
Yang--Baxter additionally requires a Hecke braid
representation. The involution
$\operatorname{diag}(1,1,1,-1)$ at $q=1$ satisfies
the quadratic relation and fails the adjacent braid relation.

For the singlet, let
\[
 j(v)=(e_0\otimes e_1-e_1\otimes e_0)\otimes v,\qquad
 \pi=\tfrac12j^{\mathsf t},\qquad R(u)=1+\hbar P/u.
\]
Then
\[
 R_{23}(u+\hbar)R_{13}(u)j
 =\frac{u+2\hbar}{u+\hbar}j,\qquad
 \pi R_{13}(u)R_{23}(u+\hbar)
 =\frac{u+2\hbar}{u+\hbar}\pi.
\]
Indeed, for $A=P_{13}$, $B=P_{23}$,
$(A+B)j=j$ and $BAj=-Aj$.
With $a=\hbar/u$, $b=\hbar/(u+\hbar)$,
the unwanted coefficient is $a-b-ab=0$.
In the reverse order it is $a-b+ab\ne0$.

The normalization preserving this cup and cap therefore obeys
\[
 f(u)f(u+\hbar)=\frac{u+\hbar}{u+2\hbar}.
\]
It has a unique solution with constant term one at infinity:
each coefficient is twice the new unknown plus earlier terms.
Writing $t=\hbar/u$ gives
\[
 F(t)=1-\tfrac12t+\tfrac58t^2-\tfrac{11}{16}t^3+\cdots,
 \qquad F(t)F(-t)=\frac1{1-t^2}.
\]

For the last identity, put $G(t)=F(t)F(-t)$.
The shift equation implies
$G(t)G(t/(1+t))=(1+t)/((1-t)(1+2t))$.
Both $G$ and $(1-t^2)^{-1}$ solve it with constant term one.
The same triangular uniqueness proves equality.
Thus the normalized exchange is formally unitary.

For fixed $\hbar\ne0$, there is no rational solution.
If $d_n$ is its order at $-2\hbar+n\hbar$, the equation
requires $d_n+d_{n+1}$ to be $-1$ at $n=0$, $1$ at $n=1$,
and zero elsewhere. Finite support forces its largest
nonzero index into $\{0,1\}$ and its smallest into
$\{1,2\}$. Only $d_1$ could remain, but it would have to
equal both $-1$ and $1$.
Meromorphic normalizations require separate analytic choices.
The algebraic cup calculation also leaves the comparison
with a chosen Yangian evaluation-module duality to its
specified coproduct and parameter conventions.

\section{Half-braidings and their probes}
\label{v4-arith:centers:sec:probes}

The objects against which a half-braiding is tested are
part of its definition. For geometric words, the allowed
maps of the labeled base matter as well.

Let $X$ be smooth and separated of finite type, with
arbitrarily many distinct closed $k$-points.
Use word objects with monotone collision equivalences,
the prefix/suffix cut tensor, and its isolated empty unit.
Require comparison maps over the fixed labeled base.
For a point $x$, let $P_x$ have zero empty coefficient
and positive coefficient $(x,\ldots,x)_*k$.

\begin{proposition}[The absolute point probes]
\label{v4-arith:centers:prop:point-probes}
If $F\star P_x\simeq P_x\star F$ for every closed point $x$,
then every positive coefficient of $F$ vanishes.
\end{proposition}

\begin{proof}
Choose the least positive $n$ with $F_n\ne0$.
In arity $n+1$, restrict to
$(X\setminus\{x\})^n\times\{x\}$.
On the left, the cut after $n$ contributes
$F_n|_{(X\setminus\{x\})^n}$.
All smaller positive cuts vanish by minimality.
The empty term has constant support and also vanishes.

On the right the remaining positive cut requires
the first coordinate to be $x$, so it vanishes too.
Thus $F_n$ restricts to zero on each
$(X\setminus\{x\})^n$.

Choose $n+1$ distinct points. Every $n$-tuple omits
at least one, so these opens form a finite cover of $X^n$.
Descent gives $F_n=0$, a contradiction.
\end{proof}

This result uses the cut tensor, all point probes,
and the fixed labeled base. It does not apply to the
symmetric convolution using all ordered decompositions.

Restricting to one point gives a positive alternative.
Objects supported at the constant tuple $x$ correspond
to pairs of complexes $(V_0,V_1)$, because the positive
collision diagram has terminal singleton.
Their tensor is
\[
 (V_0,V_1)\star(W_0,W_1)=
 (V_0\otimes W_0,\,
  V_0\otimes W_1\oplus V_1\otimes W_0).
\]

Two positive factors give zero on the cut-disjoint open.
Every nonzero iterated term therefore has at most one
positive factor. The Koszul flip on these terms supplies
a symmetry and all its coherences.
Nonzero positive objects have half-braidings in this
relative probe category.

\subsection{Charge categories and excess central objects}

Even a braided category can have central objects outside
its canonical central image.
For finite-dimensional $\mathbb Z$-graded vector spaces,
choose any bicharacter braiding.
On the unit, the scalars $T^n$ on charge $n$ define
a half-braiding for every $T\in\mathbb C^\times$.
Its hexagon is $T^{n+m}=T^nT^m$.
The canonical central image with underlying unit has
trivial half-braiding, so the examples with $T\ne1$
are additional central objects.
A fixed-level Heisenberg realization would require its
own fusion tensor and monoidal comparison.

There is a finite modular example with all coherences visible.
Let $\ell>0$ be even and $G=\mathbb Z/\ell\mathbb Z$.
For representatives $0\leq a,b<\ell$, define
\[
 c(a,b)=\left\lfloor\frac{a+b}{\ell}\right\rfloor,\qquad
 [a+b]=a+b-\ell c(a,b).
\]
On finite-dimensional $G$-graded vector spaces set
\[
 \omega(a,b,c)=(-1)^{a c(b,c)},\qquad
 R(a,b)=e^{\pi iab/\ell},\qquad
 \theta(a)=e^{\pi ia^2/\ell}.
\]
The first scalar is the associator from left to right
parenthesization, the second multiplies the interchange,
and the third is the twist.

\begin{proposition}[The finite cyclic center]
\label{v4-arith:centers:prop:cyclic-center}
These data define a modular ribbon category $\mathcal C_\ell$.
Its simple objects have dimension one.
The center has $\ell^2$ simples, indexed by $(a,j)\in G^2$,
with canonical central image $(a,0)$.
The braided functor
\[
 \mathcal C_\ell\boxtimes\mathcal C_\ell^{\mathrm{rev}}
 \longrightarrow Z(\mathcal C_\ell)
\]
is an equivalence and sends $(a,b)$ to $([a+b],-b)$.
\end{proposition}

\begin{proof}
The carry identity
\[
 c(b,c)+c([b+c],d)=c(c,d)+c(b,[c+d])
\]
expresses the quotient and remainder of $b+c+d$.
The pentagon exponent discrepancy is
$\ell c(a,b)c(c,d)$, which is even.
The two hexagons are
\[
 \begin{split}
 R(a,[b+c])
 &=\frac{\omega(b,a,c)}
 {\omega(a,b,c)\omega(b,c,a)}R(a,b)R(a,c),\\
 R([a+b],c)
 &=\frac{\omega(a,b,c)\omega(c,a,b)}
 {\omega(a,c,b)}R(a,c)R(b,c).
 \end{split}
\]
Symmetry of the carry cancels two factors in each ratio.
The remaining sign is exactly the phase from reducing
the sum modulo $\ell$. Without the associator, the first
hexagon already fails at $\ell=2$, $a=b=c=1$.

Balancing follows from
\[
 \theta([a+b])=R(a,b)R(b,a)\theta(a)\theta(b).
\]
The exponent discrepancy is
$\pi i[-2c(a,b)(a+b)+\ell c(a,b)^2]$.
Choose dual $[-a]$, evaluation one, and coevaluation
$\omega(a,[-a],a)^{-1}$.
The carry identities prove both triangles, and
$\theta([-a])=\theta(a)$ gives dual compatibility.
The ribbon dimension is
\[
 \theta(a)R(a,[-a])\omega(a,[-a],a)^{-1}=1.
\]
Double braiding is the Fourier matrix
$e^{2\pi iab/\ell}$, with inverse
$\ell^{-1}e^{-2\pi iab/\ell}$, proving nondegeneracy.

For a central object of degree $a$, divide its
half-braiding with degree $b$ by $R(a,b)$.
The hexagon makes the remaining operators a representation
of $G$. It splits into characters
$e^{2\pi ijb/\ell}$, giving the stated simple labels.
The canonical image has $j=0$.

Using the original crossing on the first tensor factor
and the inverse crossing on the second gives
$([a+b],-b)$. The associator factors are those just checked.
The mixed double braiding cancels, so this is braided.
It is bijective on simples and their scalar Hom spaces,
hence is an equivalence of semisimple categories.
\end{proof}

The even-level condition has an integral normalization.
With minimum electric charge one and $F/(2\pi)$ representing
$c_1$, the action $\ell(4\pi)^{-1}\int A\wedge dA$
changes under a four-dimensional extension by
\[
 \exp\left(2\pi i\,\frac\ell2
                  \int(F/(2\pi))^2\right).
\]
Even $\ell$ makes this one on every oriented extension.
The hyperplane line bundle on $\mathbb{CP}^2$
proves necessity. In this bosonic convention the class is
$(\ell/2)c_1^2\in H^4(BU(1);\mathbb Z)$.
The cell structure of $\mathbb{CP}^\infty$ gives
$H^*(BU(1);\mathbb Z)=\mathbb Z[c_1]$ and $H^3=0$.

For the marked boundary current with
$J(z)J(w)\sim\ell(z-w)^{-2}$, the stress tensor
$(2\ell)^{-1}:JJ:$ has double-contraction coefficient
$2\ell^2/(4\ell^2)=1/2$, hence Virasoro central charge one.
The current commutator, integral charge lattice, and
quadratic form still retain $\ell$.
An unmarked rescaling does not preserve these data.

\section{Arithmetic centers and geometric comparisons}
\label{v4-arith:centers:sec:arithmetic}

The arithmetic constructions require the same distinctions.
An ordinary center acts by natural endomorphisms.
A scalar Hochschild center is a complex with its higher
operations. A categorical center is itself a category.
An identity trace is a class in a trace recipient.
A comparison must specify which of these objects it carries.

For a stable presentable $k$-linear monoidal category
$\mathcal C$, with tensor preserving colimits separately, put
\[
 \begin{split}
 Z_{\mathrm{cat}}(\mathcal C)
 &=\operatorname{Fun}^{\mathrm L}_{
       \mathcal C\otimes\mathcal C^{\mathrm{op}}}
       (\mathcal C,\mathcal C),\\
 \operatorname{Tr}_{\mathrm{cat}}(\mathcal C)
 &=\mathcal C\otimes_{
       \mathcal C\otimes\mathcal C^{\mathrm{op}}}\mathcal C.
 \end{split}
\]
The tensor and functors are in presentable $k$-linear
categories. The first has a braided tensor structure.
The second carries a circle action.
For an ordinary algebra $A$, the scalar complexes
$R\operatorname{Hom}_{A^e}(A,A)$ and
$A\otimes_{A^e}^{\mathbf L}A$ occur one categorical level
lower. They are not identified by their shared name
Hochschild.

The precise positive geometric comparisons of
Section~\ref{v4-arith:ic:traces} apply here.
For a perfect derived stack $X$,
\[
 Z_{\mathrm{cat}}(\operatorname{QCoh}(X))
 \simeq\operatorname{QCoh}(\mathcal LX)
 \simeq\operatorname{Tr}_{\mathrm{cat}}
             (\operatorname{QCoh}(X)).
\]
For the complex reductive Hecke category
$D(B\backslash G/B)$, the center and trace identify with the
character category generated by the Harish--Chandra transform.
The same holds for its specified unipotently monodromic version.
These are the theorems of Ben-Zvi--Francis--Nadler and
Ben-Zvi--Nadler on their respective geometric carriers.
They do not identify a category with a ring of functions
on an unspecified arithmetic parameter space.

Even a smooth affine line has derived loops:
\[
 \mathcal O(\mathcal L\mathbb A^1)
 =k[x]\otimes^{\mathbf L}_{k[x_1,x_2]}k[x]
 \simeq k[x,\eta],\qquad |\eta|=-1,\quad d=0.
\]
The diagonal Koszul resolution gives this formula.
Thus smoothness of $X$ does not make its derived loop
space classical.

For finite cyclic loops of $B\mathbb G_m$,
the complex points instead have $\mu_n$ components,
each with stabilizer $\mathbb C^\times$.
A quotient of finite cyclic groups reverses the induced
mapping-stack arrow. These constructions do not form
formal neighborhoods of the identity by notation alone.
An arithmetic loop comparison must specify its source
topology and retain the actual Weil and inertia data.

\subsection{Scalar actions and the full categorical center}

Assume characteristic zero, with coefficient prime different
from the residue prime. The stated $GL_2$ Bernstein theorem
identifies the ordinary Bernstein center with functions on its
continuous Weil parameter stack.
Its normalization retains inertia and Frobenius.
The fixed-character and coefficient-prime comparisons
are different theorems.

For the latter, the ring
\[
 E_{\mathcal B}=\operatorname{End}(P_{\mathcal B})
 =\operatorname{End}_G(J_{\mathcal B})^{\mathrm{op}}
\]
of Section~\ref{v4-arith:ic:paskunas}
is pseudocompact and finite over its pseudodeformation
center. In the generic reducible case its full algebra
retains the two opposite extension representations.
The center does not replace that algebra.

A scalar action on a spectral category need not exhaust
its center. For example, objects of
$\operatorname{QCoh}(B\mathbb G_m)$ decompose into
integer weights. Each sequence $(a_n)_{n\in\mathbb Z}$
acts naturally by $a_n$ on weight $n$.
The weight-zero projector is a central operation
that is not multiplication by a scalar function on
$B\mathbb G_m$, whose global functions are $k$.
In the scalar stabilizer of a $GL_2$ parameter,
a determinant twist changes weight by two.
It cannot identify every weight sector or eliminate
these projectors.

Similarly, a matrix-algebra description on a generic
locus gives a Morita equivalence there, with its center.
It does not extend across a singular locus without
a compatible functor and its inverse.
A regular ring alone does not make extensions vanish
at every point: for $A=k[t]$,
\[
 \operatorname{Ext}^1_A(k,k)=k.
\]
Localization to $k(t)$ kills this supported module.
The generic and supported categories retain different
information.

\subsection{Gluing central actions}

Formal patching supplies a positive ordinary-center
calculation with actual overlap maps.
For
\[
 A=k[t],\quad B=k[[t]],\quad
 C=k[t,t^{-1}],\quad D=k((t)),
\]
the equivalence of finitely generated module categories
in Theorem~\ref{v4-arith:ic:w5-h:formal-patching}
uses triples $(P,Q,\eta)$ with an isomorphism over $D$.
The ordinary center is
\[
 Z(\operatorname{mod}_{\mathrm{fg}}A)
 \simeq B\times_D C=A.
\]
Indeed, a natural endomorphism on finitely generated
$R$-modules is multiplication by its value on
$1\in R$. The identity $B\cap C=A$ inside $D$
then gives the formula.

The punctured formal disk is an actual overlap.
It is not the empty intersection of a closed point
with its complementary open.

Higher centers need more than this scalar equalizer.
An equivalence of categories with effective descent
transports central operations together with their
coherences. An isomorphism of rings of parameters
does not supply the categorical descent equivalence.
For a derived external product followed by pullback,
ordinary formulas also need two Tor-independence
conditions, one for each step.
Already
\[
 k\otimes^{\mathbf L}_{k[[t]]}k\simeq k\oplus k[1]
\]
has an extra term although both factors are perfect.
An identity pullback preserves it.

For quasi-smooth characteristic-zero charts, singular
support is a cone in
\[
 \operatorname{Sing}(X)=
 \operatorname{Spec}_{X^{\mathrm{cl}}}
 \operatorname{Sym}H^1(T_X).
\]
Its linear functions act in Hochschild degree two.
It is not an unspecified closed locus in $X$ itself.
An arithmetic category with this support therefore needs
the derived charts, their coefficient theory, and the
compatible cone. The finite-stratum generator argument
also needs extension functors preserving compactness
and the actual gluing morphisms.
The localization and recollement constructions of
Chapter~\ref{v4-arith:ic:chapter} retain those extensions.

The restricted tensor theorem of
Section~\ref{v4-arith:ic:restricted-tensors}
gives another positive assembly.
Its finite-stage insertion maps preserve the specified
operadic structures, including units.
The filtered colimit transports that structure.
Local endofunctors act on cofinal diagrams containing
their places and commute through tensor interchange.
Neither scalar commutativity nor these separate-place
actions creates collision maps or raises an $E_1$
structure to $E_2$.

The operad of little intervals has an integral associative
comparison. Each ordering component is a nonempty convex
region of endpoints. Its chains therefore augment
quasi-isomorphically to $R[S_n]$ in every arity.
Substitution respects orderings.

Rectification of an algebra along this comparison,
when available, does not make the algebra formal.
The cubic example of
\eqref{v4-prologue:eq:massey} retains its nonzero triple
Massey product in a strict associative model.
For centers, the chain-level $E_2$ operations, coefficient
base, and comparison maps must likewise be retained.
A rational operadic comparison alone supplies no
integral replacement or arithmetic geometric action.

\section{The marked arithmetic comparison}
\label{v4-arith:centers:sec:marked-comparison}

Let $\mathcal C_I$ be the arithmetic fusion category of
Chapter~\ref{v4-arith:mf:chapter}, over its specified field
$F_I$. Write
\[
 U_I:Z(\mathcal C_I)\longrightarrow\mathcal C_I,
 \qquad U_I\dashv R_I,\qquad
 B_I=R_I(\mathbf1).
\]
The distinguished algebra $B_I$ is defined by adjunction:
its multiplication is adjoint to
\[
 U_I(B_I\otimes B_I)
 \simeq U_I(B_I)\otimes U_I(B_I)
 \xrightarrow{\epsilon\otimes\epsilon}\mathbf1,
\]
and its unit is adjoint to $1_{\mathbf1}$.
Its module category reconstructs the boundary category.
Consequently the comparison object is the pair
$(Z(\mathcal C_I),B_I)$ with this multiplication and
its forgetful action.

Suppose an independently defined boundary or bulk construction
provides a braided category $\mathcal B_I$ and an algebra
$\mathfrak b_I$ in it. A marked comparison consists of
a braided equivalence
\[
 \Xi_I:Z(\mathcal C_I)\xrightarrow{\sim}\mathcal B_I
\]
and an algebra isomorphism
$\xi_I:\Xi_I(B_I)\xrightarrow{\sim}\mathfrak b_I$.
Using the tensor comparison of $\Xi_I$, it must satisfy
\[
 \xi_I\circ\Xi_I(\mu_{B_I})
 =\mu_{\mathfrak b_I}\circ(\xi_I\otimes\xi_I),
 \qquad
 \xi_I\circ\Xi_I(\eta_{B_I})=\eta_{\mathfrak b_I}.
\]
The induced module-category comparison must intertwine
the central action on the specified boundary objects.
An unmarked braided equivalence does not provide these
additional identifications.

The arithmetic marking also has semilinear transport.
For an ideal $\mathfrak a$ in the domain of
\eqref{v4-arith:mf:artin-input}, put
$\sigma=\operatorname{Art}(\mathfrak a)$ and
$I'=\mathfrak a^{-1}I$.
The normalized network coefficients satisfy
$\sigma(e_{I,g})=e_{I',g}$, with the stated label transport.

A tensor transport lifting this law must be
$\sigma$-semilinear on morphisms and preserve the
marked junctions and duality.
Denote such a transport by
$\mathsf A_{\mathfrak a}:\mathcal C_I\to\mathcal C_{I'}$.
It induces a center transport and, by the adjunction,
a marked algebra isomorphism
\[
 Z(\mathsf A_{\mathfrak a})(B_I)\simeq B_{I'}.
\]
Its multiplication compatibility follows by transporting
the counit definition above.

The boundary or bulk construction must supply its own
semilinear transport
$\mathsf T_{\mathfrak a}:\mathcal B_I\to\mathcal B_{I'}$.
The required square is a coherent braided isomorphism
\[
 \mathsf T_{\mathfrak a}\Xi_I
 \simeq
 \Xi_{I'}Z(\mathsf A_{\mathfrak a}),
\]
compatible with the algebra maps $\xi_I,\xi_{I'}$,
the units, and composition of Artin transports.
If these tensor functors have not been constructed,
the scalar reciprocity identity is only their necessary
coefficient condition.

\begin{question}[The marked center comparison]
\label{v4-arith:centers:qn:arithmetic}
Construct the independently specified boundary or bulk
center, its distinguished algebra and central action,
and the equivalences $\Xi_I,\xi_I$ with their semilinear
transport compatibilities. On a field realization, retain the root and full collision
diagrams. Construct the separated products and action,
and identify the passage to the rational boundary sectors.
\end{question}

The affine $E_6$ theorem solves the tensor-category
center problem in that realization and identifies its
distinguished algebra and mirror extension.
The pointed obstruction shows why the corresponding
$U$-state algebra cannot simply act on the original
vacuum in the required pointed manner.
The derived Heisenberg example independently shows that
even the correct one-position endomorphism object may
miss action homotopies on separated configurations.
Both results constrain the comparison without replacing
its original domain.

A solution must therefore provide the stated maps,
or an invariant incompatible with every permitted
marked comparison. Equality of fusion rules, characters,
scalar networks, or ordinary centers does not decide it.
The trace comparison, the sewing of higher-genus
amplitudes, and physical bulk reconstruction require
their own operation-preserving maps as well.

\section*{References}

For the scalar center, see Keller,
\emph{Derived invariance of higher structures on the
Hochschild complex}. The higher operations use McClure--Smith,
\emph{A solution of Deligne's Hochschild cohomology conjecture},
arXiv:math/9910126.
The internal centralizer and quasi-unit results are in
Lurie, \emph{Higher Algebra}, Sections 4.4.1, 5.3.1,
and 5.4.3, with the theorem numbers specified above.

The chiral center, factorization lattices, and
Chevalley--Cousin comparison are in
Beilinson--Drinfeld, \emph{Chiral Algebras},
Sections 3.3--3.4 and 3.7.
For the differential-module pairing and projection formula,
see Drinfeld--Gaitsgory, \emph{On some finiteness questions
for algebraic stacks}, arXiv:1108.5351, Section 5.
The holonomic functors use Mili\v{c}i\'c,
\emph{Lectures on Algebraic Theory of D-Modules},
Chapter I, Section 14.

The fusion-center construction uses
Davydov--M\"uger--Nikshych--Ostrik,
\emph{The Witt group of non-degenerate braided fusion
categories}, arXiv:1009.2117, Sections 3--4.
Gelaki--Naidu--Nikshych,
\emph{Centers of graded fusion categories},
arXiv:0905.3117, Corollary 3.7. And Bischoff,
\emph{A remark on CFT realization of quantum doubles
of subfactors. Case index less than four},
arXiv:1506.02606, Corollary 4.7 and Section 5.

The vertex extension inputs are Huang--Kirillov--Lepowsky,
\emph{Braided tensor categories and extensions of vertex
operator algebras}, arXiv:1406.3420, Theorems 3.2--3.6.
Creutzig--Kanade--McRae,
\emph{Tensor categories for vertex operator superalgebra
extensions}, arXiv:1705.05017, Theorem 3.65. And their
\emph{Gluing vertex algebras}, arXiv:1906.00119,
Theorems 5.2, 5.5, 5.12 and Corollary 5.13.
The holomorphic classification is Dong--Mason,
\emph{Holomorphic vertex operator algebras of small
central charges}, arXiv:math/0203005, Theorem 1.
The hypotheses on grading and strong rationality remain
part of these applications.
